% Capacitors — Physics Upper Sixth — Cours signé OnBuch+
% Official MINESEC syllabus. Compiler avec : tectonic PHY05-capacitors.tex
\newcommand{\DOCMATIERE}{Physics}
\newcommand{\DOCNIVEAU}{Upper Sixth}
\newcommand{\DOCMODULE}{Upper Sixth — Physics}
\newcommand{\DOCLECON}{5}
\newcommand{\DOCTITRE}{Capacitors}
% OnBuch+ — préambule « cours signés » ENRICHI (compilé avec tectonic / XeLaTeX)
% Système visuel « Pop » d'OnBuch+ : fond crème (jamais blanc), bordures encre
% épaisses, ombres « dures » décalées, titres Archivo Black en capitales.
% Version MATHÉMATIQUES : graphes (pgfplots/tikz), boîtes pédagogiques
% (définition, propriété, méthode, attention, le savais-tu, exercice résolu,
% exercice, TP, bilan), page de garde et signature OnBuch+ ancrée en bas.
%
% Macros attendues dans main.tex AVANT \input{preamble} :
%   \DOCMATIERE  \DOCNIVEAU  \DOCMODULE  \DOCLECON (numéro)  \DOCTITRE
\documentclass[11pt]{article}

\usepackage{fontspec}
\usepackage[english]{babel}
\usepackage[a4paper,margin=2cm,top=3.4cm,bottom=2.8cm,headheight=1.4cm,footskip=1.3cm]{geometry}
\usepackage{amsmath,amssymb,amsfonts}
\usepackage{mathtools} % \xmapsto, \coloneqq... souvent utilisés par les modèles
\usepackage{cancel} % simplifications barrées (bilans de forces)
% \abs{z} (module, valeur absolue) et \norm{u} (norme), utilisés par les modèles
\providecommand{\abs}{}\let\abs\relax\DeclarePairedDelimiter{\abs}{\lvert}{\rvert}
\providecommand{\norm}{}\let\norm\relax\DeclarePairedDelimiter{\norm}{\lVert}{\rVert}
\usepackage[table]{xcolor} % [table] : fournit aussi \rowcolors (alternance de lignes)
\usepackage{pagecolor}
\usepackage{tikz}
\usetikzlibrary{arrows.meta,positioning,calc,shapes.geometric,shapes.misc,shapes.arrows,decorations.pathmorphing,decorations.pathreplacing,decorations.markings,patterns,shadows,mindmap,trees,fit,backgrounds,matrix,intersections,angles,quotes}
\usepackage{pgfplots}
\usepackage[version=4]{mhchem} % équations nucléaires \ce{^{235}_{92}U}
\usepackage[european,siunitx]{circuitikz} % schémas électriques
\pgfplotsset{compat=1.18}
\usepgfplotslibrary{fillbetween,groupplots} % aires entre courbes (calcul intégral)
\usepackage{enumitem}
\usepackage{fancyhdr}
% [most] charge toutes les bibliothèques tcolorbox (listings, minted, raster,
% documentation...) et gonfle inutilement la mémoire TeX sur un document long
% (~300 boîtes) : on ne charge que ce qui est réellement utilisé.
\usepackage[skins,breakable]{tcolorbox}
\usepackage{graphicx}
\usepackage{colortbl}
\usepackage{booktabs}
\usepackage{tabularx}
\usepackage{diagbox} % case d'en-tête diagonale des tableaux à double entrée
% Colonnes extensibles centrée / gauche / droite (C, L, R), souvent utilisées par les modèles
\newcolumntype{C}{>{\centering\arraybackslash}X}
\newcolumntype{L}{>{\raggedright\arraybackslash}X}
\newcolumntype{R}{>{\raggedleft\arraybackslash}X}
\usepackage{array}
\usepackage{multicol}
\usepackage{siunitx}
% parse-numbers=false : accepte aussi \SI{2,5 \times 10^{-3}}{...}
\sisetup{locale=UK,output-decimal-marker={.},group-minimum-digits=4,parse-numbers=false}
\DeclareSIUnit\ohm{\ensuremath{\symup{\Omega}}} % U+2126 absent de la police
\DeclareSIUnit\division{div} % réglages d'oscilloscope (ms/div, V/div)
\DeclareSIUnit\tr{tr} % tours (vitesse de rotation, ex. \SI{1500}{\tr\per\minute})
\usepackage{qrcode}
% \degree : commande fréquemment utilisée par les modèles pour un angle ou une
% température en dehors de \SI{}{} (non fournie par les paquets chargés).
\providecommand{\degree}{^{\circ}}
% \qed (fin de démonstration), utilisé par les modèles sans amsthm
\providecommand{\qed}{\hfill$\square$}
% \barre : double barre des valeurs interdites dans un tableau de variations (tkz-tab)
\providecommand{\barre}{\ensuremath{\|}}
% environnement proof (amsthm non chargé) utilisé par les modèles
\ifdefined\proof\else\newenvironment{proof}[1][Proof]{\par\noindent\textit{#1.}\ }{\qed\par}\fi
\usepackage{lastpage}
\usepackage{hyperref}
\hypersetup{hidelinks}

% ============================================================
% Palette « Pop » OnBuch+ — mêmes hex que la classe P (plus_ui.dart)
% ============================================================
\definecolor{poporange}{HTML}{F59321}
\definecolor{popdark}{HTML}{E07A0C}
\definecolor{popcream}{HTML}{FFF6E8}
\definecolor{popink}{HTML}{1C1714}
\definecolor{poppurple}{HTML}{7A5AE0}
\definecolor{popgreen}{HTML}{2E9E6A}
\definecolor{poppink}{HTML}{E0457B}
\definecolor{popblue}{HTML}{2F7DE1}
\definecolor{popgold}{HTML}{C98A12}
\definecolor{popmuted}{HTML}{8A7F76}
\definecolor{popline}{HTML}{EADFD2}
\definecolor{popgray}{HTML}{8A7F76}
\colorlet{popgrey}{popgray}
% Alias tolérants (le modèle nomme parfois les couleurs différemment)
\colorlet{popwhite}{white}
\colorlet{popblack}{popink}
\colorlet{popred}{poppink}
\colorlet{popyellow}{popgold}
% Teintes claires pour fonds de tableaux / zones
\colorlet{popblueL}{popblue!10!white}
\colorlet{poporangeL}{poporange!14!white}
\colorlet{popgreenL}{popgreen!12!white}
\colorlet{poppurpleL}{poppurple!12!white}
\colorlet{poppinkL}{poppink!12!white}
\colorlet{popgoldL}{popgold!14!white}

\pagecolor{popcream}

% ============================================================
% Typographie
% ============================================================
\defaultfontfeatures{Ligatures=TeX}
\setmainfont{PlusJakartaSans-Medium.ttf}[Path=../../../fonts/,BoldFont=PlusJakartaSans-Bold.ttf]
% Maths en sans-serif (Fira Math) avec chiffres et lettres droites de Plus
% Jakarta, pour que les formules se fondent dans le texte.
\usepackage[math-style=ISO,bold-style=ISO]{unicode-math}
\setmathfont{FiraMath-Regular.otf}[Path=../../../fonts/]
\setmathfont{PlusJakartaSans-Medium.ttf}[Path=../../../fonts/,range={up/num,up/{Latin,latin},\mathup}]
% (icomma retiré : décimales avec point en anglais)
% Caractères mathématiques Unicode absents des polices de texte (Plus Jakarta,
% Archivo Black) : rendus par leur équivalent mathématique, en texte comme en
% formule — sinon ils s'affichent en carré vide (ex. « Arithmétique dans ℤ »).
\usepackage{newunicodechar}
\newunicodechar{ℝ}{\ensuremath{\mathbb{R}}}
\newunicodechar{ℤ}{\ensuremath{\mathbb{Z}}}
\newunicodechar{ℂ}{\ensuremath{\mathbb{C}}}
\newunicodechar{ℕ}{\ensuremath{\mathbb{N}}}
\newunicodechar{ℚ}{\ensuremath{\mathbb{Q}}}
\newunicodechar{𝔻}{\ensuremath{\mathbb{D}}}
\newunicodechar{α}{\ensuremath{\alpha}}
\newunicodechar{β}{\ensuremath{\beta}}
\newunicodechar{γ}{\ensuremath{\gamma}}
\newunicodechar{δ}{\ensuremath{\delta}}
\newunicodechar{ε}{\ensuremath{\varepsilon}}
\newunicodechar{θ}{\ensuremath{\theta}}
\newunicodechar{λ}{\ensuremath{\lambda}}
\newunicodechar{μ}{\ensuremath{\mu}}
\newunicodechar{σ}{\ensuremath{\sigma}}
\newunicodechar{τ}{\ensuremath{\tau}}
\newunicodechar{φ}{\ensuremath{\varphi}}
\newunicodechar{ψ}{\ensuremath{\psi}}
\newunicodechar{ω}{\ensuremath{\omega}}
\newunicodechar{Σ}{\ensuremath{\Sigma}}
% majuscules produites par \MakeUppercase dans les titres (α-aminés -> Α)
\newunicodechar{Α}{\ensuremath{\alpha}}
\newunicodechar{Β}{\ensuremath{\beta}}
\newunicodechar{Γ}{\ensuremath{\Gamma}}
\newunicodechar{Δ}{\ensuremath{\Delta}}
\newunicodechar{Ε}{\ensuremath{\varepsilon}}
\newunicodechar{Θ}{\ensuremath{\Theta}}
\newunicodechar{Λ}{\ensuremath{\Lambda}}
\newunicodechar{Μ}{\ensuremath{\mu}}
\newunicodechar{Τ}{\ensuremath{\tau}}
\newunicodechar{Φ}{\ensuremath{\Phi}}
\newunicodechar{Ψ}{\ensuremath{\Psi}}
\newunicodechar{Ω}{\ensuremath{\Omega}}
\newunicodechar{↦}{\ensuremath{\mapsto}}
\newunicodechar{∈}{\ensuremath{\in}}
\newunicodechar{∉}{\ensuremath{\notin}}
\newunicodechar{∧}{\ensuremath{\wedge}}
\newunicodechar{⇔}{\ensuremath{\Leftrightarrow}}
\newunicodechar{⟺}{\ensuremath{\Longleftrightarrow}}
\newunicodechar{⇒}{\ensuremath{\Rightarrow}}
\newunicodechar{⟹}{\ensuremath{\Longrightarrow}}
\newunicodechar{∘}{\ensuremath{\circ}}
\newunicodechar{∪}{\ensuremath{\cup}}
\newunicodechar{∩}{\ensuremath{\cap}}
\newunicodechar{≡}{\ensuremath{\equiv}}
\newunicodechar{⊂}{\ensuremath{\subset}}
\newunicodechar{⊥}{\ensuremath{\perp}}
\newunicodechar{∃}{\ensuremath{\exists}}
\newunicodechar{∀}{\ensuremath{\forall}}
\newunicodechar{∛}{\ensuremath{\sqrt[3]{\,}}}
\newunicodechar{∜}{\ensuremath{\sqrt[4]{\,}}}
\newunicodechar{⃗}{\ensuremath{{}^{\to}}}
\newunicodechar{ⁿ}{\ifmmode^{n}\else\textsuperscript{\ensuremath{n}}\fi}
\newunicodechar{ˣ}{\ifmmode^{x}\else\textsuperscript{\ensuremath{x}}\fi}
\newunicodechar{ᵇ}{\ifmmode^{b}\else\textsuperscript{\ensuremath{b}}\fi}
\newunicodechar{ʳ}{\ifmmode^{r}\else\textsuperscript{\ensuremath{r}}\fi}
\newunicodechar{ᵘ}{\ifmmode^{u}\else\textsuperscript{\ensuremath{u}}\fi}
\newunicodechar{ʸ}{\ifmmode^{y}\else\textsuperscript{\ensuremath{y}}\fi}
\newunicodechar{ᵃ}{\ifmmode^{a}\else\textsuperscript{\ensuremath{a}}\fi}
\newunicodechar{ᵉ}{\ifmmode^{e}\else\textsuperscript{\ensuremath{e}}\fi}
\newunicodechar{ᵏ}{\ifmmode^{k}\else\textsuperscript{\ensuremath{k}}\fi}
\newunicodechar{ᵖ}{\ifmmode^{p}\else\textsuperscript{\ensuremath{p}}\fi}
\newunicodechar{ᴺ}{\ifmmode^{N}\else\textsuperscript{\ensuremath{N}}\fi}
\newunicodechar{ᶜ}{\ifmmode^{c}\else\textsuperscript{\ensuremath{c}}\fi}
\newunicodechar{ʰ}{\ifmmode^{h}\else\textsuperscript{\ensuremath{h}}\fi}
\newunicodechar{ᵠ}{\ifmmode^{q}\else\textsuperscript{\ensuremath{q}}\fi}
\newunicodechar{ₙ}{\ifmmode_{n}\else\textsubscript{\ensuremath{n}}\fi}
\newunicodechar{ₐ}{\ifmmode_{a}\else\textsubscript{\ensuremath{a}}\fi}
\newunicodechar{ᵢ}{\ifmmode_{i}\else\textsubscript{\ensuremath{i}}\fi}
\newunicodechar{ᵣ}{\ifmmode_{r}\else\textsubscript{\ensuremath{r}}\fi}
\newunicodechar{ₖ}{\ifmmode_{k}\else\textsubscript{\ensuremath{k}}\fi}
\newunicodechar{ᵦ}{\ifmmode_{\beta}\else\textsubscript{\ensuremath{\beta}}\fi}
\newunicodechar{ₓ}{\ifmmode_{x}\else\textsubscript{\ensuremath{x}}\fi}
\newfontfamily\popdisplay{ArchivoBlack-Regular.ttf}[Path=../../../fonts/]
\newfontfamily\poplogo{SpaceGrotesk-Bold.ttf}[Path=../../../fonts/]
\newfontfamily\popsemi{PlusJakartaSans-SemiBold.ttf}[Path=../../../fonts/]
\newfontfamily\popxbold{PlusJakartaSans-ExtraBold.ttf}[Path=../../../fonts/]
\newfontfamily\popxboldls{PlusJakartaSans-ExtraBold.ttf}[Path=../../../fonts/,LetterSpace=3.0]

\setlist[enumerate]{leftmargin=1.7em,itemsep=5pt,topsep=4pt}
\setlist[itemize]{leftmargin=1.7em,itemsep=4pt,topsep=4pt,label={\color{poporange}\textbullet}}
\setlength{\parindent}{0pt}
\setlength{\parskip}{5pt}
\renewcommand{\arraystretch}{1.25}

% pgfplots : style Pop par défaut
\pgfplotsset{
  every axis/.append style={axis line style={popink,line width=1pt},tick style={popink},
    grid style={popline,line width=0.5pt},label style={font=\small},tick label style={font=\footnotesize}},
  popcurve/.style={poporange,line width=1.8pt,no markers,smooth},
}
% Même style disponible dans un \draw TikZ simple (hors pgfplots)
\tikzset{popcurve/.style={poporange,line width=1.8pt}}
\tikzset{popgrid/.style={popline,very thin}}
\colorlet{popcurve}{poporange} % le nom du style est parfois utilisé comme couleur

% ============================================================
% Pill / squiggle
% ============================================================
\newcommand{\poppill}[3]{%
  \tikz[baseline=-0.55ex]\node[draw=popink,line width=1.2pt,rounded corners=2.6pt,%
    fill=#2,inner xsep=6.5pt,inner ysep=3pt]{%
    \popxboldls\fontsize{7}{7}\selectfont\color{#3}\MakeUppercase{#1}};%
}
\newcommand{\popsquiggle}[1]{%
  \tikz[baseline=-0.4ex]{
    \draw[#1,line width=2.6pt,line cap=round]
      plot [smooth] coordinates {(0,0.09) (0.34,0.01) (0.68,0.21) (1.02,0.01) (1.36,0.10)};
  }%
}
% Mot-clé surligné (marqueur orange sous le texte)
\newcommand{\cle}[1]{{\popxbold\color{popdark}#1}}

% ============================================================
% En-tête / pied
% ============================================================
\pagestyle{fancy}
\fancyhf{}
\renewcommand{\headrulewidth}{0pt}
\renewcommand{\footrulewidth}{0pt}
\lhead{\raisebox{-0.15ex}{\poplogo\fontsize{13}{13}\selectfont\color{popink}OnBuch\color{poporange}+}}
\rhead{\poppill{\DOCMATIERE\ \textbullet\ \DOCNIVEAU\ \textbullet\ Chapter \DOCLECON}{white}{popink}}
\lfoot{\popxbold\fontsize{7.6}{7.6}\selectfont\color{popmuted}\MakeUppercase{Signed course OnBuch+}\ \textbullet\ {\popsemi\DOCTITRE}}
\rfoot{\tikz[baseline=-0.6ex]\node[rounded corners=8.5pt,fill=popink,minimum height=17pt,minimum width=17pt,inner xsep=5pt,inner sep=0]{\popxbold\fontsize{8}{8}\selectfont\color{popcream}\thepage\,/\,\pageref*{LastPage}};}

% ============================================================
% Titres
% ============================================================
\newcommand{\chaptitre}[2]{%
  \begin{tcolorbox}[enhanced,colback=poporange,colframe=popink,boxrule=2.6pt,arc=11pt,
      drop shadow=popink,left=15pt,right=15pt,top=12pt,bottom=11pt,before skip=3pt,after skip=18pt]
    \poppill{#2}{popink}{popcream}\par\vspace{9pt}
    {\raggedright\hyphenpenalty=10000\exhyphenpenalty=10000\popdisplay\fontsize{22}{24}\selectfont\color{popcream}\MakeUppercase{#1}\par}
  \end{tcolorbox}
}
% Section numérotée : pastille orange avec le numéro + titre Archivo
\newcounter{popsec}
\newcommand{\coursec}[1]{%
  \stepcounter{popsec}\setcounter{popsub}{0}%
  \par\vspace{16pt}\noindent
  \begin{minipage}{\linewidth}
  \tikz[baseline=-0.7ex]\node[draw=popink,line width=1.6pt,fill=poporange,rounded corners=4pt,minimum size=22pt,inner sep=0]{\popdisplay\fontsize{12}{12}\selectfont\color{popcream}\thepopsec};%
  \hspace{8pt}{\popdisplay\fontsize{15}{17}\selectfont\color{popink}\MakeUppercase{#1}}\par
  \vspace{4pt}\noindent{\color{poporange}\rule{36pt}{3.6pt}}
  \end{minipage}\par\nopagebreak\vspace{8pt}
}
\newcounter{popsub}
\newcommand{\courssub}[1]{%
  \stepcounter{popsub}%
  \par\vspace{9pt}\noindent{\popxbold\fontsize{11.5}{13}\selectfont\color{popink}{\color{poporange}\thepopsec.\thepopsub}\hspace{6pt}#1}\par\nopagebreak\vspace{4pt}
}

% ============================================================
% Boîtes pédagogiques — même recette : fond blanc, bordure épaisse colorée,
% ombre dure encre, badge plein. Titre optionnel [#1].
% ============================================================
\tcbset{popbase/.style={breakable,enhanced,colback=white,boxrule=2.3pt,arc=9pt,
  shadow={3pt}{-3pt}{0pt}{popink},left=14pt,right=14pt,top=11pt,bottom=11pt,before skip=11pt,after skip=15pt,
  fontupper=\popsemi\fontsize{10.6}{14.5}\selectfont\color{popink},
  fontlower=\popsemi\fontsize{10.6}{14.5}\selectfont\color{popink}}}
% Coche dessinée (le glyphe ✓ n'existe pas dans les polices du document)
\newcommand{\popcheck}[1]{\tikz[baseline=-0.5ex]\draw[#1,line width=1.6pt,line cap=round,line join=round](-0.13,0.0)--(-0.04,-0.09)--(0.13,0.1);}
\providecommand{\coche}{\popcheck{popgreen}}
\providecommand{\resultat}[1]{\par\smallskip\noindent\fcolorbox{popgreen}{popgreenL}{\parbox{\dimexpr\linewidth-2\fboxsep-2\fboxrule\relax}{\textbf{Result.} #1}}\par\smallskip}
\newcommand{\popboxhead}[3]{\poppill{#1}{#2}{white}\ifx\relax#3\relax\else\hspace{6pt}{\popxbold\fontsize{10.6}{12}\selectfont\color{popink}#3}\fi\par\vspace{7pt}}

\newtcolorbox{aretenir}[1][]{popbase,colframe=popgreen,before upper={\popboxhead{Key points}{popgreen}{#1}}}
\newtcolorbox{exemplebox}[1][]{popbase,colframe=poporange,before upper={\popboxhead{Exemple}{poporange}{#1}}}
\newtcolorbox{definition}[1][]{popbase,colframe=popblue,before upper={\popboxhead{Definition}{popblue}{#1}}}
\newtcolorbox{propriete}[1][]{popbase,colframe=poppurple,before upper={\popboxhead{Property}{poppurple}{#1}}}
\newtcolorbox{methode}[1][]{popbase,colframe=popgold,before upper={\popboxhead{Method}{popgold}{#1}}}
\newtcolorbox{attention}[1][]{popbase,colframe=poppink,before upper={\popboxhead{Warning}{poppink}{#1}}}
\newtcolorbox{experience}[1][]{popbase,colframe=popblue,colback=popblueL,before upper={\popboxhead{Experiment}{popblue}{#1}}}
% « Le savais-tu ? » : carte violette pleine (contexte, culture, Cameroun)
\newtcolorbox{savaistu}[1][]{popbase,colback=poppurple,colframe=popink,
  fontupper=\popsemi\fontsize{10.4}{14.2}\selectfont\color{white},
  before upper={\poppill{Did you know?}{white}{poppurple}\ifx\relax#1\relax\else\hspace{6pt}{\popxbold\color{white}#1}\fi\par\vspace{7pt}}}
% Exercice résolu : énoncé en haut, \tcblower puis la solution sur fond crème
\newcounter{popexo}\newcounter{popexr}
\newtcolorbox{exoresolu}[1][]{popbase,colframe=poporange,colbacklower=popcream,
  segmentation style={popink,line width=1.2pt,dashed},
  before upper={\stepcounter{popexr}\popboxhead{Worked example \thepopexr}{poporange}{#1}},
  before lower={\poppill{Solution}{popink}{popcream}\par\vspace{6pt}}}
% Exercice (non résolu en place) — la correction est dans la partie Corrigés
\newtcolorbox{exercice}[1][]{popbase,colframe=popink,
  before upper={\stepcounter{popexo}\popboxhead{Exercise \thepopexo}{popink}{#1}}}
% Corrigé
\newtcolorbox{corrige}[1][]{popbase,colframe=popgreen,colback=popgreenL,
  before upper={\popboxhead{Solution}{popgreen}{#1}}}
% Objectifs / prérequis / bilan
\newtcolorbox{objectifs}{popbase,colframe=popink,colback=poporangeL,
  before upper={\popboxhead{Chapter objectives}{popink}{}}}
\newtcolorbox{prerequis}{popbase,colframe=popink,colback=popblueL,
  before upper={\popboxhead{Prerequisites}{popblue}{}}}
\newtcolorbox{bilan}{popbase,colframe=popink,colback=white,boxrule=2.8pt,
  before upper={\popboxhead{Fiche bilan}{poporange}{L'essentiel à connaître pour l'examen}}}
% Figure encadrée avec légende
\newtcolorbox{popfigure}[1][]{enhanced,breakable=false,colback=white,colframe=popline,boxrule=1.4pt,arc=8pt,
  left=10pt,right=10pt,top=10pt,bottom=8pt,before skip=10pt,after skip=12pt,halign=center,
  lower separated=false,halign lower=center,
  fontlower=\popsemi\fontsize{9}{11.5}\selectfont\color{popmuted},
  #1}
\newcommand{\legende}[1]{\par\vspace{6pt}{\popsemi\fontsize{9}{11.5}\selectfont\color{popmuted}#1}}

% ============================================================
% Signature OnBuch+
% ============================================================
\newcommand{\poplogoinline}[1]{{\poplogo\fontsize{#1}{#1}\selectfont\color{popink}OnBuch\color{poporange}+}}
\newcommand{\signatureonbuch}{%
  \begin{tcolorbox}[enhanced,colback=popink,colframe=popink,boxrule=2.2pt,arc=10pt,
      drop shadow=poporange,left=14pt,right=14pt,top=10pt,bottom=10pt,before skip=0pt,after skip=0pt]
    \tikz[baseline=-0.7ex]\node[circle,fill=popgreen,draw=popcream,line width=1.3pt,minimum size=22pt,inner sep=0]{\popcheck{white}};%
    \hspace{9pt}\begin{minipage}[c]{0.62\linewidth}
      {\popxbold\fontsize{12}{13}\selectfont\color{popcream}Signed\ \textbullet\ {\poplogo OnBuch\color{poporange}+}}\par\vspace{2pt}
      {\raggedright\popsemi\fontsize{8}{10.5}\selectfont\color{popline}\MakeUppercase{OnBuch+ teaching team}\\\MakeUppercase{Follows the official MINESEC syllabus}\par}
    \end{minipage}\hfill
    {\poplogo\fontsize{22}{22}\selectfont\color{popcream}OnBuch\color{poporange}+}
  \end{tcolorbox}%
}

% ============================================================
% Page de garde
% #1 titre · #2 matière · #3 niveau · #4 module · #5 n° leçon · #6 durée ·
% #7 description / objectifs (texte ou itemize)
% ============================================================
\newcommand{\pagegarde}[7]{%
  \thispagestyle{empty}
  \newgeometry{a4paper,margin=1.7cm,top=1.6cm,bottom=1.4cm}%
  \begin{tikzpicture}[remember picture,overlay]
    \fill[poporange!18] ([xshift=-3.2cm,yshift=-3.6cm]current page.north east) circle (4.6cm);
    \fill[poppurple!14] ([xshift=2.4cm,yshift=7.6cm]current page.south west) circle (3.8cm);
  \end{tikzpicture}%
  {\poplogo\fontsize{30}{30}\selectfont\color{popink}OnBuch\color{poporange}+}\hfill
  \raisebox{0.6ex}{\poppill{Signed course \textbullet\ Premium}{popink}{popcream}}\par
  \vspace{1pt}\popsquiggle{poppurple}\par
  \vspace{8pt}
  \begin{tcolorbox}[enhanced,colback=poporange,colframe=popink,boxrule=2.8pt,arc=13pt,
      drop shadow=popink,left=18pt,right=18pt,top=12pt,bottom=13pt,after skip=10pt]
    \poppill{#2 \textbullet\ #3}{popink}{popcream}\hspace{5pt}\poppill{#4}{white}{popink}\par\vspace{8pt}
    {\popdisplay\fontsize{14}{14}\selectfont\color{popink}CHAPTER #5}\par\vspace{4pt}
    {\raggedright\hyphenpenalty=10000\exhyphenpenalty=10000\popdisplay\fontsize{25}{27}\selectfont\color{popcream}\MakeUppercase{#1}\par}
    \vspace{8pt}
    {\popxbold\fontsize{9.5}{9.5}\selectfont\color{popink}\MakeUppercase{Suggested time: #6}}
  \end{tcolorbox}
  \begin{tcolorbox}[enhanced,colback=white,colframe=popink,boxrule=2.2pt,arc=10pt,
      drop shadow=popink,left=16pt,right=16pt,top=10pt,bottom=10pt,after skip=8pt]
    \poppill{In this chapter}{poppurple}{white}\par\vspace{5pt}
    \popsemi\fontsize{9.3}{12.6}\selectfont\color{popink}#7
  \end{tcolorbox}
  \noindent\begin{minipage}[t]{0.487\linewidth}
    \begin{tcolorbox}[enhanced,colback=popgreen,colframe=popink,boxrule=2.2pt,arc=10pt,
        drop shadow=popink,left=11pt,right=11pt,top=10pt,bottom=10pt,height=3.1cm,valign=center]
      \tcbox[colback=white,colframe=popink,boxrule=1.3pt,arc=5pt,boxsep=0pt,nobeforeafter,
        left=3pt,right=3pt,top=3pt,bottom=3pt,valign=center]{\qrcode[height=1.9cm]{https://play.google.com/store/apps/details?id=com.luvvix.onbuch&hl=en}}%
      \hspace{8pt}\begin{minipage}[c]{0.5\linewidth}
        {\popdisplay\fontsize{11}{12.5}\selectfont\color{white}\MakeUppercase{Download OnBuch}}\par\vspace{4pt}
        {\raggedright\popsemi\fontsize{8.4}{11}\selectfont\color{white}Free \textbullet\ Google Play\\AI tutor, past papers, results\par}
      \end{minipage}
    \end{tcolorbox}
  \end{minipage}\hfill
  \begin{minipage}[t]{0.487\linewidth}
    \begin{tcolorbox}[enhanced,colback=poppurple,colframe=popink,boxrule=2.2pt,arc=10pt,
        drop shadow=popink,left=11pt,right=11pt,top=10pt,bottom=10pt,height=3.1cm,valign=center]
      \tikz[baseline=-0.7ex]\node[circle,fill=white,draw=popink,line width=1.3pt,minimum size=1.6cm,inner sep=0]{\popdisplay\fontsize{24}{24}\selectfont\color{poppurple}+};%
      \hspace{8pt}\begin{minipage}[c]{0.6\linewidth}
        {\popdisplay\fontsize{11}{12.5}\selectfont\color{white}\MakeUppercase{Go OnBuch+}}\par\vspace{4pt}
        {\raggedright\popsemi\fontsize{8.4}{11}\selectfont\color{white}All signed courses, unlimited AI tutor, PDF sheets — OnBuch+ tab in the app\par}
      \end{minipage}
    \end{tcolorbox}
  \end{minipage}
  \vspace{8pt}
  \popcontenu
  \vfill
  \signatureonbuch
  \restoregeometry
  \newpage
}

% Rangée « Ce cours contient » (page de garde)
\newcommand{\poptile}[3]{% #1 couleur #2 gros texte #3 légende
  \begin{tcolorbox}[enhanced,colback=#1,colframe=popink,boxrule=2pt,arc=9pt,shadow={2.5pt}{-2.5pt}{0pt}{popink},
      left=6pt,right=6pt,top=9pt,bottom=9pt,halign=center,valign=center,height=1.75cm,width=\linewidth,nobeforeafter]
    {\popdisplay\fontsize{12.5}{13}\selectfont\color{white}\MakeUppercase{#2}}\par\vspace{3pt}
    {\centering\popsemi\fontsize{7.8}{9.5}\selectfont\color{white}#3\par}
  \end{tcolorbox}}
\newcommand{\popcontenu}{%
  \par\noindent{\popxboldls\fontsize{7.5}{7.5}\selectfont\color{popmuted}THIS SIGNED COURSE CONTAINS}\par\vspace{7pt}
  \noindent\begin{minipage}[t]{0.235\linewidth}\poptile{popblue}{Course}{complete, illustrated, step by step}\end{minipage}\hfill
  \begin{minipage}[t]{0.235\linewidth}\poptile{poporange}{Methods}{and worked examples}\end{minipage}\hfill
  \begin{minipage}[t]{0.235\linewidth}\poptile{poppink}{Exercises}{exam style + solutions}\end{minipage}\hfill
  \begin{minipage}[t]{0.235\linewidth}\poptile{popgreen}{Summary sheet}{the essentials to remember}\end{minipage}\par}

% Sommaire
\newtcolorbox{sommaire}{enhanced,colback=white,colframe=popink,boxrule=2.2pt,arc=10pt,
  drop shadow=popink,left=18pt,right=18pt,top=13pt,bottom=13pt,before skip=6pt,after skip=20pt,
  before upper={\poppill{Contents}{poppurple}{white}\par\vspace{9pt}\popsemi\fontsize{10.6}{14.5}\selectfont\color{popink}}}

% Signature de fin de cours — poussée tout en bas de la dernière page
\newcommand{\findecours}{%
  \par\vfill\nopagebreak
  \begin{center}\popsquiggle{poporange}\end{center}\vspace{4pt}
  \signatureonbuch
}
\AtBeginDocument{\def\llbracket{\mathopen{[\mkern-3.5mu[}}\def\rrbracket{\mathclose{]\mkern-3.5mu]}}} % intervalles d'entiers (glyphe ⟦ absent de la police)
% Glyphes absents de Fira Math : redéfinis à partir de glyphes disponibles
\AtBeginDocument{%
  \def\setminus{\mathbin{\backslash}}\let\smallsetminus\setminus
  \def\vdots{\mathord{\vbox{\baselineskip3.2pt\lineskiplimit0pt\kern4pt\hbox{.}\hbox{.}\hbox{.}}}}%
  \def\ddots{\mathinner{\mkern1mu\raise6.5pt\hbox{.}\mkern2mu\raise3.6pt\hbox{.}\mkern2mu\raise0.7pt\hbox{.}\mkern1mu}}%
  \def\triangle{\mathord{\scalebox{0.9}{$\symup{\Delta}$}}}%
  \def\nabla{\mathord{\raisebox{\depth}{\scalebox{1}[-1]{$\symup{\Delta}$}}}}%
  \def\checkmark{\popcheck{popdark}}%
  \def\star{\mathbin{\ast}}\def\bigstar{\mathord{\ast}}%
  \def\diamond{\mathbin{\scalebox{0.6}{$\Diamond$}}}%
  \def\bigcup{\mathop{\mathchoice{\raisebox{-0.3em}{\scalebox{1.7}{$\cup$}}}{\scalebox{1.25}{$\cup$}}{\cup}{\cup}}\displaylimits}%
  \def\bigcap{\mathop{\mathchoice{\raisebox{-0.3em}{\scalebox{1.7}{$\cap$}}}{\scalebox{1.25}{$\cap$}}{\cap}{\cap}}\displaylimits}%
  \def\lesseqqgtr{\mathrel{\lessgtr}}\def\bigoplus{\mathop{\oplus}\displaylimits}%
}
\providecommand{\ssi}{if and only if} % abréviation fréquente
\AtBeginDocument{\providecommand{\wideparen}{\overparen}} % arc de cercle

%% FIGFIX-BEGIN v3 — correctifs globaux des figures (docs/onbuch-generation/tools/figfix)
\usepackage{adjustbox}\usepackage{etoolbox}
% toute figure TikZ/pgfplots plus large que la ligne est réduite à la largeur disponible
\BeforeBeginEnvironment{tikzpicture}{\begin{adjustbox}{max width=\linewidth}}
\AfterEndEnvironment{tikzpicture}{\end{adjustbox}}
% légendes pgfplots lisibles par-dessus les courbes
\pgfplotsset{every axis/.append style={legend style={fill=white,fill opacity=0.92,text opacity=1,draw=black!15,inner sep=3pt}}}
% étiquettes d'arêtes (syntaxe edge["..."]) avec fond blanc
\tikzset{every edge quotes/.append style={fill=white,inner sep=1.2pt}}
% bulles de cartes mentales : texte à la ligne dans la bulle, bulle un peu plus grande
\tikzset{every concept/.append style={text width=2.0cm,align=center,minimum size=2.3cm,inner sep=2pt}}
%% FIGFIX-END
\begin{document}

\pagegarde{Capacitors}{Physics}{Upper Sixth}{Upper Sixth — Physics}{5}{4 periods}{This chapter studies capacitors: how they store charge and energy in an electric field, how their capacitance depends on geometry and dielectric materials, and how they behave in series and parallel combinations. It ends with the time-dependent charging and discharging of capacitors through resistors, a cornerstone topic of the GCE Advanced Level Physics examination.
\vspace{4pt}
\begin{multicols}{2}\small
\begin{itemize}
\item By the end of this chapter I can define electric potential and potential difference and compute the work done in moving a charge in a uniform electric field.
\item By the end of this chapter I can describe an electric dipole in a uniform field and calculate the torque acting on it.
\item By the end of this chapter I can identify capacitors by their circuit symbols and physical appearance, and state the factors affecting capacitance.
\item By the end of this chapter I can define capacitance, relative permittivity, and use C = ε₀εᵣA/d for a parallel-plate capacitor.
\item By the end of this chapter I can measure capacitance experimentally using a charge–voltage method.
\item By the end of this chapter I can derive and apply the equivalent capacitance of series and parallel combinations.
\end{itemize}\end{multicols}}

\begin{sommaire}
\begin{multicols}{2}
\begin{enumerate}
\item Electric potential, work done and the electric dipole
\item Capacitors, capacitance and permittivity
\item Combinations of capacitors and energy stored
\item Charging and discharging through a resistor: the time constant
\item Integration activity
\item Methods and worked examples
\item Exercises
\item Solutions to the exercises
\item Summary sheet
\end{enumerate}
\end{multicols}
\end{sommaire}

\begin{objectifs}
\begin{itemize}
\item By the end of this chapter I can define electric potential and potential difference and compute the work done in moving a charge in a uniform electric field.
\item By the end of this chapter I can describe an electric dipole in a uniform field and calculate the torque acting on it.
\item By the end of this chapter I can identify capacitors by their circuit symbols and physical appearance, and state the factors affecting capacitance.
\item By the end of this chapter I can define capacitance, relative permittivity, and use C = ε₀εᵣA/d for a parallel-plate capacitor.
\item By the end of this chapter I can measure capacitance experimentally using a charge–voltage method.
\item By the end of this chapter I can derive and apply the equivalent capacitance of series and parallel combinations.
\item By the end of this chapter I can calculate the charge and energy stored in any combination of capacitors connected to a d.c. supply.
\item By the end of this chapter I can describe and sketch the charging and discharging of a capacitor through a resistor and define the time constant τ = RC.
\item By the end of this chapter I can use the exponential equations of charging and discharging and interpret them at t = 0 and as t → ∞.
\end{itemize}
\end{objectifs}

\begin{prerequis}
\begin{itemize}
\item Electric charge, Coulomb's law and the concept of electric field (Lower Sixth, electrostatics)
\item Uniform electric field between parallel plates, E = V/d (Lower Sixth)
\item Current, voltage, resistance and Ohm's law; simple d.c. circuits (Lower Sixth / O Level)
\item Exponential and natural logarithm functions (Upper Sixth Mathematics)
\item Work and energy concepts from mechanics (Lower Sixth)
\end{itemize}
\end{prerequis}

\newpage

\coursec{Electric potential, work done and the electric dipole}

During the long dry-season evenings in Bamenda or Yaoundé, many families rely on rechargeable lamps and phone chargers that keep shining for a while even after the mains power from ENEO cuts off. Inside these devices, a component called a \cle{capacitor} quietly stores electric energy and releases it when needed. The same principle explains why a camera flash can fire instantly, why a defibrillator can restart a heart, and why unplugging a television does not make its circuits safe to touch immediately. How much charge can such a device store, how fast does it fill up, and how quickly does it empty? This chapter answers these questions quantitatively.

Before we can describe how a capacitor stores energy, we need a precise way of talking about the \emph{energy per unit charge} at a point in an electric field. That quantity is the \cle{electric potential}. This first section therefore builds the essential vocabulary: potential, potential difference, the work done in moving a charge, and the behaviour of an electric dipole in a field.

\courssub{Electric potential at a point}

\textbf{Intuition.} Imagine a hilly landscape. A ball placed on a slope "wants" to roll downhill: the height tells you how much gravitational energy each kilogram of the ball possesses. In the same way, an electric field is a kind of invisible landscape for charges. A positive charge placed in the field "wants" to move in a certain direction, and the \emph{electric potential} tells you how much electrical energy each coulomb of charge possesses at that point. Just as we measure height relative to sea level, we measure potential relative to a reference point, chosen to be \emph{infinity}, where the potential is taken as zero.

\begin{definition}[Electric potential]
The \cle{electric potential} $V$ at a point in an electric field is the \textbf{work done per unit positive charge} in bringing a small positive test charge from infinity to that point, without accelerating it and without disturbing the field:
\[
V = \frac{W_{\infty \to P}}{q}.
\]
The SI unit of potential is the \cle{volt} (V): $1\ \mathrm{V} = 1\ \mathrm{J\,C^{-1}}$. Potential is a \textbf{scalar} quantity.
\end{definition}

\begin{definition}[Potential difference]
The \cle{potential difference} (p.d.) between two points $A$ and $B$ is the work done per unit positive charge in moving a charge from $A$ to $B$:
\[
V_{AB} = V_B - V_A = \frac{W_{A\to B}}{q}.
\]
It is also measured in volts. In practice, all circuit measurements are potential \emph{differences}; a voltmeter connected between two points reads $V_{AB}$.
\end{definition}

\begin{attention}[Infinity is just a reference]
Saying "the potential at $P$ is $+50\ \mathrm{V}$" means: it takes $50\ \mathrm{J}$ of work per coulomb to bring a positive charge from infinity to $P$ (against the field of a positive source charge, for example). Only potential \emph{differences} have physical consequences; the choice of zero is a convention, exactly like choosing sea level for altitudes.
\end{attention}

\begin{savaistu}[Potential in Cameroon's high-voltage lines]
The high-voltage transmission lines you see crossing the countryside (e.g., from Song Loulou to Yaoundé) operate at potentials of $225\ \mathrm{kV}$ or $90\ \mathrm{kV}$ relative to earth. This high potential reduces the current for a given power, minimising $I^2R$ losses over hundreds of kilometres. The "potential" of the line is always a potential \emph{difference} with respect to the ground.
\end{savaistu}

\courssub{Work done in moving a charge: $W = qV$}

From the definition above, the work done in moving a charge $q$ through a potential difference $V$ follows immediately.

\begin{propriete}[Work done and potential difference]
When a charge $q$ moves from a point $A$ to a point $B$ between which the p.d. is $V = V_B - V_A$, the work done \textbf{by the electric field} is
\[
W_{\text{field}} = q(V_A - V_B) = -qV,
\]
while the work that must be supplied \textbf{by an external agent} to move the charge slowly (without change of kinetic energy) is
\[
W_{\text{ext}} = q(V_B - V_A) = qV.
\]
A positive charge released freely in a field accelerates from \textbf{high} potential to \textbf{low} potential, converting electrical potential energy into kinetic energy. A negative charge does the opposite.
\end{propriete}

\begin{attention}[Work done BY the field versus BY an external agent]
This is the most frequent source of sign errors at the GCE. Ask yourself: \emph{who is doing the work?}
\begin{itemize}
\item If the field pushes the charge (charge moving the way the field "wants" it to go), the work done \textbf{by the field} is \textbf{positive} and the charge speeds up.
\item If you push the charge \emph{against} the field, the work done \textbf{by you} (the external agent) is positive, and the work done by the field is negative.
\end{itemize}
Always write down which work you are computing before substituting signs for $q$ and $V$.
\end{attention}

\begin{exemplebox}[Moving a charge between two points]
A charge $q = +2.0\ \mathrm{\mu C}$ is moved slowly from a point $A$ where $V_A = 100\ \mathrm{V}$ to a point $B$ where $V_B = 250\ \mathrm{V}$. The p.d. is $V = V_B - V_A = 150\ \mathrm{V}$. The external agent must supply
\[
W_{\text{ext}} = qV = 2.0\times 10^{-6} \times 150 = 3.0\times 10^{-4}\ \mathrm{J},
\]
while the field does work $W_{\text{field}} = -3.0\times 10^{-4}\ \mathrm{J}$ (it opposes the motion, since a positive charge is being pushed towards higher potential).
\end{exemplebox}

\courssub{Field and potential: $E = V/d$ in a uniform field}

Consider two large parallel metal plates separated by a distance $d$, one at potential $V$ and the other at potential $0$. Between them the field is \textbf{uniform}: the field lines are parallel and equally spaced. Move a charge $q$ from the low-potential plate to the high-potential plate, a distance $d$ against the field. The external force needed is $F = qE$ (constant, since $E$ is uniform), so the work supplied is
\[
W_{\text{ext}} = F\,d = qEd.
\]
But by definition this same work equals $qV$. Hence $qEd = qV$, which gives the key result.

\begin{propriete}[Relation between field and potential in a uniform field]
Between two points separated by a distance $d$ measured \textbf{along the field direction} in a uniform field,
\[
\boxed{E = \frac{V}{d}}
\]
The electric field points from \textbf{high} potential to \textbf{low} potential. This relation gives an alternative SI unit for $E$: the volt per metre, $\mathrm{V\,m^{-1}}$, which is identical to $\mathrm{N\,C^{-1}}$.
\end{propriete}

More generally, the field is related to how \emph{rapidly} the potential changes with position: $E = -\dfrac{dV}{dx}$. The minus sign says the field points in the direction of \emph{decreasing} potential, just as a ball rolls downhill.

\begin{popfigure}
\begin{tikzpicture}[scale=1.0]
% plates
\fill[popblueL] (0,-2.2) rectangle (0.25,2.2);
\fill[poppink!40] (6.75,-2.2) rectangle (7,2.2);
\node[above] at (0.12,2.2) {$+Q$ plate, $V$};
\node[above] at (6.87,2.2) {$-Q$ plate, $0$};
% field lines
\foreach \y in {-1.5,-0.75,0,0.75,1.5}{
  \draw[-latex,popblue,thick] (0.4,\y) -- (6.6,\y);
}
\node[popblue] at (3.5,1.85) {$\vec{E}$ (uniform)};
% equipotentials
\foreach \x/\v in {1.75/{}, 3.5/{}, 5.25/{}}{
  \draw[dashed,popdark] (\x,-2.0) -- (\x,2.0);
}
\node[popdark,below] at (1.75,-2.0) {$\tfrac{3V}{4}$};
\node[popdark,below] at (3.5,-2.0) {$\tfrac{V}{2}$};
\node[popdark,below] at (5.25,-2.0) {$\tfrac{V}{4}$};
% charge path
\draw[-latex,poporange,very thick] (1.0,-1.7) .. controls (2.5,-0.9) and (4.0,-1.3) .. (5.9,-0.4);
\fill[poporange] (1.0,-1.7) circle (0.09);
\node[poporange,below left] at (1.0,-1.7) {$q$ at $A$};
\fill[poporange] (5.9,-0.4) circle (0.09);
\node[poporange,right] at (5.9,-0.4) {$B$};
% d arrow
\draw[<->,popmuted] (0.12,-2.5) -- (6.87,-2.5);
\node[popmuted,below] at (3.5,-2.5) {$d$};
\end{tikzpicture}
\legende{Uniform field between parallel plates. Solid arrows: field lines (high to low potential). Dashed lines: equipotentials, everywhere perpendicular to the field. The work done on charge $q$ depends only on the potentials at $A$ and $B$, not on the path taken.}
\end{popfigure}

\begin{popfigure}
\begin{center}
\begin{tikzpicture}
\begin{axis}[
  width=0.75\linewidth, height=6cm,
  xlabel={distance $x$ from the positive plate (m)},
  ylabel={potential $V$ (V)},
  xmin=0, xmax=0.11, ymin=0, ymax=110,
  xtick={0,0.02,0.04,0.06,0.08,0.1},
  ytick={0,20,40,60,80,100},
  grid=both, grid style={popline!40},
]
\addplot[popblue, very thick, domain=0:0.1] {100 - 1000*x};
\addplot[mark=*, poporange, only marks] coordinates {(0.04,60)};
\node[popdark, anchor=west] at (axis cs:0.045,60) {$V = V_0 - Ex$};
\end{axis}
\end{tikzpicture}
\end{center}
\legende{Potential versus position between parallel plates ($V_0 = 100\ \mathrm{V}$, $d = 0.10\ \mathrm{m}$, so $E = 1000\ \mathrm{V\,m^{-1}}$). The potential falls \emph{linearly}; the magnitude of the gradient is the field strength.}
\end{popfigure}

\begin{experience}[Mapping equipotentials (laboratory)]
\textbf{Apparatus:} Conductive paper (teledeltos), power supply ($0-12\ \mathrm{V}$), two metal electrodes (bars or points), voltmeter, connecting wires.
\textbf{Procedure:}
\begin{enumerate}
\item Place the conductive paper on a flat board. Fix the two electrodes on the paper and connect them to the power supply.
\item Set the p.d. between electrodes to about $6\ \mathrm{V}$.
\item Use the voltmeter probe (negative lead fixed on one electrode) to find points on the paper where the reading is $1\ \mathrm{V}$, $2\ \mathrm{V}$, $3\ \mathrm{V}$, etc. Mark these points.
\item Join points of equal potential to draw equipotential lines. Draw field lines perpendicular to them.
\end{enumerate}
\textbf{Observation:} For parallel bar electrodes, equipotentials are straight, parallel, equally spaced lines (uniform field). For point electrodes, they are curved, spacing increases with distance (radial field).
\end{experience}

\begin{aretenir}[Equipotential surfaces]
An \cle{equipotential surface} is a surface on which the potential is the same everywhere.
\begin{itemize}
\item No work is done in moving a charge \emph{along} an equipotential ($W = q\Delta V = 0$).
\item Field lines always cross equipotentials at \textbf{right angles}.
\item Between parallel plates, the equipotentials are planes parallel to the plates and equally spaced (uniform field).
\item Around a point charge, the equipotentials are concentric spheres, getting farther apart as the field weakens.
\end{itemize}
\end{aretenir}

\courssub{Kinetic energy gained by an accelerated charge}

When a charge is \emph{free} to move, energy conservation replaces the external agent: the electrical potential energy lost becomes kinetic energy. This is the working principle of the \cle{electron gun} inside old television tubes, oscilloscopes and X-ray tubes, and it reappears constantly in later chapters (motion of charges in fields, particle accelerators).

\begin{methode}[Solving acceleration-of-charge problems with $\tfrac{1}{2}mv^2 = qV$]
\begin{enumerate}
\item Identify the charge $q$ (magnitude \emph{and} sign) and the p.d. $V$ through which it is accelerated. Use the \emph{magnitude} of the p.d. that causes the speed-up.
\item State energy conservation: loss of electrical potential energy = gain of kinetic energy:
\[
\frac{1}{2}mv^2 - \frac{1}{2}mv_0^2 = |q|\,V.
\]
If the particle starts from rest, $v_0 = 0$.
\item Solve for the unknown: $v = \sqrt{\dfrac{2|q|V}{m}}$.
\item Check: every numerical answer must carry its SI unit and sensible significant figures.
\end{enumerate}
\end{methode}

\begin{exoresolu}[Electron accelerated through 2.0 kV]
An electron (charge $e = 1.60\times 10^{-19}\ \mathrm{C}$, mass $m_e = 9.11\times 10^{-31}\ \mathrm{kg}$) is emitted from rest at the cathode of an electron gun and accelerated through a p.d. of $2.0\ \mathrm{kV}$. Calculate (a) the kinetic energy gained, in joules and in electronvolts; (b) the final speed of the electron.
\tcblower
\textbf{(a)} The electron, being negative, is accelerated from low to high potential; the magnitude of the p.d. is $V = 2.0\times 10^{3}\ \mathrm{V}$. Energy conservation gives
\[
E_k = eV = 1.60\times 10^{-19} \times 2.0\times 10^{3} = 3.2\times 10^{-16}\ \mathrm{J}.
\]
By definition of the \cle{electronvolt} ($1\ \mathrm{eV} = e \times 1\ \mathrm{V} = 1.60\times 10^{-19}\ \mathrm{J}$), this is simply $E_k = 2.0\ \mathrm{keV}$.

\textbf{(b)} Since the electron starts from rest,
\[
\frac{1}{2}m_e v^2 = eV
\quad\Longrightarrow\quad
v = \sqrt{\frac{2eV}{m_e}} = \sqrt{\frac{2\times 3.2\times 10^{-16}}{9.11\times 10^{-31}}}.
\]
\[
v = \sqrt{7.03\times 10^{14}} = 2.7\times 10^{7}\ \mathrm{m\,s^{-1}}.
\]
This is about $9\%$ of the speed of light: electron guns produce genuinely fast particles from modest voltages.
\end{exoresolu}

\courssub{The electric dipole}

Many molecules, such as water, are electrically neutral overall yet behave as if a small positive charge and a small negative charge sit at slightly different positions. Such a system is an \cle{electric dipole}, and understanding it explains why charged objects attract neutral pieces of paper and why some liquids rotate the direction of an applied field.

\begin{definition}[Electric dipole and dipole moment]
An \cle{electric dipole} consists of two equal and opposite point charges $+q$ and $-q$ separated by a fixed distance $2a$ (the \emph{separation}). Its \cle{electric dipole moment} is the vector
\[
\vec{p} = q\,\vec{d},
\qquad p = q \times (\text{separation}) = 2aq,
\]
directed from the \textbf{negative} charge to the \textbf{positive} charge. The SI unit of $p$ is the coulomb-metre ($\mathrm{C\,m}$).
\end{definition}

\courssub{Torque on a dipole in a uniform field}

Place the dipole in a uniform field $\vec{E}$, with its axis making an angle $\theta$ with the field. The force on $+q$ is $\vec{F}_+ = q\vec{E}$ (along the field) and the force on $-q$ is $\vec{F}_- = -q\vec{E}$ (opposite the field). The two forces are equal, opposite and parallel but act at different points: they form a \textbf{couple}. The net force is zero (the dipole does not translate), but there is a net \textbf{torque}.

Taking moments about the centre of the dipole: each force has magnitude $qE$ and lever arm $a\sin\theta$, and both moments act in the same sense of rotation, so
\[
\tau = 2\times qE\,a\sin\theta = (2aq)E\sin\theta.
\]

\begin{propriete}[Torque on an electric dipole — proof examinable]
A dipole of moment $p$ placed at an angle $\theta$ to a uniform field $E$ experiences a torque
\[
\boxed{\tau = pE\sin\theta}
\]
which tends to \textbf{align the dipole with the field}. In vector form, $\vec{\tau} = \vec{p}\times\vec{E}$.
\begin{itemize}
\item $\theta = 0$ (dipole parallel to $\vec{E}$): $\tau = 0$, \textbf{stable equilibrium}. A small displacement produces a restoring torque.
\item $\theta = 180^{\circ}$ (dipole antiparallel): $\tau = 0$, \textbf{unstable equilibrium}. The slightest disturbance flips the dipole round.
\item $\theta = 90^{\circ}$: maximum torque, $\tau_{\max} = pE$.
\end{itemize}
\end{propriete}

\begin{popfigure}
\begin{tikzpicture}[scale=1.1]
% field
\foreach \y in {-1.4,-0.7,0,0.7,1.4}{
  \draw[-latex,popblue!70] (-0.4,\y) -- (5.4,\y);
}
\node[popblue] at (5.0,1.7) {$\vec{E}$};
% dipole axis at angle theta
\coordinate (C) at (2.5,0);
\coordinate (P) at (3.9,1.12);
\coordinate (N) at (1.1,-1.12);
\draw[popdark, thick] (N) -- (P);
\fill[popgreen] (P) circle (0.14);
\node[white] at (P) {\tiny $+$};
\node[popdark, above right=1pt] at (P) {$+q$};
\fill[popink] (N) circle (0.14);
\node[white] at (N) {\tiny $-$};
\node[popdark, below left=1pt] at (N) {$-q$};
% forces
\draw[-latex,poporange,very thick] (P) -- ++(1.3,0) node[right] {$\vec{F}_+ = q\vec{E}$};
\draw[-latex,poporange,very thick] (N) -- ++(-1.3,0) node[left] {$\vec{F}_- = -q\vec{E}$};
% angle
\draw[popmuted] (C) ++(1.6,0) arc (0:35:1.6);
\node[popmuted] at (3.9,0.28) {$\theta$};
\draw[dashed,popmuted] (1.0,0) -- (4.4,0);
% centre and moment
\fill (C) circle (0.05);
\node[below] at (2.5,-0.08) {centre};
\draw[-latex,popgreen,thick] (1.35,-0.95) -- (3.65,0.95);
\node[popgreen] at (2.0,0.75) {$\vec{p}$};
% torque annotation
\node[popdark, align=center, text width=4.2cm] at (2.5,-2.1) {couple $\Rightarrow$ torque\\ $\tau = pE\sin\theta$ (into page)};
\end{tikzpicture}
\legende{A dipole at angle $\theta$ in a uniform field. The forces $q\vec{E}$ and $-q\vec{E}$ form a couple: zero net force, but a torque $\tau = pE\sin\theta$ that rotates the dipole towards alignment with $\vec{E}$.}
\end{popfigure}

\begin{exoresolu}[Torque on a water-like dipole]
A molecule is modelled as a dipole with charges $\pm 1.0\times 10^{-20}\ \mathrm{C}$ separated by $4.0\times 10^{-10}\ \mathrm{m}$. It sits in a uniform field $E = 5.0\times 10^{5}\ \mathrm{V\,m^{-1}}$ making an angle of $30^{\circ}$ with the field. Find (a) the dipole moment; (b) the torque on the molecule; (c) the maximum possible torque.
\tcblower
\textbf{(a)} $p = q \times \text{separation} = 1.0\times 10^{-20}\times 4.0\times 10^{-10} = 4.0\times 10^{-30}\ \mathrm{C\,m}$.

\textbf{(b)} $\tau = pE\sin\theta = 4.0\times 10^{-30}\times 5.0\times 10^{5}\times \sin 30^{\circ} = 1.0\times 10^{-24}\ \mathrm{N\,m}$.

\textbf{(c)} $\tau_{\max} = pE = 2.0\times 10^{-24}\ \mathrm{N\,m}$, occurring at $\theta = 90^{\circ}$.
\end{exoresolu}

\textbf{Extension — work done in rotating a dipole.} To rotate a dipole slowly from alignment ($\theta = 0$) to an angle $\theta$, an external agent must work against the restoring torque. The work required is
\[
W = \int_0^{\theta} pE\sin\theta'\,d\theta' = pE(1 - \cos\theta),
\]
so the dipole has its \emph{lowest} potential energy when aligned with the field (stable equilibrium) and its highest when antiparallel. This explains why polar molecules in a liquid align with an applied field, an idea we shall meet again when we study dielectrics inside capacitors. (This integral treatment is an extension: at GCE level, only the qualitative conclusion is required.)

\begin{exercice}[Practice]
\begin{enumerate}
\item Two parallel plates are $5.0\ \mathrm{mm}$ apart and the p.d. between them is $200\ \mathrm{V}$. Calculate the electric field strength between the plates.
\item A proton ($q = +e$, $m = 1.67\times 10^{-27}\ \mathrm{kg}$) is accelerated from rest through $5.0\ \mathrm{kV}$. Find its final kinetic energy and speed.
\item A dipole of moment $2.0\times 10^{-29}\ \mathrm{C\,m}$ lies in a field of $1.0\times 10^{6}\ \mathrm{V\,m^{-1}}$. At what angle is the torque half its maximum value?
\end{enumerate}
\end{exercice}

\begin{corrige}[Exercise 1]
\textbf{1.} $E = V/d = 200/(5.0\times 10^{-3}) = 4.0\times 10^{4}\ \mathrm{V\,m^{-1}}$.

\textbf{2.} $E_k = eV = 1.60\times 10^{-19}\times 5.0\times 10^{3} = 8.0\times 10^{-16}\ \mathrm{J}$; then $v = \sqrt{2E_k/m} = \sqrt{1.6\times 10^{-15}/1.67\times 10^{-27}} = 9.8\times 10^{5}\ \mathrm{m\,s^{-1}}$.

\textbf{3.} $\tau_{\max} = pE$, and $\tau = pE\sin\theta = \tfrac{1}{2}pE$ gives $\sin\theta = 0.5$, so $\theta = 30^{\circ}$ (or $150^{\circ}$).
\end{corrige}

\begin{aretenir}[Key points of this section]
\begin{itemize}
\item Potential $V$ = work per unit charge brought from infinity; unit: the volt ($1\ \mathrm{V} = 1\ \mathrm{J\,C^{-1}}$).
\item Work done in moving $q$ through a p.d.: $W = qV$; watch the sign conventions (field vs external agent).
\item Uniform field: $E = V/d$, directed from high to low potential; equipotentials are perpendicular to field lines.
\item Freely accelerated charge: $\tfrac{1}{2}mv^2 = qV$ — the energy method behind electron guns.
\item Dipole: $p = q\times$separation, directed from $-$ to $+$; torque $\tau = pE\sin\theta$; stable equilibrium when aligned with $\vec{E}$.
\end{itemize}
\end{aretenir}

With the language of potential and energy now in place, we are ready for the central object of this chapter: the capacitor, a device engineered precisely to separate charge and store energy in an electric field.

\coursec{Capacitors, capacitance and permittivity}

In the previous section we saw how electric potential allows us to quantify the work done when a charge moves in an electric field. We now turn to a device that \emph{stores} electric charge and energy by maintaining a potential difference between two conductors: the capacitor. From the flash unit of a camera in a Yaoundé photography studio to the smoothing circuit inside your phone charger, capacitors are everywhere in modern electronics.

\courssub{What is a capacitor?}

A \cle{capacitor} consists of two conductors (called \cle{plates}) separated by an insulating material known as a \cle{dielectric}. When a voltage source is connected across the plates, one plate acquires a charge $+Q$ and the other $-Q$; the net charge of the capacitor as a whole remains zero. The dielectric prevents charge from flowing directly between the plates, so the charge is stored on the surfaces facing each other.

\begin{popfigure}
\begin{tikzpicture}[scale=1.2]
% Plates
\draw[popblue, thick, fill=popblueL] (0,0) rectangle (4,0.3);
\draw[popblue, thick, fill=popblueL] (0,1.5) rectangle (4,1.8);
% Dielectric
\draw[popgreen, thick, fill=popgreenL, opacity=0.5] (0,0.3) rectangle (4,1.5);
% Labels
\node[align=center] at (2, -0.5) {Plate 1 ($-Q$)};
\node[align=center] at (2, 2.1) {Plate 2 ($+Q$)};
\node[align=center, popgreen] at (2, 0.9) {Dielectric ($\varepsilon_r$)} ;
% Dimensions
\draw[<->, popdark] (4.2,0) -- (4.2,0.3) node[midway, right] {$d$};
\draw[<->, popdark] (0,-0.3) -- (4,-0.3) node[midway, below] {$A$};
% Field lines
\foreach \y in {0.5, 1.0, 1.3} {
    \draw[->, poporange, thick] (2, \y+0.15) -- (2, \y-0.15);
}
\end{tikzpicture}
\legende{Parallel-plate capacitor: two conducting plates of area $A$ separated by distance $d$, with a dielectric of relative permittivity $\varepsilon_r$. The electric field is uniform between the plates (fringing ignored).}
\end{popfigure}

Everyday examples abound:
\begin{itemize}
    \item \textbf{Camera flash:} A capacitor charged to several hundred volts discharges rapidly through a xenon tube, producing an intense burst of light.
    \item \textbf{Radio tuning:} A variable capacitor in the LC circuit selects the desired frequency.
    \item \textbf{Phone charger / power supply:} Electrolytic capacitors smooth the rectified AC voltage, reducing ripple.
    \item \textbf{Power factor correction:} Large capacitors in industrial installations (e.g., at the ENEO substation in Douala) improve the power factor of inductive loads.
\end{itemize}

\courssub{Circuit symbols and identification}

Capacitors are represented in circuit diagrams by specific symbols. It is essential to recognise them and to identify real components by their markings.

\begin{popfigure}
\begin{tikzpicture}[scale=1.2, european]
% Fixed capacitor
\draw (0,0) to[C, l={$C$}] (2,0);
\node at (1, -0.8) {(a) Fixed capacitor};
% Polarised (electrolytic) capacitor
\draw (3,0) to[pC, l={$C$}, l_=+] (5,0);
\node at (4, -0.8) {(b) Polarised (electrolytic)};
% Variable capacitor
\draw (6,0) to[vC, l={$C$}] (8,0);
\node at (7, -0.8) {(c) Variable capacitor};
\end{tikzpicture}
\legende{Standard circuit symbols for capacitors (European style). The polarised capacitor has a `+' mark on the positive terminal; the variable capacitor is shown with an arrow through the plates.}
\end{popfigure}

\textbf{Identification of real components:}
\begin{itemize}
    \item \textbf{Fixed capacitors (ceramic, film, mica):} Usually marked with a three-digit code: the first two digits are the significant figures, the third is the multiplier (power of ten) in picofarads. Example: \texttt{104} means $10 \times 10^4\ \mathrm{pF} = 100\ \mathrm{nF} = 0.1\ \mu\mathrm{F}$. A letter may indicate tolerance (J = $\pm5\%$, K = $\pm10\%$, M = $\pm20\%$).
    \item \textbf{Electrolytic capacitors:} Cylindrical with a clearly marked negative stripe (or positive `+' sign). The capacitance (in $\mu\mathrm{F}$) and voltage rating (e.g., $470\ \mu\mathrm{F}\ 25\mathrm{V}$) are printed on the body. \emph{Warning:} Reversing the polarity can cause the capacitor to overheat and explode.
    \item \textbf{Variable capacitors:} Air-spaced or plastic-film types with a rotating shaft; used in tuning circuits. Capacitance varies typically from a few pF to a few hundred pF.
\end{itemize}

\begin{attention}[Polarity of electrolytic capacitors]
Always observe the polarity markings when connecting an electrolytic capacitor in a DC circuit. Connecting it backwards can lead to catastrophic failure (venting, explosion). In AC or signal circuits, use non-polarised types (ceramic, film, or bipolar electrolytic).
\end{attention}

\courssub{Definition of capacitance and the farad}

The ability of a capacitor to store charge per unit voltage is its \cle{capacitance}.

\begin{definition}[Capacitance]
The capacitance $C$ of a capacitor is defined as the ratio of the magnitude of charge $Q$ on \emph{one} plate to the potential difference $V$ between the plates:
\[
C = \frac{Q}{V}
\]
The SI unit of capacitance is the \cle{farad} (symbol F): $1\ \mathrm{F} = 1\ \mathrm{C\,V^{-1}}$.
\end{definition}

\begin{attention}[Charge $Q$ in $C = Q/V$]
$Q$ is the charge on \textbf{one} plate only (the positive plate carries $+Q$, the negative plate $-Q$). It is \textbf{not} the sum of the charges on both plates (which would be zero). This is a frequent source of error in examinations.
\end{attention}

The farad is a very large unit. Practical capacitors use submultiples:
\[
1\ \mu\mathrm{F} = 10^{-6}\ \mathrm{F},\quad
1\ \mathrm{nF} = 10^{-9}\ \mathrm{F},\quad
1\ \mathrm{pF} = 10^{-12}\ \mathrm{F}.
\]

\begin{aretenir}[Unit conversion]
Always convert $\mu\mathrm{F}$, $\mathrm{nF}$ and $\mathrm{pF}$ to farads before substituting into formulas. For example, $4.7\ \mu\mathrm{F} = 4.7 \times 10^{-6}\ \mathrm{F}$. Forgetting this conversion is one of the most common mistakes in GCE Advanced Level papers.
\end{aretenir}

\courssub{The parallel-plate capacitor: factors affecting capacitance}

The simplest capacitor to analyse is the parallel-plate capacitor: two large conducting plates of area $A$, separated by a small distance $d$ (so that edge effects are negligible). Let the dielectric between the plates have relative permittivity $\varepsilon_r$.

\begin{propriete}[Derivation of capacitance for a parallel-plate capacitor (examinable)]
Using Gauss's law, consider a Gaussian pillbox that encloses part of the positive plate. The electric flux passes only through the face inside the capacitor (field is zero inside the conductor and negligible outside). If $\sigma = Q/A$ is the surface charge density:
\[
\oint \vec{E} \cdot d\vec{A} = \frac{Q_{\text{enc}}}{\varepsilon_0 \varepsilon_r} \quad\Longrightarrow\quad E A = \frac{\sigma A}{\varepsilon_0 \varepsilon_r}
\]
Hence the uniform electric field between the plates is
\[
E = \frac{\sigma}{\varepsilon_0 \varepsilon_r} = \frac{Q}{\varepsilon_0 \varepsilon_r A}.
\]
The potential difference is $V = E d$ (since the field is uniform). Therefore
\[
V = \frac{Q d}{\varepsilon_0 \varepsilon_r A} \quad\Longrightarrow\quad
C = \frac{Q}{V} = \frac{\varepsilon_0 \varepsilon_r A}{d}.
\]
\end{propriete}

\begin{definition}[Relative permittivity (dielectric constant)]
The \cle{relative permittivity} $\varepsilon_r$ of a material is defined as the ratio of the capacitance of a capacitor with that material as dielectric to the capacitance of the same capacitor with vacuum (or air, to a very good approximation) as dielectric:
\[
\varepsilon_r = \frac{C_{\text{dielectric}}}{C_{\text{vacuum}}}.
\]
For a parallel-plate capacitor,
\[
C = \frac{\varepsilon_0 \varepsilon_r A}{d},
\]
where $\varepsilon_0 = 8.85 \times 10^{-12}\ \mathrm{F\,m^{-1}}$ is the permittivity of free space.
\end{definition}

\begin{propriete}[Factors affecting the capacitance of a parallel-plate capacitor]
\begin{itemize}
    \item \textbf{Plate area $A$:} $C \propto A$. Larger plates store more charge for a given voltage.
    \item \textbf{Separation $d$:} $C \propto 1/d$. Bringing plates closer increases the field for the same charge, reducing $V$ and increasing $C$.
    \item \textbf{Dielectric material ($\varepsilon_r$):} $C \propto \varepsilon_r$. Inserting a dielectric increases capacitance by a factor $\varepsilon_r$.
\end{itemize}
These relationships hold qualitatively for other geometries as well (cylindrical, spherical capacitors).
\end{propriete}

\courssub{Action of a dielectric: polarisation}

When a dielectric is placed in an electric field, its molecules become \cle{polarised}: positive and negative charges within each molecule are slightly displaced, creating induced dipoles. In a polar dielectric (e.g., water), permanent dipoles align with the field; in a non-polar dielectric (e.g., polyethylene), dipoles are induced.

The net effect is a layer of induced charge on the dielectric surfaces adjacent to the plates. This induced charge produces an electric field \emph{opposing} the original field, reducing the net field inside the dielectric by a factor $\varepsilon_r$. Since $V = E d$, the potential difference for a given charge $Q$ decreases, so $C = Q/V$ increases.

\begin{savaistu}[Dielectric strength]
Every dielectric has a maximum electric field it can withstand without breaking down (becoming conductive). This is the \cle{dielectric strength}, typically in $\mathrm{V\,m^{-1}}$ or $\mathrm{kV\,mm^{-1}}$. For air it is about $3 \times 10^6\ \mathrm{V\,m^{-1}}$; for mica $\sim 10^8\ \mathrm{V\,m^{-1}}$. Exceeding this limit causes a spark or permanent damage. Capacitors are rated with a maximum working voltage.
\end{savaistu}

Typical values of $\varepsilon_r$ at room temperature:
\begin{center}
\begin{tabularx}{\linewidth}{|l|X|}\hline
\textbf{Material} & \textbf{Relative permittivity $\varepsilon_r$} \\ \hline
Vacuum & 1 (by definition) \\
Air (STP) & 1.0006 $\approx$ 1 \\
Paper (impregnated) & 2--4 \\
Mica & 5--7 \\
Glass & 5--10 \\
Ceramics (Class 1) & 10--100 \\
Ceramics (Class 2, high-K) & 1000--5000 \\
Water (pure, 20°C) & 80 \\ \hline
\end{tabularx}
\end{center}

\courssub{Measurement of capacitance}

The most direct method uses the definition $C = Q/V$. We can measure $Q$ and $V$ simultaneously.

\begin{methode}[Measuring capacitance by constant-current charging]
\begin{enumerate}
    \item Set up the circuit: a DC supply, a switch, a microammeter (or coulombmeter) in series with the capacitor $C$, a variable resistor to maintain constant current, and a voltmeter across $C$.
    \item Close the switch and adjust the variable resistor to give a \emph{constant} charging current $I$ (the current will naturally decrease as the capacitor charges; the resistor must be adjusted continuously, or a constant-current source used).
    \item Start a stopwatch as the switch is closed. Record the time $t$ when the voltmeter reads a chosen voltage $V$.
    \item The charge delivered is $Q = I t$ (since $I$ is constant).
    \item Calculate $C = Q/V = I t / V$.
    \item Repeat for several $(V, t)$ pairs and plot $Q$ vs $V$; the gradient gives $C$.
\end{enumerate}
\end{methode}

\begin{popfigure}
\begin{tikzpicture}[scale=1.0, european]
% Circuit for measuring capacitance
\draw (0,0) to[battery1, l={$E$}] (0,2)
      to[short] (2,2)
      to[switch, l={$S$}] (4,2)
      to[ammeter, l={$I$}] (6,2)
      to[R, l={$R_{\text{var}}$}] (8,2)
      to[C, l={$C$}] (8,0)
      to[short] (0,0);
% Voltmeter across C
\draw (8,2) to[open, v^={$V$}, o-o] (8,0);
\node at (1, 2.5) {DC supply};
\node at (5, 2.5) {$\mu$A};
\node at (7, 2.5) {$R_{\text{var}}$};
\node at (8.5, 1) {$C$};
\end{tikzpicture}
\legende{Experimental circuit for measuring capacitance by constant-current charging. The variable resistor $R_{\text{var}}$ is adjusted to keep the current $I$ constant. The voltmeter measures the voltage $V$ across the capacitor.}
\end{popfigure}

\begin{methode}[Alternative: Reed switch method (common in school laboratories)]
\begin{enumerate}
    \item A reed switch driven at a known frequency $f$ alternately charges the capacitor from a battery of voltage $V$ and discharges it through a microammeter.
    \item The average current measured is $I = Q f = C V f$.
    \item Hence $C = I / (V f)$.
\end{enumerate}
\end{methode}

\begin{popfigure}
\begin{tikzpicture}
\begin{axis}[
    width=10cm, height=6cm,
    axis lines=middle,
    xlabel={$V$ / V}, ylabel={$Q$ / $\mu$C},
    xmin=0, xmax=12,
    ymin=0, ymax=120,
    xtick={0,3,6,9,12},
    ytick={0,30,60,90,120},
    grid=major, grid style={popline},
    title style={font=\small}, title={$Q$--$V$ graph for a capacitor}
]
\addplot[popblue, thick, mark=*] coordinates {
    (0,0)
    (3, 25)
    (6, 50)
    (9, 75)
    (12, 100)
};
\addplot[poporange, dashed, thick] coordinates {
    (0,0) (12, 100)
};
\node[poporange, anchor=south west] at (axis cs:10, 80) {Gradient $= C$};
\end{axis}
\end{tikzpicture}
\legende{$Q$ versus $V$ graph for a capacitor. The graph is a straight line through the origin; its gradient equals the capacitance $C$.}
\end{popfigure}

\courssub{Worked examples}

\begin{exoresolu}[Capacitance of a parallel-plate capacitor]
A parallel-plate capacitor has plates of area $A = 200\ \mathrm{cm^2}$ separated by a distance $d = 0.50\ \mathrm{mm}$. The space between the plates is filled with mica of relative permittivity $\varepsilon_r = 6.0$. Calculate:
\begin{enumerate}
    \item the capacitance $C$,
    \item the charge on each plate when the capacitor is connected to a $12\ \mathrm{V}$ battery,
    \item the electric field strength between the plates.
\end{enumerate}
\tcblower
\textbf{Solution:}
\begin{enumerate}
    \item Convert to SI units:
    $A = 200\ \mathrm{cm^2} = 200 \times 10^{-4}\ \mathrm{m^2} = 0.0200\ \mathrm{m^2}$,
    $d = 0.50\ \mathrm{mm} = 0.50 \times 10^{-3}\ \mathrm{m} = 5.0 \times 10^{-4}\ \mathrm{m}$.
    \[
    C = \frac{\varepsilon_0 \varepsilon_r A}{d}
      = \frac{(8.85 \times 10^{-12}) \times 6.0 \times 0.0200}{5.0 \times 10^{-4}}
      = \frac{1.062 \times 10^{-12}}{5.0 \times 10^{-4}}
      = 2.124 \times 10^{-9}\ \mathrm{F} = 2.12\ \mathrm{nF}.
    \]
    \item $Q = C V = (2.124 \times 10^{-9}) \times 12 = 2.55 \times 10^{-8}\ \mathrm{C} = 25.5\ \mathrm{nC}$.
    \item $E = V/d = 12 / (5.0 \times 10^{-4}) = 2.4 \times 10^4\ \mathrm{V\,m^{-1}}$.
    (Alternatively, $E = \sigma/(\varepsilon_0\varepsilon_r) = Q/(\varepsilon_0\varepsilon_r A)$ gives the same result.)
\end{enumerate}
\end{exoresolu}

\begin{exoresolu}[Effect of dielectric insertion on an isolated capacitor]
A parallel-plate capacitor with air dielectric ($\varepsilon_r = 1$) is charged to a potential difference $V_0$ and then disconnected from the supply. A dielectric slab of relative permittivity $\varepsilon_r = 4$ is inserted completely between the plates. Determine the new:
\begin{enumerate}
    \item capacitance,
    \item potential difference,
    \item electric field,
    \item energy stored.
\end{enumerate}
Express each in terms of the initial values $C_0$, $V_0$, $E_0$, $U_0$.
\tcblower
\textbf{Solution:}
Since the capacitor is isolated, the charge $Q$ remains constant.
\begin{enumerate}
    \item $C = \varepsilon_r C_0 = 4 C_0$.
    \item $V = Q/C = Q/(\varepsilon_r C_0) = V_0/\varepsilon_r = V_0/4$.
    \item $E = V/d = (V_0/4)/d = E_0/4$.
    \item $U = \frac{1}{2} Q V = \frac{1}{2} Q (V_0/4) = U_0/4$.
    \emph{Note:} The energy decreases. The ``missing'' energy is the work done by the electric field in pulling the dielectric into the capacitor (the dielectric is attracted between the plates).
\end{enumerate}
\end{exoresolu}

\courssub{Summary of key points}

\begin{aretenir}[Key points for Section 2]
\begin{itemize}
    \item A capacitor stores charge on two conductors separated by a dielectric. Circuit symbols distinguish fixed, polarised (electrolytic) and variable capacitors.
    \item Capacitance $C = Q/V$; unit: farad (F). Submultiples: $\mu\mathrm{F}$, $\mathrm{nF}$, $\mathrm{pF}$.
    \item For a parallel-plate capacitor: $C = \varepsilon_0 \varepsilon_r A / d$. Capacitance increases with plate area $A$ and $\varepsilon_r$, decreases with separation $d$.
    \item Relative permittivity $\varepsilon_r = C_{\text{dielectric}} / C_{\text{vacuum}}$. Typical values: air $\approx 1$, paper $2$--$4$, mica $5$--$7$, ceramics up to thousands.
    \item A dielectric polarises in an electric field, reducing the net field and increasing capacitance.
    \item Capacitance can be measured by charging with a constant current $I$ for time $t$ ($Q = It$) and measuring $V$, then $C = It/V$. A $Q$--$V$ graph is a straight line through the origin with gradient $C$.
    \item \textbf{Exam traps:} $Q$ is charge on one plate only; always convert submultiples to farads; observe polarity of electrolytic capacitors; distinguish between isolated capacitor (constant $Q$) and capacitor connected to battery (constant $V$) when a dielectric is inserted.
\end{itemize}
\end{aretenir}

\begin{exercice}[Practice questions]
\begin{enumerate}
    \item A $10\ \mu\mathrm{F}$ capacitor is charged to $24\ \mathrm{V}$. Calculate the charge stored on one plate and the energy stored. (Energy formula will be covered in the next section.)
    \item A parallel-plate capacitor with air dielectric has capacitance $C_0$. The plate separation is halved and the space is filled with a dielectric of $\varepsilon_r = 4$. What is the new capacitance in terms of $C_0$?
    \item In a constant-current charging experiment, a current of $50\ \mu\mathrm{A}$ flows for $2.0\ \mathrm{s}$ and the voltage across the capacitor reaches $5.0\ \mathrm{V}$. Determine the capacitance.
    \item Explain why the capacitance of a capacitor increases when a dielectric is inserted, in terms of polarisation and the electric field.
    \item A ceramic capacitor is marked \texttt{473K}. State its capacitance and tolerance.
    \item A parallel-plate capacitor is charged to $100\ \mathrm{V}$ and then disconnected. The plate separation is doubled. What is the new voltage? (Assume air dielectric.)
    \item Describe how you would use a reed switch and a microammeter to measure the capacitance of a capacitor. State the formula used.
\end{enumerate}
\end{exercice}

\begin{corrige}[Exercise 2]
\begin{enumerate}
    \item $Q = CV = 10 \times 10^{-6} \times 24 = 2.4 \times 10^{-4}\ \mathrm{C} = 240\ \mu\mathrm{C}$.
    \item $C = \varepsilon_0 \varepsilon_r A / d$. Halving $d$ doubles $C$; multiplying by $\varepsilon_r = 4$ quadruples it. New $C = 8 C_0$.
    \item $Q = It = 50 \times 10^{-6} \times 2.0 = 1.0 \times 10^{-4}\ \mathrm{C}$. $C = Q/V = 1.0 \times 10^{-4} / 5.0 = 2.0 \times 10^{-5}\ \mathrm{F} = 20\ \mu\mathrm{F}$.
    \item The dielectric polarises, creating induced surface charges that produce an opposing field. The net field $E = E_0/\varepsilon_r$ decreases, so for the same charge $Q$, $V = Ed$ decreases, hence $C = Q/V$ increases by factor $\varepsilon_r$.
    \item Code \texttt{473}: $47 \times 10^3\ \mathrm{pF} = 47\ \mathrm{nF} = 0.047\ \mu\mathrm{F}$. Letter \texttt{K} means $\pm 10\%$ tolerance.
    \item Isolated capacitor $\Rightarrow Q$ constant. $C \propto 1/d$, so doubling $d$ halves $C$. $V = Q/C$, so $V$ doubles: new voltage $= 200\ \mathrm{V}$.
    \item The reed switch charges the capacitor from a battery of voltage $V$ at frequency $f$, then discharges it through the microammeter. Average current $I = Q f = C V f$. Hence $C = I/(V f)$. Measure $I$, $V$, and $f$ to find $C$.
\end{enumerate}
\end{corrige}

\coursec{Combinations of capacitors and energy stored}

In the previous section, we established that a single capacitor is characterised by its capacitance $C$, defined as the ratio of the charge $Q$ stored on one plate to the potential difference $V$ across it: $C = Q/V$. We also saw how the geometry of the plates and the dielectric material (through the relative permittivity $\varepsilon_r$) determine this capacitance. In real circuits — whether in a power supply smoothing circuit in a Douala electronics workshop or the flash unit of a camera — capacitors are rarely used alone. They are combined in networks. Understanding how to replace a network by a single equivalent capacitor, and how energy is distributed and stored, is essential for both the GCE examination and practical design.

\courssub{Capacitors in parallel}

Imagine connecting three capacitors side by side so that all their positive plates are joined to one conductor (at potential $V_+$) and all their negative plates to another conductor (at potential $V_-$). This is a \textbf{parallel combination}.

\begin{popfigure}
\begin{tikzpicture}[scale=0.9, european]
% Parallel combination
\draw[thick] (0,2) -- (5,2) node[above, midway] {$V_+$};
\draw[thick] (0,0) -- (5,0) node[below, midway] {$V_-$};
% Capacitor 1
\draw (1,2) to[C=$C_1$, *-*] (1,0);
\node at (1,1.3) {$+Q_1$};
\node at (1,-0.3) {$-Q_1$};
% Capacitor 2
\draw (2.5,2) to[C=$C_2$, *-*] (2.5,0);
\node at (2.5,1.3) {$+Q_2$};
\node at (2.5,-0.3) {$-Q_2$};
% Capacitor 3
\draw (4,2) to[C=$C_3$, *-*] (4,0);
\node at (4,1.3) {$+Q_3$};
\node at (4,-0.3) {$-Q_3$};
% Battery
\draw (5,0) to[battery1, l={$V$}] (5,2);
% Equivalent
\draw[dashed, thick] (6,2) -- (8,2);
\draw[dashed, thick] (6,0) -- (8,0);
\draw (7,2) to[C=$C_{\text{eq}}$, *-*] (7,0);
\node at (7,1.3) {$+Q = Q_1+Q_2+Q_3$};
\node at (7,-0.3) {$-Q$};
\draw (8,0) to[battery1, l={$V$}] (8,2);
\end{tikzpicture}
\legende{Capacitors in parallel: same potential difference $V$, charges add up.}
\end{popfigure}

\noindent\textbf{Key observation:} The potential difference across each capacitor is exactly the same and equal to the supply voltage $V$. However, the charge stored on each capacitor depends on its individual capacitance: $Q_1 = C_1V$, $Q_2 = C_2V$, $Q_3 = C_3V$.

The total charge $Q$ drawn from the supply is the sum of the charges on each capacitor:
\[ Q = Q_1 + Q_2 + Q_3 = C_1V + C_2V + C_3V = (C_1 + C_2 + C_3)V. \]

By definition, the equivalent capacitance $C_{\text{eq}}$ of the combination is $C_{\text{eq}} = Q/V$. Therefore:

\begin{propriete}[Equivalent Capacitance in Parallel]
For capacitors connected in parallel, the potential difference across each is the same, and the equivalent capacitance is the \textbf{sum} of the individual capacitances:
\[ C_{\text{eq}} = C_1 + C_2 + C_3 + \dots \]
The equivalent capacitance is \textbf{always larger} than the largest individual capacitance. This is physically intuitive: connecting plates in parallel effectively increases the total plate area $A$, and since $C \propto A$, the capacitance increases.
\end{propriete}

\begin{exemplebox}[Parallel Network Calculation]
Three capacitors of capacitances $2\,\mu\mathrm{F}$, $3\,\mu\mathrm{F}$ and $5\,\mu\mathrm{F}$ are connected in parallel across a $12\,\mathrm{V}$ battery. Find:
\begin{enumerate}
\item The equivalent capacitance.
\item The charge on each capacitor.
\item The total charge drawn from the battery.
\end{enumerate}
\textbf{Solution:}
\begin{enumerate}
\item $C_{\text{eq}} = 2 + 3 + 5 = 10\,\mu\mathrm{F}$.
\item Since $V = 12\,\mathrm{V}$ across each:
$Q_1 = C_1V = 2 \times 12 = 24\,\mu\mathrm{C}$,
$Q_2 = 3 \times 12 = 36\,\mu\mathrm{C}$,
$Q_3 = 5 \times 12 = 60\,\mu\mathrm{C}$.
\item $Q_{\text{total}} = C_{\text{eq}}V = 10 \times 12 = 120\,\mu\mathrm{C}$.
(Check: $24 + 36 + 60 = 120\,\mu\mathrm{C}$.)
\end{enumerate}
\end{exemplebox}

\courssub{Capacitors in series}

Now connect the capacitors end-to-end: the negative plate of $C_1$ to the positive plate of $C_2$, the negative of $C_2$ to the positive of $C_3$, and so on. The outer plates connect to the battery terminals.

\begin{popfigure}
\begin{tikzpicture}[scale=0.9, european]
% Series combination
\draw[thick] (0,1) -- (0.5,1);
\draw (0.5,1) to[C=$C_1$, *-*] (1.5,1);
\draw (1.5,1) to[C=$C_2$, *-*] (2.5,1);
\draw (2.5,1) to[C=$C_3$, *-*] (3.5,1);
\draw[thick] (3.5,1) -- (4,1);
% Charges
\node at (1,1.4) {$+Q$}; \node at (1,0.6) {$-Q$};
\node at (2,1.4) {$+Q$}; \node at (2,0.6) {$-Q$};
\node at (3,1.4) {$+Q$}; \node at (3,0.6) {$-Q$};
% Voltages
\node[above] at (1,1.6) {$V_1$};
\node[above] at (2,1.6) {$V_2$};
\node[above] at (3,1.6) {$V_3$};
% Battery
\draw (0,1) to[battery1, l_={$V$}] (0,-1);
\draw[thick] (4,1) -- (4,-1);
\draw[thick] (0,-1) -- (4,-1);
% Equivalent
\draw[dashed, thick] (5,1) -- (6.5,1);
\draw (5.75,1) to[C=$C_{\text{eq}}$, *-*] (5.75,-1);
\node at (5.75,1.4) {$+Q$}; \node at (5.75,0.6) {$-Q$};
\draw (6.5,1) to[battery1, l_={$V$}] (6.5,-1);
\end{tikzpicture}
\legende{Capacitors in series: same charge $Q$ on each, potential differences add up.}
\end{popfigure}

\noindent\textbf{Key observation:} When the battery is connected, electrons flow onto the first plate ($-Q$) and off the last plate ($+Q$). The isolated plates in between (the negative plate of $C_1$ and positive plate of $C_2$, etc.) cannot exchange charge with the outside circuit. By conservation of charge and induction, the magnitude of charge $Q$ on \textbf{every} capacitor in a series chain is identical.

The potential differences add up: $V = V_1 + V_2 + V_3$.
Since $V_1 = Q/C_1$, $V_2 = Q/C_2$, $V_3 = Q/C_3$, we have:
\[ V = \frac{Q}{C_1} + \frac{Q}{C_2} + \frac{Q}{C_3} = Q\left(\frac{1}{C_1} + \frac{1}{C_2} + \frac{1}{C_3}\right). \]

The equivalent capacitance satisfies $V = Q/C_{\text{eq}}$, so:

\begin{propriete}[Equivalent Capacitance in Series]
For capacitors connected in series, the charge $Q$ on each capacitor is the same, and the reciprocal of the equivalent capacitance is the sum of the reciprocals of the individual capacitances:
\[ \frac{1}{C_{\text{eq}}} = \frac{1}{C_1} + \frac{1}{C_2} + \frac{1}{C_3} + \dots \]
The equivalent capacitance is \textbf{always smaller} than the smallest individual capacitance. Physically, the effective plate separation $d$ increases (the gaps add up), and since $C \propto 1/d$, the capacitance decreases.
\end{propriete}

\begin{exemplebox}[Series Network Calculation]
Three capacitors of $2\,\mu\mathrm{F}$, $3\,\mu\mathrm{F}$ and $6\,\mu\mathrm{F}$ are connected in series across a $12\,\mathrm{V}$ battery. Find:
\begin{enumerate}
\item The equivalent capacitance.
\item The charge on each capacitor.
\item The potential difference across each capacitor.
\end{enumerate}
\textbf{Solution:}
\begin{enumerate}
\item $\frac{1}{C_{\text{eq}}} = \frac{1}{2} + \frac{1}{3} + \frac{1}{6} = \frac{3+2+1}{6} = 1 \quad\Rightarrow\quad C_{\text{eq}} = 1\,\mu\mathrm{F}$.
\item $Q = C_{\text{eq}}V = 1 \times 12 = 12\,\mu\mathrm{C}$. This same charge is on $C_1$, $C_2$ and $C_3$.
\item $V_1 = Q/C_1 = 12/2 = 6\,\mathrm{V}$,
$V_2 = 12/3 = 4\,\mathrm{V}$,
$V_3 = 12/6 = 2\,\mathrm{V}$.
Check: $6 + 4 + 2 = 12\,\mathrm{V}$.
\end{enumerate}
\end{exemplebox}

\courssub{Comparison with resistors: a classic trap}

\begin{attention}[Resistors vs. Capacitors -- Rules are Reversed!]
This is one of the most frequent sources of errors in the GCE examination. Memorise this table:
\begin{center}
\begin{tabularx}{\linewidth}{|l|X|X|}\hline
\textbf{Combination} & \textbf{Resistors ($R$)} & \textbf{Capacitors ($C$)} \\ \hline
\textbf{Series} & $R_{\text{eq}} = R_1 + R_2 + \dots$ (Increases) & $\dfrac{1}{C_{\text{eq}}} = \dfrac{1}{C_1} + \dfrac{1}{C_2} + \dots$ (Decreases) \\ \hline
\textbf{Parallel} & $\dfrac{1}{R_{\text{eq}}} = \dfrac{1}{R_1} + \dfrac{1}{R_2} + \dots$ (Decreases) & $C_{\text{eq}} = C_1 + C_2 + \dots$ (Increases) \\ \hline
\end{tabularx}
\end{center}
\textbf{Mnemonic:} \emph{``Capacitors in parallel add like resistors in series (simple sum). Capacitors in series add like resistors in parallel (reciprocal sum).''}
\end{attention}

\courssub{Mixed networks: a step-by-step reduction method}

Real circuits often contain a mixture of series and parallel connections. The strategy is to simplify the network step by step, replacing identified sub-combinations by their equivalent capacitance, until a single $C_{\text{eq}}$ remains.

\begin{methode}[Reducing a Capacitor Network (4 Steps)]
\begin{enumerate}
\item \textbf{Identify} a sub-combination that is \emph{purely} series or purely parallel (no other branches connected to the intermediate nodes).
\item \textbf{Replace} that sub-combination by its equivalent capacitance ($C_{\text{eq}}$). Redraw the circuit diagram -- this is crucial for clarity.
\item \textbf{Repeat} steps 1 and 2 until only one equivalent capacitor remains across the supply terminals.
\item \textbf{Work backwards} (if individual charges/voltages are needed): Start from the total $Q_{\text{total}} = C_{\text{eq}}V_{\text{supply}}$ and distribute charge/voltage back through the steps using the rules: \emph{Series $\rightarrow$ same $Q$; Parallel $\rightarrow$ same $V$}.
\end{enumerate}
\end{methode}

\begin{popfigure}
\begin{tikzpicture}[scale=0.85, european]
% Original Network
\node[align=center] at (2.5, 4.5) {\textbf{Original Network}};
\draw (0,4) to[C, l={$C_1=4\,\mu\mathrm{F}$}, *-*] (0,2);
\draw (0,4) -- (1,4);
\draw (1,4) to[C, l={$C_2=6\,\mu\mathrm{F}$}, *-] (1,3) node[right] {A};
\draw (1,3) to[C, l={$C_3=3\,\mu\mathrm{F}$}, -*] (1,2) node[right] {B};
\draw (1,2) -- (0,2);
\draw (0,2) -- (-0.5,2);
\draw (0,4) -- (-0.5,4);
\draw (-0.5,2) to[battery1, l={$V=12\,\mathrm{V}$}] (-0.5,4);

% Step 1: Parallel C2, C3
\begin{scope}[xshift=7cm]
\node[align=center] at (2.5, 4.5) {\textbf{Step 1: $C_2 \parallel C_3$}};
\draw (0,4) to[C, l={$C_1=4\,\mu\mathrm{F}$}, *-*] (0,2);
\draw (0,4) -- (1,4);
\draw (1,4) to[C, l={$C_{23}=9\,\mu\mathrm{F}$}, *-*] (1,2);
\draw (1,2) -- (0,2);
\draw (0,2) -- (-0.5,2);
\draw (0,4) -- (-0.5,4);
\draw (-0.5,2) to[battery1, l={$V=12\,\mathrm{V}$}] (-0.5,4);
\node at (2.2,3) {$C_{23}=6+3=9\,\mu\mathrm{F}$};
\end{scope}

% Step 2: Series C1, C23
\begin{scope}[xshift=14cm]
\node[align=center] at (2.5, 4.5) {\textbf{Step 2: $C_1$ series $C_{23}$}};
\draw (0,4) to[C, l={$C_{\mathrm{eq}}$}, *-*] (0,2);
\draw (0,2) -- (-0.5,2);
\draw (0,4) -- (-0.5,4);
\draw (-0.5,2) to[battery1, l={$V=12\,\mathrm{V}$}] (-0.5,4);
\node at (2.2,3) {$\dfrac{1}{C_{\mathrm{eq}}} = \dfrac{1}{4} + \dfrac{1}{9} = \dfrac{13}{36}$};
\node at (2.2,2.3) {$C_{\mathrm{eq}} = \dfrac{36}{13} \approx 2.77\,\mu\mathrm{F}$};
\end{scope}
\end{tikzpicture}
\legende{Step-by-step reduction of a mixed network: $C_2$ and $C_3$ in parallel, then that combination in series with $C_1$.}
\end{popfigure}

\begin{exoresolu}[Mixed Network Analysis]
In the network shown above ($C_1=4\,\mu\mathrm{F}$, $C_2=6\,\mu\mathrm{F}$, $C_3=3\,\mu\mathrm{F}$, $V=12\,\mathrm{V}$), determine:
\begin{enumerate}
\item The total charge supplied by the battery.
\item The charge on and potential difference across each capacitor.
\end{enumerate}
\tcblower
\textbf{Solution:}
\begin{enumerate}
\item From the reduction: $C_{\text{eq}} = 36/13\,\mu\mathrm{F}$.
$Q_{\text{total}} = C_{\text{eq}}V = \frac{36}{13} \times 12 = \frac{432}{13} \approx 33.2\,\mu\mathrm{C}$.
\item \textbf{Work backwards:}
\begin{itemize}
\item \textbf{Step 2 (Series $C_1$ and $C_{23}$):} Same charge $Q_{\text{total}}$ flows through both.
$Q_1 = Q_{23} = 33.2\,\mu\mathrm{C}$.
$V_1 = Q_1/C_1 = 33.2/4 = 8.31\,\mathrm{V}$.
$V_{23} = Q_{23}/C_{23} = 33.2/9 = 3.69\,\mathrm{V}$.
(Check: $8.31 + 3.69 \approx 12\,\mathrm{V}$.)
\item \textbf{Step 1 (Parallel $C_2$ and $C_3$):} Same voltage $V_{23} = 3.69\,\mathrm{V}$ across both.
$Q_2 = C_2 V_{23} = 6 \times 3.69 = 22.1\,\mu\mathrm{C}$.
$Q_3 = C_3 V_{23} = 3 \times 3.69 = 11.1\,\mu\mathrm{C}$.
(Check: $Q_2 + Q_3 = 33.2\,\mu\mathrm{C} = Q_{23}$.)
\end{itemize}
\end{enumerate}
\end{exoresolu}

\courssub{Energy stored in a capacitor}

When a capacitor is charged, work is done by the battery (or external agent) to move charge against the growing potential difference. This work is stored as electrostatic potential energy in the electric field between the plates.

Consider the charging process. At an intermediate stage, the capacitor has charge $q$ and potential difference $v = q/C$. To add a small charge $\mathrm{d}q$, the work done is $\mathrm{d}W = v\,\mathrm{d}q = \frac{q}{C}\,\mathrm{d}q$. The total work to charge from $0$ to $Q$ is:
\[ W = \int_0^Q \frac{q}{C}\,\mathrm{d}q = \frac{1}{C}\left[\frac{q^2}{2}\right]_0^Q = \frac{Q^2}{2C}. \]

This is exactly the area under the $V$-$Q$ graph (a straight line through the origin).

\begin{popfigure}
\begin{tikzpicture}
\begin{axis}[
    width=10cm, height=7cm,
    axis lines=middle,
    xlabel={$Q$ ($\mu\mathrm{C}$)}, ylabel={$V$ (V)},
    xmin=0, xmax=12, ymin=0, ymax=13,
    xtick={4,8,12}, ytick={4,8,12},
    domain=0:12,
    samples=2,
    clip=false
]
\addplot[popblue, thick] {x}; % V = Q/C, assume C=1 for simplicity
\addplot[poporange, fill=poporange, fill opacity=0.3, draw=none] coordinates {(0,0) (12,0) (12,12) (0,0)} \closedcycle;
\node[poporange] at (6,4) {$W = \frac{1}{2}QV$};
\draw[dashed] (12,12) -- (12,0) node[below] {$Q$};
\draw[dashed] (12,12) -- (0,12) node[left] {$V$};
\end{axis}
\end{tikzpicture}
\legende{The $V$-$Q$ graph for a capacitor. The shaded triangular area represents the energy stored $W = \frac{1}{2}QV$.}
\end{popfigure}

Using $Q = CV$, we obtain three equivalent forms for the stored energy:

\begin{aretenir}[Energy Stored in a Capacitor]
The electrostatic energy $W$ stored in a charged capacitor can be expressed in three equivalent forms:
\[ W = \frac{1}{2}QV = \frac{1}{2}CV^2 = \frac{Q^2}{2C} \]
\textbf{Units:} Joule ($\mathrm{J}$). \\
\textbf{Where is the energy?} It resides in the \textbf{electric field} between the plates. For a parallel-plate capacitor, the energy density (energy per unit volume) is $u = \frac{1}{2}\varepsilon_0\varepsilon_r E^2$. This concept is fundamental in electromagnetism.
\end{aretenir}

\begin{exemplebox}[Energy Calculations]
A $10\,\mu\mathrm{F}$ capacitor is charged to $100\,\mathrm{V}$.
\begin{enumerate}
\item Calculate the energy stored.
\item The capacitor is then disconnected from the supply and the plate separation is doubled (capacitance halves). What is the new energy? Explain the change.
\end{enumerate}
\textbf{Solution:}
\begin{enumerate}
\item $W = \frac{1}{2}CV^2 = 0.5 \times 10\times10^{-6} \times 100^2 = 0.05\,\mathrm{J} = 50\,\mathrm{mJ}$.
\item Disconnected $\Rightarrow$ Charge $Q$ is constant. $C' = C/2 = 5\,\mu\mathrm{F}$.
$W' = \frac{Q^2}{2C'} = \frac{Q^2}{2(C/2)} = \frac{Q^2}{C} = 2 \times \frac{Q^2}{2C} = 2W = 0.10\,\mathrm{J}$.
The energy \textbf{doubles}. The extra energy comes from the \textbf{work done by the external agent} pulling the plates apart against the attractive electrostatic force.
\end{enumerate}
\end{exemplebox}

\courssub{Energy in combinations and dissipation}

\paragraph{Energy in series and parallel combinations.}
Since energy is a scalar, the total energy stored in a network is simply the sum of the energies stored in each capacitor: $W_{\text{total}} = \sum \frac{1}{2}C_iV_i^2 = \sum \frac{Q_i^2}{2C_i}$. This must equal the energy stored in the equivalent capacitor: $W_{\text{eq}} = \frac{1}{2}C_{\text{eq}}V_{\text{supply}}^2$. This provides a useful check for your calculations.

\paragraph{Energy dissipation during charging.}
A surprising and examinable result: when a capacitor is charged from a battery of constant voltage $V$ through a resistor (or even ideal wires with radiation resistance), the battery supplies energy $W_{\text{batt}} = QV = CV^2$. The capacitor stores only $W_{\text{cap}} = \frac{1}{2}CV^2$. \textbf{Half the energy supplied by the battery is dissipated as heat in the connecting wires/resistor (or radiated)}. This is independent of the resistance value!

\begin{attention}[The ``Missing Half'' Energy]
\textbf{Never} assume the energy supplied by the battery equals the energy stored in the capacitor.
\begin{itemize}
\item Battery supplies: $W_{\text{supplied}} = QV = CV^2$.
\item Capacitor stores: $W_{\text{stored}} = \frac{1}{2}CV^2$.
\item Energy dissipated (heat/radiation): $W_{\text{dissipated}} = \frac{1}{2}CV^2$.
\end{itemize}
This is a favourite GCE multiple-choice and explanation question.
\end{attention}

\paragraph{Connecting charged capacitors (qualitative).}
If a charged capacitor $C_1$ (charge $Q_0$, voltage $V_0$) is connected in parallel to an uncharged capacitor $C_2$, charge redistributes until the voltages equalise. The final voltage is $V_f = Q_0/(C_1+C_2)$. The final energy is $W_f = \frac{1}{2}(C_1+C_2)V_f^2 = \frac{1}{2}\frac{Q_0^2}{C_1+C_2}$. The initial energy was $W_i = \frac{1}{2}\frac{Q_0^2}{C_1}$. Since $C_1+C_2 > C_1$, $W_f < W_i$. The ``lost'' energy is dissipated as heat in the connecting wires (and spark if contact is made abruptly). \textbf{Charge is conserved; energy is not conserved in the capacitor system alone.}

\begin{exoresolu}[Energy Redistribution]
A $4\,\mu\mathrm{F}$ capacitor is charged to $200\,\mathrm{V}$. It is then disconnected from the supply and connected in parallel with an uncharged $2\,\mu\mathrm{F}$ capacitor. Find:
\begin{enumerate}
\item The final potential difference across the combination.
\item The initial and final total energy stored.
\item The energy dissipated.
\end{enumerate}
\tcblower
\textbf{Solution:}
\begin{enumerate}
\item Initial charge: $Q_0 = C_1V_0 = 4\times10^{-6} \times 200 = 800\,\mu\mathrm{C}$.
After connection (parallel): $V_f = \frac{Q_0}{C_1+C_2} = \frac{800}{6} = 133.3\,\mathrm{V}$.
\item $W_i = \frac{1}{2}C_1V_0^2 = 0.5 \times 4\times10^{-6} \times 200^2 = 0.08\,\mathrm{J}$.
$W_f = \frac{1}{2}(C_1+C_2)V_f^2 = 0.5 \times 6\times10^{-6} \times (133.3)^2 \approx 0.0533\,\mathrm{J}$.
\item $W_{\text{dissipated}} = W_i - W_f = 0.08 - 0.0533 = 0.0267\,\mathrm{J}$.
\end{enumerate}
\end{exoresolu}

\begin{savaistu}[Supercapacitors in Yaound\'e Traffic Lights?]
Modern ``supercapacitors'' (or ultracapacitors) can have capacitances of thousands of Farads (e.g., $3000\,\mathrm{F}$) in a small can. They bridge the gap between electrolytic capacitors and batteries. In Cameroon, they are increasingly used in solar-powered street lighting systems (like those along the Yaound\'e-Nsimalen highway) to provide high-current bursts for LED lamps and to smooth the intermittent solar input, lasting far longer than batteries in the harsh tropical heat. Their energy density is still lower than Li-ion batteries, but their power density and cycle life are unmatched.
\end{savaistu}

\begin{exercice}[Practice Problems]
\begin{enumerate}
\item Three capacitors $C_1=2\,\mu\mathrm{F}$, $C_2=4\,\mu\mathrm{F}$, $C_3=6\,\mu\mathrm{F}$ are connected: $C_1$ and $C_2$ in series, and this series combination in parallel with $C_3$. A $12\,\mathrm{V}$ battery is connected across the whole network. Find (a) the equivalent capacitance, (b) the charge on $C_3$, (c) the voltage across $C_1$.
\item A $10\,\mu\mathrm{F}$ capacitor charged to $50\,\mathrm{V}$ is disconnected and then connected in series with an uncharged $20\,\mu\mathrm{F}$ capacitor (the two free outer plates are \textbf{not} connected to each other). Find the final charge on each and the final voltage across each. \textit{Hint: Is there a closed circuit for charge to flow?}
\item Prove that the total energy stored in two capacitors connected in parallel across a battery $V$ is exactly $\frac{1}{2}C_{\text{eq}}V^2$. Show that the energy supplied by the battery is $C_{\text{eq}}V^2$ and deduce the heat dissipated.
\end{enumerate}
\end{exercice}

\begin{corrige}[Exercise 1]
\begin{enumerate}
\item (a) $C_{12} = \frac{C_1C_2}{C_1+C_2} = \frac{8}{6} = \frac{4}{3}\,\mu\mathrm{F}$. $C_{\text{eq}} = C_{12} + C_3 = \frac{4}{3} + 6 = \frac{22}{3} \approx 7.33\,\mu\mathrm{F}$.
(b) $C_3$ is directly across battery: $V_3 = 12\,\mathrm{V}$. $Q_3 = C_3V_3 = 6 \times 12 = 72\,\mu\mathrm{C}$.
(c) Voltage across series branch $V_{12} = 12\,\mathrm{V}$. Charge on series branch $Q_{12} = C_{12}V_{12} = \frac{4}{3} \times 12 = 16\,\mu\mathrm{C}$. This is the charge on $C_1$ and $C_2$. $V_1 = Q_{12}/C_1 = 16/2 = 8\,\mathrm{V}$.
\item Isolated series connection (open circuit): The two inner plates are joined and isolated from the outer plates. The outer plates are not connected to each other. There is \textbf{no closed circuit} for charge to flow. Therefore, \textbf{no redistribution occurs}. The charged capacitor retains its charge $Q_1 = 500\,\mu\mathrm{C}$ and voltage $V_1 = 50\,\mathrm{V}$. The uncharged capacitor remains with $Q_2 = 0$ and $V_2 = 0\,\mathrm{V}$.
\item Parallel combination: $W_{\text{total}} = \frac{1}{2}C_1V^2 + \frac{1}{2}C_2V^2 = \frac{1}{2}(C_1+C_2)V^2 = \frac{1}{2}C_{\text{eq}}V^2$. Battery supplies charge $Q_{\text{total}} = C_{\text{eq}}V$ at constant $V$: $W_{\text{batt}} = Q_{\text{total}}V = C_{\text{eq}}V^2$. Dissipated $= W_{\text{batt}} - W_{\text{stored}} = \frac{1}{2}C_{\text{eq}}V^2$.
\end{enumerate}
\end{corrige}

\coursec{Charging and discharging through a resistor: the time constant}

In the previous section we saw how capacitors combine and how they store energy when connected to a d.c. supply. We implicitly assumed the charging happened instantly. In reality, every circuit has resistance — the internal resistance of the battery, the resistance of connecting wires, or a deliberately added resistor. This resistance limits the current and makes the charging (and discharging) a gradual process governed by a single, crucial parameter: the \cle{time constant}. Understanding this transient behaviour is essential not only for the GCE examination but for designing timing circuits, camera flashes, and life-saving defibrillators.

\courssub{The RC circuit: qualitative description}

Consider a capacitor $C$ connected in series with a resistor $R$ and a d.c. source of e.m.f. $E$ (negligible internal resistance). A two-way switch allows us to connect the capacitor either to the supply (position A: charging) or to a short circuit through the resistor (position B: discharging).

\begin{popfigure}
\begin{tikzpicture}[european, thick, scale=1.1]
\draw (0,0) to[battery1, l_={$E$}, invert] (0,3) node[above] {A};
\draw (0,3) -- (2,3) to[R, l={$R$}] (4,3);
\draw (4,3) to[C, l={$C$}] (4,0);
\draw (4,0) -- (0,0);
\draw (2,3) -- (2,4) node[above] {Charge};
\draw[dashed] (2,3) -- (2,2) node[below] {Discharge};
\node at (1, 1.5) {$S$};
\draw[<->] (4.5, 0.5) -- (4.5, 2.5) node[midway, right] {$V_C$};
\draw[->] (3, 3.3) -- (3, 3.6) node[above] {$I$};
\end{tikzpicture}
\legende{RC series circuit with a two-way switch. Position A (solid): charging from source $E$. Position B (dashed): discharging through $R$.}
\end{popfigure}

\textbf{Charging (Switch at A):} At the instant the switch closes ($t=0$), the capacitor is uncharged ($V_C=0$). The full e.m.f. $E$ appears across the resistor, so the current is maximum: $I_0 = E/R$. As charge accumulates on the plates, $V_C$ rises, opposing the supply. The voltage across the resistor ($V_R = E - V_C$) falls, so the current decreases. Eventually, $V_C \to E$, $V_R \to 0$, and $I \to 0$. The capacitor acts like an open circuit.

\textbf{Discharging (Switch at B):} With the capacitor fully charged to $Q_0 = CE$ ($V_C = E$), flipping the switch to B connects $C$ directly across $R$. The capacitor drives a current through the resistor in the \emph{opposite direction} to the charging current. $V_C$ and $I$ fall exponentially to zero. The capacitor acts like a temporary battery whose voltage decays.

\courssub{The time constant $\tau = RC$}

During these transients, the product $RC$ appears naturally. It is called the \cle{time constant} and denoted $\tau$ (tau).

\begin{definition}[Time constant $\tau$]
For a series RC circuit, the time constant is defined as
\[
\tau = R C
\]
where $R$ is the resistance in ohms ($\Omega$) and $C$ is the capacitance in farads (F).
The unit of $\tau$ is the \cle{second} (s), since $\Omega \cdot \mathrm{F} = \mathrm{V\,A^{-1}} \cdot \mathrm{C\,V^{-1}} = \mathrm{C\,A^{-1}} = \mathrm{s}$.
\end{definition}

\begin{propriete}[Physical meaning of $\tau$]
The time constant $\tau$ has three equivalent practical interpretations:
\begin{enumerate}
    \item \textbf{Charging:} It is the time taken for the charge (or voltage) to rise to $(1 - e^{-1}) \approx 63.2\%$ of its final maximum value.
    \item \textbf{Discharging:} It is the time taken for the charge (or voltage) to fall to $e^{-1} \approx 36.8\%$ of its initial value.
    \item \textbf{Initial slope (tangent method):} It is the time the process would take to complete if the initial rate of change were maintained constant (the tangent at $t=0$ intercepts the final value at $t=\tau$).
\end{enumerate}
\end{propriete}

\courssub{Mathematical equations for charging and discharging}

We derive the equations using Kirchhoff's voltage law (KVL) and the definition of current $I = \dfrac{dQ}{dt}$.

\textbf{1. Charging (Switch at A):}
KVL: $E = V_R + V_C = IR + \dfrac{Q}{C}$.
Since $I = \dfrac{dQ}{dt}$, we have the differential equation:
\[
E = R \frac{dQ}{dt} + \frac{Q}{C} \quad \Rightarrow \quad \frac{dQ}{dt} = \frac{E}{R} - \frac{Q}{RC} = \frac{1}{RC}(EC - Q).
\]
Separating variables and integrating from $t=0$ ($Q=0$) to $t$ ($Q$):
\[
\int_0^Q \frac{dQ'}{EC - Q'} = \int_0^t \frac{dt'}{RC} \quad \Rightarrow \quad \Big[-\ln(EC - Q')\Big]_0^Q = \frac{t}{RC}.
\]
\[
-\ln(EC - Q) + \ln(EC) = \frac{t}{RC} \quad \Rightarrow \quad \ln\left(\frac{EC}{EC - Q}\right) = \frac{t}{RC}.
\]
Solving for $Q$, and noting $Q_0 = CE$ (final charge) and $V_0 = E$ (final voltage):
\[
Q = Q_0 \left(1 - e^{-t/RC}\right), \qquad V = V_0 \left(1 - e^{-t/RC}\right).
\]
Current $I = \dfrac{dQ}{dt} = \dfrac{Q_0}{RC}e^{-t/RC} = \dfrac{V_0}{R}e^{-t/RC} = I_0 e^{-t/RC}$ where $I_0 = V_0/R = E/R$.

\textbf{2. Discharging (Switch at B):}
KVL: $0 = V_R + V_C = IR + \dfrac{Q}{C}$. Current flows \emph{out of} the capacitor, so $I = -\dfrac{dQ}{dt}$ (if $Q$ is charge on the positive plate).
\[
0 = -R \frac{dQ}{dt} + \frac{Q}{C} \quad \Rightarrow \quad \frac{dQ}{dt} = -\frac{Q}{RC}.
\]
Integrating from $t=0$ ($Q=Q_0$) to $t$ ($Q$):
\[
\int_{Q_0}^Q \frac{dQ'}{Q'} = -\int_0^t \frac{dt'}{RC} \quad \Rightarrow \quad \ln\frac{Q}{Q_0} = -\frac{t}{RC}.
\]
Hence:
\[
Q = Q_0 e^{-t/RC}, \qquad V = V_0 e^{-t/RC}.
\]
Current magnitude $I = \left|\frac{dQ}{dt}\right| = \frac{Q_0}{RC}e^{-t/RC} = I_0 e^{-t/RC}$. The \emph{direction} is opposite to the charging current.

\begin{propriete}[The six exponential equations]
Let $\tau = RC$, $Q_0 = CV_0$ (maximum charge), $I_0 = V_0/R$ (maximum current magnitude).
\begin{center}
\begin{tabularx}{\linewidth}{|l|X|X|}\hline
\rowcolor{popblueL}
\textbf{Quantity} & \textbf{Charging (from 0)} & \textbf{Discharging (from $Q_0, V_0$)} \\ \hline
Charge $Q$ & $Q = Q_0 (1 - e^{-t/\tau})$ & $Q = Q_0 e^{-t/\tau}$ \\ \hline
Voltage $V$ & $V = V_0 (1 - e^{-t/\tau})$ & $V = V_0 e^{-t/\tau}$ \\ \hline
Current $I$ & $I = I_0 e^{-t/\tau}$ & $I = -I_0 e^{-t/\tau}$ (opposite direction) \\ \hline
\end{tabularx}
\end{center}
\end{propriete}

\begin{propriete}[Behaviour at limits]
\begin{itemize}
    \item \textbf{At $t=0$ (charging):} $Q=0$, $V=0$, $I=I_0$ (maximum). The uncharged capacitor behaves like a \textbf{short circuit} (wire).
    \item \textbf{As $t \to \infty$ (charging):} $Q \to Q_0$, $V \to V_0$, $I \to 0$. The fully charged capacitor behaves like an \textbf{open circuit} (infinite resistance).
    \item \textbf{At $t=0$ (discharging):} $Q=Q_0$, $V=V_0$, $I=-I_0$ (maximum, opposite direction to charging).
    \item \textbf{As $t \to \infty$ (discharging):} $Q \to 0$, $V \to 0$, $I \to 0$.
\end{itemize}
\end{propriete}

\begin{popfigure}
\begin{tikzpicture}
\begin{axis}[
    width=\linewidth, height=7cm,
    axis lines=middle, xlabel=$t$, ylabel={$V, I$ (normalised)},
    xmin=0, xmax=5, ymin=-0.1, ymax=1.1,
    xtick={1,2,3,4,5}, xticklabels={$\tau$, $2\tau$, $3\tau$, $4\tau$, $5\tau$},
    ytick={0.368, 0.632, 1}, yticklabels={$0.37V_0$, $0.63V_0$, $V_0$},
    legend style={at={(0.98,0.98)}, anchor=north east, font=\small},
    domain=0:5, samples=100,
]
\addplot[popblue, thick] {1 - exp(-x)} node[pos=0.7, right] {$V_{\text{charge}}$};
\addplot[poporange, thick, dashed] {exp(-x)} node[pos=0.3, right] {$I_{\text{charge}}$};
\draw[popmuted, dashed] (axis cs:1,0) -- (axis cs:1,0.632);
\draw[popmuted, dashed] (axis cs:0,0.632) -- (axis cs:1,0.632);
\draw[popmuted, dashed] (axis cs:1,0) -- (axis cs:1,0.368);
\draw[popmuted, dashed] (axis cs:0,0.368) -- (axis cs:1,0.368);
\node[popmuted, font=\small] at (axis cs:1, -0.08) {$\tau$};
\end{axis}
\end{tikzpicture}
\legende{Charging curves: Voltage $V(t)$ rises to $63\%$ of $V_0$ at $t=\tau$; Current $I(t)$ falls to $37\%$ of $I_0$ at $t=\tau$.}
\end{popfigure}

\begin{popfigure}
\begin{tikzpicture}
\begin{axis}[
    width=\linewidth, height=7cm,
    axis lines=middle, xlabel=$t$, ylabel={$Q, V$ (normalised)},
    xmin=0, xmax=5, ymin=-0.1, ymax=1.1,
    xtick={1,2,3,4,5}, xticklabels={$\tau$, $2\tau$, $3\tau$, $4\tau$, $5\tau$},
    ytick={0.368, 1}, yticklabels={$0.37Q_0$, $Q_0$},
    domain=0:5, samples=100,
]
\addplot[popgreen, thick] {exp(-x)} node[pos=0.5, above right] {$Q = Q_0 e^{-t/\tau}$};
\addplot[popred, dashed, thick, domain=0:1.2] {1 - x} node[pos=1.1, right] {Tangent at $t=0$};
\draw[popmuted, dashed] (axis cs:1,0) -- (axis cs:1,0.368);
\draw[popmuted, dashed] (axis cs:0,0.368) -- (axis cs:1,0.368);
\node[popmuted, font=\small] at (axis cs:1, -0.08) {$\tau$};
\node[popmuted, font=\small] at (axis cs:1.1, 0) {$\tau$};
\end{axis}
\end{tikzpicture}
\legende{Discharging curve $Q(t)$. The tangent at $t=0$ (gradient $-Q_0/\tau$) intercepts the time axis at $t=\tau$. At $t=\tau$, $Q = 0.37Q_0$.}
\end{popfigure}

\courssub{Graphical analysis and determination of $\tau$}

In the laboratory (Paper 3 Practical), you will often determine $\tau$ from experimental data. Two graphical methods are standard.

\begin{methode}[Finding $\tau$ from a $V$-$t$ or $Q$-$t$ graph (Direct reading)]
\begin{enumerate}
    \item Plot the experimental points ($V$ vs $t$ for charging or discharging) and draw a smooth curve.
    \item \textbf{Charging:} Find the final voltage $V_0$ (plateau). Calculate $0.63 V_0$. Read the time $t$ corresponding to this voltage on the curve. This $t$ is $\tau$.
    \item \textbf{Discharging:} Find the initial voltage $V_0$. Calculate $0.37 V_0$. Read the time $t$ corresponding to this voltage. This $t$ is $\tau$.
    \item \textbf{Tangent method (works for both):} Draw the tangent to the curve at $t=0$. Find where this tangent intersects the horizontal asymptote ($V_0$ for charging, $0$ for discharging). The time coordinate of this intersection is $\tau$.
\end{enumerate}
\end{methode}

\begin{methode}[Finding $\tau$ from a $\ln V$ vs $t$ graph (Linearisation) — Core Practical Skill]
For discharging: $V = V_0 e^{-t/RC}$. Taking natural logarithms:
\[
\ln V = \ln V_0 - \frac{1}{RC} t.
\]
This is a straight line equation $y = c + mx$ with:
\begin{itemize}
    \item $y = \ln V$ (vertical axis)
    \item $x = t$ (horizontal axis)
    \item Gradient $m = -\dfrac{1}{RC} = -\dfrac{1}{\tau}$
    \item Intercept $c = \ln V_0$
\end{itemize}
\textbf{Procedure:}
\begin{enumerate}
    \item Record $V$ and $t$ during discharge.
    \item Calculate $\ln V$ for each reading.
    \item Plot $\ln V$ (y-axis) against $t$ (x-axis).
    \item Draw the best-fit straight line.
    \item Determine the gradient $m$ (negative).
    \item Calculate $\tau = -\dfrac{1}{m}$ and $C = \dfrac{\tau}{R}$ (if $R$ is known).
\end{enumerate}
\end{methode}

\begin{popfigure}
\begin{tikzpicture}
\begin{axis}[
    width=\linewidth, height=7cm,
    axis lines=middle, xlabel=$t\ (\mathrm{s})$, ylabel=$\ln(V/V_0)$,
    xmin=0, xmax=5, ymin=-2.5, ymax=0.5,
    xtick={1,2,3,4,5}, xticklabels={$\tau$, $2\tau$, $3\tau$, $4\tau$, $5\tau$},
    ytick={-1, 0}, yticklabels={$-1$, $0$},
    domain=0:5, samples=50,
    legend style={at={(0.98,0.98)}, anchor=north east, font=\small},
]
\addplot[poppurple, thick] {-x} node[pos=0.8, below right] {$\ln V = \ln V_0 - t/\tau$};
\draw[popmuted, thick] (axis cs:1,0) -- (axis cs:2,0) node[midway, below] {$\Delta t = \tau$};
\draw[popmuted, thick] (axis cs:2,0) -- (axis cs:2,-1) node[midway, right] {$\Delta \ln V = -1$};
\node[popmuted, font=\small, rotate=90] at (axis cs:2.2, -0.5) {Gradient $= -1/\tau$};
\end{axis}
\end{tikzpicture}
\legende{Linearised discharge graph: $\ln V$ vs $t$. The gradient is $-1/\tau$. A gradient triangle of $\Delta t = \tau$ gives $\Delta \ln V = -1$.}
\end{popfigure}

\courssub{Worked examples}

\begin{exoresolu}[Charging a capacitor]
A $\SI{10}{\mu F}$ capacitor is charged through a $\SI{100}{k\Omega}$ resistor from a $\SI{12}{V}$ d.c. supply.
\begin{enumerate}
    \item Calculate the time constant $\tau$.
    \item Find the charge, voltage, and current at $t = \tau$.
    \item How long does it take for the voltage to reach $\SI{9.0}{V}$?
\end{enumerate}
\tcblower
\textbf{Solution:}
\begin{enumerate}
    \item $\tau = RC = (\SI{100e3}{\Omega}) \times (\SI{10e-6}{F}) = \SI{1.0}{s}$.
    \item $Q_0 = CV_0 = \SI{10e-6}{F} \times \SI{12}{V} = \SI{120}{\mu C}$.
        $I_0 = V_0/R = \SI{12}{V} / \SI{100e3}{\Omega} = \SI{120}{\mu A}$.
        At $t = \tau$:
        $Q = Q_0(1 - e^{-1}) = 120 \times 0.632 = \SI{75.8}{\mu C}$.
        $V = V_0(1 - e^{-1}) = 12 \times 0.632 = \SI{7.58}{V}$.
        $I = I_0 e^{-1} = 120 \times 0.368 = \SI{44.2}{\mu A}$.
    \item Use $V = V_0(1 - e^{-t/\tau})$.
        $9.0 = 12(1 - e^{-t/1.0}) \Rightarrow 1 - e^{-t} = 0.75 \Rightarrow e^{-t} = 0.25$.
        $-t = \ln(0.25) = -\ln 4 \Rightarrow t = \ln 4 \approx \SI{1.39}{s}$.
\end{enumerate}
\end{exoresolu}

\begin{exoresolu}[Analysing discharge data]
A capacitor discharges through a $\SI{22}{k\Omega}$ resistor. The following data is recorded:
\begin{center}
\begin{tabular}{|c|c|c|c|c|c|}\hline
$t\ (\mathrm{s})$ & 0 & 10 & 20 & 30 & 40 \\ \hline
$V\ (\mathrm{V})$ & 10.0 & 6.1 & 3.7 & 2.2 & 1.4 \\ \hline
\end{tabular}
\end{center}
Determine the capacitance $C$ using a graphical method.
\tcblower
\textbf{Solution:}
We linearise: $\ln V = \ln V_0 - t/RC$.
Calculate $\ln V$:
\begin{center}
\begin{tabular}{|c|c|c|c|c|c|}\hline
$t\ (\mathrm{s})$ & 0 & 10 & 20 & 30 & 40 \\ \hline
$\ln V$ & 2.30 & 1.81 & 1.31 & 0.79 & 0.34 \\ \hline
\end{tabular}
\end{center}
Plotting $\ln V$ vs $t$ gives a straight line. Using the first and last points for gradient:
$m = \dfrac{0.34 - 2.30}{40 - 0} = \dfrac{-1.96}{40} = -0.049\ \mathrm{s^{-1}}$.
$\tau = -\dfrac{1}{m} = \dfrac{1}{0.049} \approx \SI{20.4}{s}$.
$C = \dfrac{\tau}{R} = \dfrac{20.4}{22 \times 10^3} = 9.27 \times 10^{-4}\ \mathrm{F} = \SI{927}{\mu F}$.
(Using all points for a best-fit line would give a more accurate result, typically $\approx \SI{930}{\mu F}$).
\end{exoresolu}

\courssub{Applications of RC circuits}

\begin{savaistu}[RC circuits in everyday life]
\begin{itemize}
    \item \textbf{Camera flash units:} A battery charges a capacitor slowly through a high resistance (high $\tau$). Triggering the flash discharges the capacitor through the xenon tube with very low resistance (tiny $\tau \approx \mathrm{ms}$), delivering a huge power pulse.
    \item \textbf{Defibrillators:} A capacitor charged to several kV discharges through the patient's chest in a controlled time ($\tau \approx 5-10\ \mathrm{ms}$) to reset the heart rhythm. The energy stored ($\frac{1}{2}CV^2$) is typically $200-360\ \mathrm{J}$.
    \item \textbf{Smoothing / Filter circuits:} In power supplies, a large capacitor after a rectifier charges quickly (low $R$) and discharges slowly through the load (high $R_{\text{load}}$), reducing ripple voltage. The condition $RC \gg T$ (period of AC) is essential.
    \item \textbf{Car indicator flashers:} An RC timer circuit controls a relay. The capacitor charges until a threshold voltage triggers the relay to switch the lights and discharge the capacitor. The flash rate is set by $\tau$.
\end{itemize}
\end{savaistu}

\courssub{Common mistakes and examination tips}

\begin{attention}[Current direction in discharge]
During discharge, the current flows \textbf{from the positive plate of the capacitor, through the resistor, to the negative plate}. This is \emph{opposite} to the charging current direction. In equations, if charging current is defined as positive, discharge current is $I = -I_0 e^{-t/RC}$. Always state the direction clearly in exams.
\end{attention}

\begin{attention}[Confusing the equations]
\begin{itemize}
    \item Charging: $V$ and $Q$ \textbf{rise} ($1 - e^{-t/\tau}$); $I$ \textbf{falls} ($e^{-t/\tau}$).
    \item Discharging: $V$, $Q$, and $I$ \textbf{all fall} ($e^{-t/\tau}$).
    \item Do not write $V = V_0 e^{-t/\tau}$ for charging voltage!
\end{itemize}
\end{attention}

\begin{attention}[The "5$\tau$" rule]
For all practical purposes, a capacitor is considered fully charged or fully discharged after $5\tau$ (reaches $>99\%$ or $<1\%$). In multiple-choice questions, "fully charged" implies $t \gg \tau$ or $t = 5\tau$.
\end{attention}

\begin{attention}[Logarithm base]
The linearisation $\ln V = \ln V_0 - t/RC$ uses \textbf{natural logarithms} (base $e$). Using $\log_{10}$ gives a gradient of $-1/(2.303 RC)$. Check the question: if it asks for a graph of $\log V$ vs $t$, clarify the base or use $\ln$ (standard in physics).
\end{attention}

\courssub{Consolidation exercises}

\begin{exercice}[Exercise 1: Charging transients]
A $\SI{4.7}{\mu F}$ capacitor is connected in series with a $\SI{1.0}{M\Omega}$ resistor and a $\SI{6.0}{V}$ battery.
\begin{enumerate}
    \item Calculate the time constant.
    \item Determine the maximum charge and maximum current.
    \item Find the voltage across the capacitor and the current in the circuit $\SI{5.0}{s}$ after closing the switch.
    \item Sketch graphs of $V_C(t)$ and $I(t)$ for $0 \le t \le \SI{20}{s}$, marking $\tau$ and the $63\%$ point.
\end{enumerate}
\end{exercice}

\begin{exercice}[Exercise 2: Discharge analysis]
A charged capacitor is discharged through a $\SI{47}{k\Omega}$ resistor. The voltage falls from $\SI{12.0}{V}$ to $\SI{4.4}{V}$ in $\SI{30}{s}$.
\begin{enumerate}
    \item Calculate the time constant $\tau$.
    \item Determine the capacitance $C$.
    \item What percentage of the initial energy remains after $\SI{30}{s}$?
\end{enumerate}
\end{exercice}

\begin{exercice}[Exercise 3: Practical linearisation]
In an experiment to determine the capacitance of an unknown capacitor, a student discharges it through a $\SI{100}{k\Omega}$ resistor and records the voltage $V$ at time $t$.
\begin{center}
\begin{tabular}{|c|c|c|c|c|c|c|}\hline
$t\ (\mathrm{s})$ & 0 & 5 & 10 & 15 & 20 & 25 \\ \hline
$V\ (\mathrm{V})$ & 10.0 & 6.1 & 3.7 & 2.2 & 1.4 & 0.8 \\ \hline
\end{tabular}
\end{center}
\begin{enumerate}
    \item Complete a table of $\ln V$ values.
    \item Plot $\ln V$ against $t$ (on graph paper) and draw the best-fit line.
    \item Determine the gradient and hence the capacitance $C$.
    \item Explain why the $\ln V$ vs $t$ graph is preferred over the $V$ vs $t$ graph for determining $C$ accurately.
\end{enumerate}
\end{exercice}

\courssub{Answers to consolidation exercises}

\begin{corrige}[Exercise 1]
\begin{enumerate}
    \item $\tau = RC = 1.0 \times 10^6 \times 4.7 \times 10^{-6} = \SI{4.7}{s}$.
    \item $Q_0 = CV_0 = 4.7 \times 10^{-6} \times 6.0 = \SI{28.2}{\mu C}$. $I_0 = V_0/R = 6.0 / 1.0 \times 10^6 = \SI{6.0}{\mu A}$.
    \item $t = \SI{5.0}{s} \approx 1.06\tau$.
        $V = 6.0(1 - e^{-5.0/4.7}) = 6.0(1 - e^{-1.064}) = 6.0(1 - 0.345) = \SI{3.93}{V}$.
        $I = 6.0 \times 10^{-6} \times e^{-1.064} = \SI{2.07}{\mu A}$.
    \item Graphs: $V$ rises from 0 to 6V asymptotically, passing through $\SI{3.78}{V}$ at $t=\SI{4.7}{s}$. $I$ falls from $\SI{6.0}{\mu A}$ to 0, passing through $\SI{2.21}{\mu A}$ at $t=\SI{4.7}{s}$.
\end{enumerate}
\end{corrige}

\begin{corrige}[Exercise 2]
\begin{enumerate}
    \item $V = V_0 e^{-t/\tau} \Rightarrow 4.4 = 12.0 e^{-30/\tau}$.
        $e^{-30/\tau} = 4.4/12.0 = 0.3667$.
        $-30/\tau = \ln(0.3667) \approx -1.003$.
        $\tau = 30 / 1.003 \approx \SI{29.9}{s}$.
    \item $C = \tau/R = 29.9 / 47000 = \SI{636}{\mu F}$.
    \item Energy $U \propto V^2$. $U/U_0 = (V/V_0)^2 = (0.3667)^2 = 0.134 = \SI{13.4}{\%}$.
\end{enumerate}
\end{corrige}

\begin{corrige}[Exercise 3]
\begin{enumerate}
    \item $\ln V$: 2.30, 1.81, 1.31, 0.79, 0.34, -0.22.
    \item Graph: Straight line with negative slope.
    \item Gradient $m \approx \dfrac{-0.22 - 2.30}{25 - 0} = \dfrac{-2.52}{25} = -0.1008\ \mathrm{s^{-1}}$.
        $\tau = -1/m = 9.92\ \mathrm{s}$.
        $C = \tau/R = 9.92 / 100000 = \SI{99.2}{\mu F}$.
    \item The $V$-$t$ graph is a steep curve; reading the time for a specific voltage (e.g., $\SI{3.7}{V}$) has large uncertainty. The $\ln V$-$t$ graph is a straight line; the gradient uses \emph{all} data points, averaging out random errors, giving a much more precise determination of $\tau$ and hence $C$.
\end{enumerate}
\end{corrige}

\begin{aretenir}[Key points for revision]
\begin{itemize}
    \item RC circuit: Resistor limits current, causing exponential charge/discharge.
    \item Time constant $\tau = RC$ (seconds). Physical meaning: $63\%$ rise (charge) or $37\%$ fall (discharge); tangent intercept time.
    \item Charging: $Q=Q_0(1-e^{-t/\tau})$, $V=V_0(1-e^{-t/\tau})$, $I=I_0 e^{-t/\tau}$. At $t=0$: short circuit. $t\to\infty$: open circuit.
    \item Discharging: $Q=Q_0 e^{-t/\tau}$, $V=V_0 e^{-t/\tau}$, $I=-I_0 e^{-t/\tau}$ (reverse direction).
    \item Graphical methods: $63\%/37\%$ rule, tangent at $t=0$, $\ln V$ vs $t$ straight line (gradient $=-1/\tau$).
    \item Applications: Flash units (high energy, short $\tau$), smoothing (long $\tau$), timers, defibrillators.
\end{itemize}
\end{aretenir}

\newpage

\coursec{Integration activity}

\courssub{Aims of the activity}

This integration activity brings together everything you have learned in this chapter in a single, realistic engineering task. By the end of it, you should be able to:

\begin{itemize}
\item combine capacitors in series and in parallel and compute the \cle{equivalent capacitance} of a network;
\item calculate the \cle{charge} $Q = CV$ and the \cle{energy} $W = \tfrac{1}{2}CV^{2}$ stored in a capacitor bank connected to a d.c. supply;
\item use the \cle{discharge equation} $V = V_{0}\,e^{-t/RC}$ and the \cle{time constant} $\tau = RC$ to predict how long a capacitor can power a load;
\item choose components from a limited ``catalogue'' to meet a design specification, and justify every choice with a calculation;
\item interpret the limits $t = 0$ and $t \to \infty$ of the discharge curve and sketch the full $V(t)$ graph;
\item communicate a technical design clearly, including its assumptions and limitations, as expected in the problem-solving questions of the GCE Advanced Level paper.
\end{itemize}

\courssub{Situation}

\begin{exemplebox}[The emergency lamp of the Ngoa-Ekelle household]
During the rainy season in Yaoundé, power cuts are frequent in the evening, exactly when students are revising. A young technician, Laetitia, wants to build a small \textbf{emergency lamp delay circuit} for her family's house in Ngoa-Ekelle. The idea is simple: while the mains adaptor is working, a bank of capacitors is charged by a $12\ \mathrm{V}$ d.c. supply; the instant the power fails, a relay disconnects the supply and the capacitor bank discharges through a small white LED lamp, keeping it lit for a short time so that nobody is left in complete darkness.

\textbf{Design specification.} The LED lamp must remain usefully lit for \textbf{at least $30\ \mathrm{s}$} after the power cut. The white LED used behaves as follows (idealised model):
\begin{itemize}
\item it conducts and lights as soon as the voltage across it reaches $V_{\text{on}} = 3.0\ \mathrm{V}$;
\item it stops emitting useful light when the voltage across the capacitor bank falls to the \textbf{cut-off voltage} $V_{c} = 3.0\ \mathrm{V}$;
\item while lit, it draws a current that we model as flowing through an \textbf{equivalent series resistor} $R$ (the resistor Laetitia must choose), so the discharge obeys the usual law
\[
V(t) = V_{0}\,e^{-t/(RC)}, \qquad V_{0} = 12\ \mathrm{V};
\]
\item the LED needs an average power of about $P_{\text{LED}} \approx 30\ \mathrm{mW}$ during the $30\ \mathrm{s}$ to give useful light.
\end{itemize}

\textbf{Laetitia's catalogue} (the only components available in her drawer):

\begin{center}
\begin{tabularx}{\linewidth}{|l|X|X|}
\hline
\rowcolor{popblueL} \textbf{Component} & \textbf{Available values} & \textbf{Remarks} \\
\hline
Electrolytic capacitors & $100\ \mu\mathrm{F}$, $220\ \mu\mathrm{F}$, $470\ \mu\mathrm{F}$, $1000\ \mu\mathrm{F}$ & all rated $16\ \mathrm{V}$ maximum; several of each value available \\
\hline
Resistors & $1\ \mathrm{k}\Omega$, $2.2\ \mathrm{k}\Omega$, $4.7\ \mathrm{k}\Omega$, $10\ \mathrm{k}\Omega$, $22\ \mathrm{k}\Omega$, $47\ \mathrm{k}\Omega$, $100\ \mathrm{k}\Omega$, $220\ \mathrm{k}\Omega$, $470\ \mathrm{k}\Omega$, $1\ \mathrm{M}\Omega$ & $\pm 5\ \%$ tolerance \\
\hline
Supply & $12\ \mathrm{V}$ d.c. adaptor & charges the bank before the cut \\
\hline
\end{tabularx}
\end{center}

\textbf{Safety constraint.} No capacitor may ever be subjected to a voltage greater than its $16\ \mathrm{V}$ rating.

Working in groups, you must produce Laetitia's design: the capacitor arrangement, the resistor, all supporting calculations, a sketch of $V(t)$, and a 5-minute pitch presenting your solution and its limitations.
\end{exemplebox}

\courssub{Tasks}

\begin{enumerate}
\item \textbf{Choosing the arrangement.} A single $1000\ \mu\mathrm{F}$ capacitor is not enough to store the required energy. Propose a parallel combination of at least two catalogue capacitors giving an equivalent capacitance $C$ between $2000\ \mu\mathrm{F}$ and $2500\ \mu\mathrm{F}$. Compute $C$. Explain, in two sentences, why a \emph{parallel} (and not series) arrangement is appropriate here, referring to the voltage rating and to the way capacitances combine.
\item \textbf{Charge and energy stored.} The bank is fully charged by the $12\ \mathrm{V}$ supply. Calculate the charge $Q_{0}$ stored and the energy $W_{0}$ available in the bank.
\item \textbf{Is the energy sufficient?} Estimate the energy the LED needs to shine for $30\ \mathrm{s}$ at $30\ \mathrm{mW}$. Compare with $W_{0}$ and conclude. Why should the \emph{usable} energy actually be taken as $\tfrac{1}{2}C\left(V_{0}^{2} - V_{c}^{2}\right)$ rather than $W_{0}$? Recompute the comparison with this corrected value.
\item \textbf{Choosing the resistor.} The lamp must still be lit at $t = 30\ \mathrm{s}$, i.e. $V(30\ \mathrm{s}) \geq V_{c} = 3.0\ \mathrm{V}$. Using $V(t) = V_{0}e^{-t/(RC)}$, show that the time constant must satisfy $\tau = RC \geq \dfrac{t}{\ln(V_{0}/V_{c})}$, and compute the minimum value of $\tau$. Deduce the \emph{minimum} acceptable value of $R$, then choose the \emph{smallest} catalogue resistor that satisfies the condition. Justify why this choice is a good compromise (a larger $R$ would give longer light but too little current for the LED to be usefully bright).
\item \textbf{Verification.} With your chosen $R$ and $C$, compute $\tau$, the voltage $V(30\ \mathrm{s})$, and the time $t_{c}$ at which the voltage actually reaches $3.0\ \mathrm{V}$. Confirm the specification is met.
\item \textbf{Interpreting the limits.} State the values of $V(t)$ and of the discharge current $I(t) = \dfrac{V_{0}}{R}e^{-t/(RC)}$ at $t = 0$ and as $t \to \infty$. Compute the initial current $I_{0}$ and check it is reasonable for a small LED circuit.
\item \textbf{Sketching the curve.} Sketch $V(t)$ for $0 \leq t \leq 3\tau$, marking $V_{0}$, $V_{c}$, $\tau$ and $t_{c}$. Recall the graphical meaning of $\tau$ (tangent at the origin, value $0.37\,V_{0}$ at $t = \tau$).
\item \textbf{Critical discussion.} Give \emph{two} reasons why the real lamp might shine for a slightly shorter time than your prediction (think of the idealisations in the model). Suggest one improvement to the design.
\end{enumerate}

\courssub{Model answer}

\textbf{1. Choosing the arrangement.}\par
We connect one $1000\ \mu\mathrm{F}$, one $470\ \mu\mathrm{F}$ and one $1000\ \mu\mathrm{F}$ capacitor in parallel. In parallel, the capacitances add:
\[
C = C_{1} + C_{2} + C_{3} = 1000 + 470 + 1000 = 2470\ \mu\mathrm{F} = 2.47\times 10^{-3}\ \mathrm{F}.
\]
This lies in the target band $2000$--$2500\ \mu\mathrm{F}$. A parallel arrangement is appropriate because every capacitor then sits directly across the same $12\ \mathrm{V}$ supply, which is safely below the $16\ \mathrm{V}$ rating, while the capacitances add to give a larger total; in series the equivalent capacitance would be \emph{smaller} than the smallest capacitor, which is the opposite of what we need for energy storage at a fixed voltage.

\textbf{2. Charge and energy stored.}\par
\[
Q_{0} = CV_{0} = 2.47\times 10^{-3} \times 12 = 2.96\times 10^{-2}\ \mathrm{C} \approx 30\ \mathrm{mC}.
\]
\[
W_{0} = \tfrac{1}{2}CV_{0}^{2} = \tfrac{1}{2}\times 2.47\times 10^{-3}\times 12^{2} = 0.178\ \mathrm{J} \approx 0.18\ \mathrm{J}.
\]

\textbf{3. Is the energy sufficient?}\par
Energy needed by the LED:
\[
W_{\text{LED}} = P_{\text{LED}}\,t = 30\times 10^{-3} \times 30 = 0.90\ \mathrm{J}.
\]
This is \emph{larger} than $W_{0} = 0.18\ \mathrm{J}$: with the catalogue capacitors, the bank cannot supply $30\ \mathrm{mW}$ continuously for $30\ \mathrm{s}$. The honest conclusion is that the LED must draw \emph{less} average power (a high-efficiency LED at reduced brightness, or intermittent flashing), or the specification on brightness must be relaxed. Moreover, the usable energy is only the part released between $12\ \mathrm{V}$ and the cut-off $3.0\ \mathrm{V}$:
\[
W_{\text{usable}} = \tfrac{1}{2}C\left(V_{0}^{2} - V_{c}^{2}\right) = \tfrac{1}{2}\times 2.47\times 10^{-3}\times(144 - 9) = 0.167\ \mathrm{J}.
\]
This corresponds to an average available power over $30\ \mathrm{s}$ of $0.167/30 \approx 5.6\ \mathrm{mW}$: enough for a small high-brightness LED glowing dimly, which is acceptable for an emergency ``find-your-candle'' lamp. Laetitia therefore keeps the design but states clearly in her pitch that the lamp gives \emph{reduced} brightness, not full illumination.

\begin{attention}[A trap worth noticing]
Comparing $W_{0}$ with $P_{\text{LED}}\,t$ and concluding ``it works'' without checking the numbers is the classic error. Always do the numerical comparison, and always subtract the energy left in the capacitor at the cut-off voltage: the charge remaining below $V_{c}$ is useless to the load.
\end{attention}

\textbf{4. Choosing the resistor.}\par
We require $V(t) \geq V_{c}$ at $t = 30\ \mathrm{s}$:
\[
V_{0}\,e^{-t/(RC)} \geq V_{c}
\;\Longrightarrow\;
e^{-t/(RC)} \geq \frac{V_{c}}{V_{0}}
\;\Longrightarrow\;
\frac{t}{RC} \leq \ln\!\frac{V_{0}}{V_{c}}
\;\Longrightarrow\;
\tau = RC \geq \frac{t}{\ln(V_{0}/V_{c})}.
\]
Numerically, $\ln(12/3.0) = \ln 4 = 1.386$, so
\[
\tau_{\min} = \frac{30}{1.386} = 21.6\ \mathrm{s}.
\]
Hence the resistance must satisfy
\[
R \geq \frac{\tau_{\min}}{C} = \frac{21.6}{2.47\times 10^{-3}} = 8.74\ \mathrm{k}\Omega.
\]
The catalogue resistors are $1\ \mathrm{k}\Omega$, $2.2\ \mathrm{k}\Omega$, $4.7\ \mathrm{k}\Omega$, $10\ \mathrm{k}\Omega$, $22\ \mathrm{k}\Omega$, \dots\ The smallest value that meets $R \geq 8.74\ \mathrm{k}\Omega$ is $R = 10\ \mathrm{k}\Omega$. A larger $R$ would indeed discharge the bank more slowly (longer light), but the initial LED current $I_{0}=V_{0}/R$ would become too small to produce useful light; $10\ \mathrm{k}\Omega$ gives a safe and adequate initial current of $1.2\ \mathrm{mA}$ (see Task 6), making it the best compromise.

\textbf{5. Verification.}\par
\[
\tau = RC = 10^{4}\times 2.47\times 10^{-3} = 24.7\ \mathrm{s}.
\]
\[
V(30\ \mathrm{s}) = 12\,e^{-30/24.7} = 12\,e^{-1.215} = 12\times 0.297 = 3.6\ \mathrm{V} \geq 3.0\ \mathrm{V}. \checkmark
\]
The actual cut-off time:
\[
t_{c} = \tau\,\ln\frac{V_{0}}{V_{c}} = 24.7\times 1.386 = 34.2\ \mathrm{s} \geq 30\ \mathrm{s}. \checkmark
\]
The specification is met, with a safety margin of about $4\ \mathrm{s}$.

\textbf{6. Interpreting the limits.}\par
At $t = 0$: $V = V_{0} = 12\ \mathrm{V}$ (the capacitor behaves like a full reservoir) and the current is maximal:
\[
I_{0} = \frac{V_{0}}{R} = \frac{12}{10^{4}} = 1.2\ \mathrm{mA},
\]
a safe current for a small LED. As $t \to \infty$: $V \to 0$ and $I \to 0$ exponentially; the capacitor is fully discharged and the lamp is dark. The light fades gradually rather than switching off abruptly — a gentle warning that the energy is running out.

\textbf{7. Sketching the curve.}\par

\begin{popfigure}
\begin{center}
\begin{tikzpicture}
\begin{axis}[
  width=11cm, height=6.5cm,
  axis lines=left,
  xlabel={$t$ (s)}, ylabel={$V$ (V)},
  xmin=0, xmax=75, ymin=0, ymax=13,
  xtick={0,24.7,34.2,49.4,74.1},
  xticklabels={$0$,$\tau$,$t_c$,$2\tau$,$3\tau$},
  ytick={3,4.4,12},
  yticklabels={$V_c$,$0.37V_0$,$V_0$},
  grid=both, grid style={popline, very thin},
]
\addplot[domain=0:75, samples=120, thick, popblue] {12*exp(-x/24.7)};
\addplot[dashed, popmuted] coordinates {(24.7,0) (24.7,4.45)};
\addplot[dashed, popmuted] coordinates {(0,4.45) (24.7,4.45)};
\addplot[dashed, poporange] coordinates {(34.2,0) (34.2,3)};
\addplot[dashed, poporange] coordinates {(0,3) (34.2,3)};
\addplot[domain=0:30, thick, popgreen] {12 - (12/24.7)*x};
\end{axis}
\end{tikzpicture}
\end{center}
\legende{Discharge of the bank through $R = 10\ \mathrm{k}\Omega$: $V(t) = 12e^{-t/24.7}$. The tangent at the origin (green) meets the time axis at $t = \tau$; the voltage crosses $V_c = 3.0\ \mathrm{V}$ at $t_c \approx 34\ \mathrm{s}$.}
\end{popfigure}

\textbf{8. Critical discussion.}\par
Possible reasons the real lamp shines for less than $34\ \mathrm{s}$:
\begin{itemize}
\item \textbf{Component tolerances and leakage:} electrolytic capacitors have tolerances of typically $\pm 20\ \%$ and a small leakage current, so the real stored charge is less than computed; the resistors are only $\pm 5\ \%$.
\item \textbf{Non-ideal LED model:} the LED does not draw current exactly like a fixed resistor; near the cut-off its current falls faster, and internal losses in the relay and wiring dissipate extra energy.
\end{itemize}
\textbf{Improvement:} add a fourth capacitor in parallel (e.g. another $1000\ \mu\mathrm{F}$, giving $C = 3470\ \mu\mathrm{F}$ and $t_{c} \approx 48\ \mathrm{s}$), or use a flashing-LED circuit which draws power only in short pulses and can stretch the useful light well beyond a minute.

\begin{aretenir}[What this activity consolidates]
\begin{itemize}
\item Parallel capacitances add: $C = C_{1}+C_{2}+\dots$; all share the supply voltage, which must stay below each rating.
\item Stored quantities at voltage $V_{0}$: $Q_{0} = CV_{0}$, $W_{0} = \tfrac{1}{2}CV_{0}^{2}$; usable energy down to cut-off: $\tfrac{1}{2}C(V_{0}^{2}-V_{c}^{2})$.
\item Discharge law $V = V_{0}e^{-t/RC}$ with $\tau = RC$; reaching a fraction $V_{c}/V_{0}$ takes $t_{c} = \tau\ln(V_{0}/V_{c})$.
\item A real design ends with a numerical verification \emph{and} an honest discussion of the model's limits.
\end{itemize}
\end{aretenir}

\newpage

\coursec{Methods and worked examples}

\coursec{Capacitors}

\courssub{Skill 1 — Work done and potential energy of a charge in a uniform field}

\begin{definition}[Electric Potential and Potential Difference]
The \cle{electric potential} $V$ at a point is the work done per unit positive charge in bringing a small test charge from infinity to that point. The \cle{potential difference} (p.d.) $V_{AB} = V_A - V_B$ between two points A and B is the work done per unit charge in moving a charge from B to A:
\[ V_{AB} = \frac{W_{B\to A}}{q}. \]
Unit: the volt ($\mathrm{V} = \mathrm{J\,C^{-1}}$). In a uniform field $\vec{E}$, $V_{AB} = E\,d\cos\theta$ where $d$ is the displacement from B to A and $\theta$ the angle between $\vec{E}$ and the displacement.
\end{definition}

\begin{propriete}[Work and Energy in a Uniform Field]
For a charge $q$ displaced from A to B in a uniform electric field $\vec{E}$:
\begin{itemize}
\item Work done by the electric force: $W_{AB} = q\,E\,d\cos\theta = q(V_A - V_B) = -q\,V_{BA}$.
\item Change in electric potential energy: $\Delta E_p = E_{p,B} - E_{p,A} = -W_{AB} = q(V_B - V_A)$.
\item Work done by an external agent moving the charge slowly (no change in kinetic energy): $W_{\text{ext}} = \Delta E_p = q(V_B - V_A)$.
\end{itemize}
\textbf{Examinable derivation:} $W = \int \vec{F}\cdot\mathrm{d}\vec{l} = q\int \vec{E}\cdot\mathrm{d}\vec{l} = q(V_A - V_B)$.
\end{propriete}

\begin{methode}[Work and potential in a uniform electric field]
\begin{enumerate}
\item \textbf{Draw the field.} Sketch the plates or the region, mark the direction of $\vec{E}$ (from $+$ plate to $-$ plate) and the displacement $\vec{d}$ of the charge.
\item \textbf{Compute the field if needed:} $E = \dfrac{V}{d}$ between parallel plates separated by $d$ at p.d. $V$.
\item \textbf{Force on the charge:} $\vec{F}=q\vec{E}$, magnitude $F=|q|E$. Direction: same as $\vec{E}$ if $q>0$, opposite if $q<0$.
\item \textbf{Work of the electric force} for a displacement from A to B:
\[ W_{AB} = q\,E\,d\cos\theta = q(V_A - V_B) = -q\,V_{BA} \]
where $\theta$ is the angle between $\vec{F}$ (or $\vec{E}$ for $q>0$) and the displacement. If the charge moves \emph{along} the force, the field does positive work; if it moves \emph{against} the force, the field does negative work.
\item \textbf{External work} to move the charge slowly (no change of kinetic energy): $W_{\text{ext}} = -W_{AB} = \Delta E_p = q(V_B - V_A)$.
\item \textbf{Check the sign} physically: a positive charge released in a field accelerates along $\vec{E}$ and loses potential energy; a negative charge accelerates opposite to $\vec{E}$.
\end{enumerate}
\end{methode}

\begin{attention}[Sign Conventions — Classic GCE Traps]
\begin{itemize}
\item $W_{\text{field}} = q\Delta V$ where $\Delta V = V_{\text{initial}} - V_{\text{final}}$. Do not confuse with $W_{\text{ext}} = q(V_{\text{final}} - V_{\text{initial}})$.
\item An electron ($q=-e$) moving \emph{from} the negative plate \emph{to} the positive plate moves \emph{against} $\vec{E}$ but \emph{along} $\vec{F}$ (since $\vec{F}=-e\vec{E}$). The field does \emph{positive} work, the electron gains kinetic energy.
\item "Potential energy of the charge" means $E_p = qV$. It decreases when the field does positive work.
\end{itemize}
\end{attention}

\begin{exoresolu}[Moving a charge between charged plates]
Two horizontal parallel plates are separated by $d = \SI{5.0}{cm}$ and maintained at a p.d. of $V = \SI{200}{V}$, the upper plate being positive. An electron ($e = \SI{1.60e-19}{C}$, $m_e = \SI{9.11e-31}{kg}$) is moved slowly from the lower plate to the upper plate. (a) Find the electric field between the plates. (b) Calculate the work done by the electric force during this displacement. (c) Calculate the work that an external agent must supply. (d) If the electron is instead released from rest at the upper plate, find its speed on reaching the lower plate.
\tcblower
\textbf{(a) Field.} For parallel plates,
\[ E = \frac{V}{d} = \frac{200}{5.0\times10^{-2}} = \SI{4.0e3}{V.m^{-1}} \]
directed \emph{downwards} (from the $+$ plate to the $-$ plate).

\textbf{(b) Work of the electric force.} The electron carries $q=-e$, so the electric force on it points \emph{upwards} (opposite to $\vec{E}$). The displacement is also upwards, so the field does \emph{positive} work:
\[ W = |q|\,E\,d = (1.60\times10^{-19})(4.0\times10^{3})(5.0\times10^{-2}) = \SI{3.2e-17}{J}. \]
Equivalently, $W = |q|V = e \times 200\ \mathrm{V} = \SI{3.2e-17}{J}$ — a useful shortcut: the work only depends on the p.d. crossed, not on the path.

\textbf{(c) External work.} Moving the electron \emph{slowly} means $\Delta E_k = 0$, so
\[ W_{\text{ext}} = -W = \SI{-3.2e-17}{J}. \]
The external agent actually \emph{holds back} the electron; the field does the pushing.

\textbf{(d) Released from the upper plate.} The electron is negative, the force on it is \emph{upwards}, so released at the upper plate it accelerates \emph{upwards} — it cannot reach the lower plate. The intended physical situation is release from the \textbf{lower} plate. Then energy conservation gives
\[ \tfrac{1}{2}m_e v^2 = eV \;\Rightarrow\; v = \sqrt{\frac{2eV}{m_e}} = \sqrt{\frac{2(1.60\times10^{-19})(200)}{9.11\times10^{-31}}} = \SI{8.4e6}{m.s^{-1}}. \]
\textbf{Lesson:} always check the sign of the charge before deciding the direction of motion.
\end{exoresolu}

\courssub{Skill 2 — Torque on an electric dipole}

\begin{definition}[Electric Dipole Moment]
An \cle{electric dipole} consists of two equal and opposite charges $+q$ and $-q$ separated by a distance $d$. Its \cle{dipole moment} is the vector
\[ \vec{p} = q\,\vec{d} \]
where $\vec{d}$ points from $-q$ to $+q$. Unit: $\mathrm{C\,m}$.
\end{definition}

\begin{propriete}[Dipole in a Uniform Electric Field]
\begin{itemize}
\item Net force: $\vec{F}_{\text{net}} = \vec{0}$ (forces $+q\vec{E}$ and $-q\vec{E}$ cancel).
\item Torque (couple): $\vec{\tau} = \vec{p} \wedge \vec{E}$, magnitude $\tau = pE\sin\theta$, where $\theta$ is the angle between $\vec{p}$ and $\vec{E}$. The torque tends to align $\vec{p}$ with $\vec{E}$.
\item Potential energy: $U = -\vec{p}\cdot\vec{E} = -pE\cos\theta$. Stable equilibrium at $\theta=0$ ($U_{\min}=-pE$); unstable at $\theta=180^\circ$ ($U_{\max}=+pE$).
\item Work done by an external agent rotating the dipole from $\theta_1$ to $\theta_2$ (slowly):
\[ W_{\text{ext}} = \Delta U = pE(\cos\theta_1 - \cos\theta_2). \]
\end{itemize}
\textbf{Examinable:} Derivation of $\tau = pE\sin\theta$ from the couple of forces $F = qE$ with lever arm $d\sin\theta$.
\end{propriete}

\begin{methode}[Dipole in a uniform field]
\begin{enumerate}
\item Identify the dipole: two charges $+q$ and $-q$ separated by a fixed distance $d$.
\item Compute the \cle{electric dipole moment}: $p = qd$, directed from $-q$ to $+q$.
\item The net force is zero, but the forces form a \textbf{couple}. The torque magnitude is
\[ \tau = pE\sin\theta \]
where $\theta$ is the angle between $\vec{p}$ and $\vec{E}$.
\item Stable equilibrium: $\theta = 0$ ($\vec{p}$ parallel to $\vec{E}$). Unstable equilibrium: $\theta = 180°$.
\item Work done by an external agent rotating the dipole from $\theta_1$ to $\theta_2$:
\[ W = pE(\cos\theta_1 - \cos\theta_2). \]
\end{enumerate}
\end{methode}

\begin{exoresolu}[Torque on a dipole]
An electric dipole consists of charges $\pm \SI{2.0}{\mu C}$ separated by $\SI{4.0}{mm}$. It is placed in a uniform field $E = \SI{5.0e4}{V.m^{-1}}$. (a) Calculate the maximum torque on the dipole. (b) Calculate the work needed to rotate it from stable equilibrium to $\theta = 60°$.
\tcblower
\textbf{(a)} Dipole moment: $p = qd = (2.0\times10^{-6})(4.0\times10^{-3}) = \SI{8.0e-9}{C.m}$.
Maximum torque occurs at $\theta = 90°$:
\[ \tau_{\max} = pE = (8.0\times10^{-9})(5.0\times10^{4}) = \SI{4.0e-4}{N.m}. \]
\textbf{(b)} From $\theta_1 = 0$ to $\theta_2 = 60°$:
\[ W = pE(\cos 0 - \cos 60°) = 4.0\times10^{-4}\,(1 - 0.5) = \SI{2.0e-4}{J}. \]
The work is positive: the field tends to align the dipole, so the external agent must do work against it.
\end{exoresolu}

\courssub{Skill 3 — Capacitance, parallel-plate formula and permittivity}

\begin{definition}[Capacitance]
A \cle{capacitor} is a device that stores electric charge and energy. It consists of two conductors (plates) carrying equal and opposite charges $\pm Q$. The \cle{capacitance} $C$ is defined as
\[ C = \frac{Q}{V} \]
where $Q$ is the magnitude of the charge on \emph{one} plate and $V$ is the potential difference between the plates. Unit: the \cle{farad} ($\mathrm{F} = \mathrm{C\,V^{-1}}$). Practical submultiples: $\mu\mathrm{F}$ ($10^{-6}$), $\mathrm{nF}$ ($10^{-9}$), $\mathrm{pF}$ ($10^{-12}$).
\end{definition}

\begin{popfigure}
\begin{tikzpicture}[european, scale=0.8]
% Circuit symbols for capacitors
\draw (0,0) to[C=$C$] (2,0) node[right] {Fixed capacitor};
\draw (0,-1.5) to[pC, l_={$C$}, *-*] (2,-1.5) node[right] {Polarised (electrolytic)};
\draw (0,-3) to[vC, l_={$C$}] (2,-3) node[right] {Variable capacitor};
\end{tikzpicture}
\legende{Standard circuit symbols for capacitors (IEC / European style).}
\end{popfigure}

\begin{aretenir}[Types of Capacitors (Identification)]
\begin{itemize}
\item \textbf{Ceramic / Film:} Non-polarised, small values (pF–µF), high frequency, low losses.
\item \textbf{Electrolytic (Aluminium/Tantalum):} Polarised (marked +/−), large values (µF–F), used for smoothing/filtering. \textbf{Must be connected with correct polarity.}
\item \textbf{Variable:} Air-spaced or plastic dielectric, capacitance adjusted by rotating plates (old radio tuners).
\end{itemize}
\end{aretenir}

\begin{propriete}[Parallel-Plate Capacitor Formula]
For two parallel plates of area $A$ (overlap area), separation $d$, with a dielectric of relative permittivity $\varepsilon_r$ filling the gap:
\[ C = \frac{\varepsilon_0 \varepsilon_r A}{d} \]
where $\varepsilon_0 = \SI{8.85e-12}{F.m^{-1}}$ (permittivity of free space). $\varepsilon_r \ge 1$ (dimensionless); $\varepsilon_r \approx 1$ for air/vacuum.
\textbf{Factors affecting $C$:} $C \propto A$, $C \propto 1/d$, $C \propto \varepsilon_r$.
\textbf{Examinable derivation:} $E = \sigma/(\varepsilon_0\varepsilon_r) = Q/(\varepsilon_0\varepsilon_r A)$, $V = Ed = Qd/(\varepsilon_0\varepsilon_r A) \Rightarrow C = Q/V = \varepsilon_0\varepsilon_r A/d$.
\end{propriete}

\begin{methode}[Computing and measuring capacitance]
\begin{enumerate}
\item \textbf{Definition:} $C = \dfrac{Q}{V}$, with $Q$ the charge on \emph{one} plate and $V$ the p.d. between the plates.
\item \textbf{Parallel-plate formula:}
\[ C = \frac{\varepsilon_0 \varepsilon_r A}{d} \]
$A$: area of overlap of the plates; $d$: separation; $\varepsilon_0 = \SI{8.85e-12}{F.m^{-1}}$; $\varepsilon_r$: relative permittivity of the dielectric.
\item \textbf{Factors affecting $C$:} $C$ increases with area $A$, decreases with separation $d$, and is multiplied by $\varepsilon_r$ when a dielectric fills the gap.
\item \textbf{Measurement:} Charge the capacitor to a known p.d. $V$, then discharge it through a ballistic galvanometer (or measure $Q$ with a coulombmeter); $C = Q/V$. Alternatively, use the discharge curve through a known resistor (Skill 7).
\item \textbf{Insertion of a dielectric:}
\begin{itemize}
\item At constant $V$ (battery connected): $C \to \varepsilon_r C$, $Q \to \varepsilon_r Q$, $U \to \varepsilon_r U$. Battery supplies energy.
\item At constant $Q$ (battery disconnected): $C \to \varepsilon_r C$, $V \to V/\varepsilon_r$, $U \to U/\varepsilon_r$. Energy decreases (mechanical work done pulling dielectric in).
\end{itemize}
\end{enumerate}
\end{methode}

\begin{attention}[Dielectric Insertion — Constant $V$ vs Constant $Q$]
This is a \textbf{classic GCE trap}. Read the question carefully:
\begin{itemize}
\item "Battery remains connected" $\Rightarrow$ $V$ constant $\Rightarrow$ $Q$ and $U$ increase by factor $\varepsilon_r$.
\item "Battery disconnected before insertion" $\Rightarrow$ $Q$ constant $\Rightarrow$ $V$ and $U$ decrease by factor $\varepsilon_r$.
\end{itemize}
The energy difference comes from the battery (constant $V$) or from mechanical work done by the field pulling the dielectric in (constant $Q$).
\end{attention}

\begin{savaistu}[Relative Permittivity Values]
Typical $\varepsilon_r$: Vacuum = 1 (exact), Air $\approx$ 1.0006, Paper $\approx$ 3.5, Mica $\approx$ 5–7, Glass $\approx$ 5–10, Water $\approx$ 80, Barium Titanate $\approx$ 1200. High-$\varepsilon_r$ ceramics allow tiny SMD capacitors of several $\mu\mathrm{F}$ in your phone.
\end{savaistu}

\begin{exoresolu}[Parallel-plate capacitor with a dielectric]
A parallel-plate capacitor has plates of area $A = \SI{200}{cm^2}$ separated by $d = \SI{1.0}{mm}$ of air ($\varepsilon_r \approx 1$). (a) Calculate its capacitance. (b) It is connected to a $\SI{12}{V}$ battery; find the charge stored. (c) The battery remaining connected, a sheet of mica ($\varepsilon_r = 6.0$) is inserted to fill the gap. Find the new capacitance, charge and energy stored. Comment.
\tcblower
\textbf{(a)} Convert units: $A = 200\times10^{-4}\ \mathrm{m^2} = \SI{2.0e-2}{m^2}$, $d = \SI{1.0e-3}{m}$.
\[ C = \frac{\varepsilon_0 A}{d} = \frac{(\SI{8.85e-12}{F.m^{-1}})(\SI{2.0e-2}{m^2})}{\SI{1.0e-3}{m}} = \SI{1.77e-10}{F} \approx \SI{177}{pF}. \]
\textbf{(b)} $Q = CV = (\SI{1.77e-10}{F})(\SI{12}{V}) = \SI{2.12e-9}{C}$.

\textbf{(c)} With mica, $C' = \varepsilon_r C = 6.0 \times \SI{177}{pF} = \SI{1.06}{nF}$.
The battery stays connected, so $V = \SI{12}{V}$ is unchanged:
\[ Q' = C'V = 6.0\,Q = \SI{1.27e-8}{C}. \]
Energy: $U = \tfrac12 C V^2$ gives initially $U = \tfrac12(\SI{1.77e-10}{F})(\SI{144}{V^2}) = \SI{1.27e-8}{J}$, and finally
\[ U' = \tfrac12 C' V^2 = 6.0\,U = \SI{7.6e-8}{J}. \]
\textbf{Comment:} at fixed $V$, both $Q$ and $U$ are multiplied by $\varepsilon_r$. The extra energy is supplied by the battery (half of the battery's work goes into the capacitor, half into mechanical work pulling the dielectric in). Had the battery been disconnected first, $Q$ would be fixed and the energy would \emph{decrease} by a factor $\varepsilon_r$ — a classic GCE trap.
\end{exoresolu}

\courssub{Skill 4 — Capacitors in series and in parallel}

\begin{propriete}[Combination Rules]
\begin{itemize}
\item \textbf{Parallel:} Same p.d. $V$ across each. Charges add: $Q_{\text{tot}} = Q_1 + Q_2 + \dots$
\[ C_{\text{eq}} = C_1 + C_2 + \dots \quad \text{(adds, like resistors in series)} \]
\item \textbf{Series:} Same charge $Q$ on each (connected end-to-end, no other path). Voltages add: $V = V_1 + V_2 + \dots$ with $V_i = Q/C_i$.
\[ \frac{1}{C_{\text{eq}}} = \frac{1}{C_1} + \frac{1}{C_2} + \dots \quad \text{(like resistors in parallel)} \]
For two capacitors in series: $C_{\text{eq}} = \dfrac{C_1 C_2}{C_1 + C_2}$.
\item \textbf{Voltage division in series:} $V_1 = \dfrac{C_{\text{eq}}}{C_1}V = \dfrac{C_2}{C_1+C_2}V$. The \emph{smallest} capacitor takes the \emph{largest} voltage.
\end{itemize}
\end{propriete}

\begin{methode}[Reducing combinations of capacitors]
\begin{enumerate}
\item \textbf{Identify the topology:} capacitors are in \cle{series} if they carry the same charge (connected end-to-end, no node between them); in \cle{parallel} if they share the same p.d. (connected between the same two nodes).
\item \textbf{Reduce step by step:} start from the branch furthest from the source, replace series/parallel groups by their equivalent, until a single $C_{\text{eq}}$ remains.
\item \textbf{Work back:} find total charge $Q_{\text{tot}} = C_{\text{eq}}V_{\text{supply}}$, then distribute charges and voltages through the network using the rules above.
\item \textbf{Check:} Kirchhoff's voltage law (sum of $V$ around any loop = 0) and charge conservation at nodes.
\end{enumerate}
\end{methode}

\begin{exoresolu}[Mixed network]
Three capacitors $C_1 = \SI{2}{\mu F}$, $C_2 = \SI{4}{\mu F}$ and $C_3 = \SI{12}{\mu F}$ are connected as follows: $C_2$ and $C_3$ in parallel, the combination in series with $C_1$, across a $\SI{24}{V}$ d.c. supply. Find (a) the equivalent capacitance, (b) the charge on each capacitor, (c) the p.d. across each capacitor.
\tcblower
\textbf{(a)} Parallel branch: $C_{23} = C_2 + C_3 = 4 + 12 = \SI{16}{\mu F}$.
Series with $C_1$:
\[ C_{\text{eq}} = \frac{C_1 C_{23}}{C_1 + C_{23}} = \frac{2 \times 16}{2 + 16} = \frac{32}{18} = \SI{1.78}{\mu F}. \]
\textbf{(b)} Charge drawn from the supply: $Q = C_{\text{eq}}V = (\SI{1.78e-6}{F})(\SI{24}{V}) = \SI{4.27e-5}{C} = \SI{42.7}{\mu C}$.
Series elements carry the same charge: $Q_1 = Q_{23} = \SI{42.7}{\mu C}$.
The p.d. across the parallel branch: $V_{23} = Q_{23}/C_{23} = 42.7/16 = \SI{2.67}{V}$.
Hence $Q_2 = C_2 V_{23} = 4 \times 2.67 = \SI{10.7}{\mu C}$ and $Q_3 = C_3 V_{23} = 12 \times 2.67 = \SI{32.0}{\mu C}$.
Check: $Q_2 + Q_3 = 42.7\ \mu\mathrm{C} = Q_1$. \checkmark

\textbf{(c)} $V_1 = Q_1/C_1 = 42.7/2 = \SI{21.3}{V}$; $V_2 = V_3 = \SI{2.67}{V}$.
Check: $V_1 + V_{23} = 21.3 + 2.67 = \SI{24}{V}$. \checkmark
\textbf{Reflex:} always finish with Kirchhoff checks — they cost ten seconds and catch most errors.
\end{exoresolu}

\courssub{Skill 5 — Energy stored in capacitors}

\begin{propriete}[Energy Stored in a Charged Capacitor]
Three equivalent forms — choose the one matching the given data:
\[ U = \tfrac12 QV = \tfrac12 C V^2 = \frac{Q^2}{2C}. \]
\textbf{Examinable derivation:} During charging, the work to add charge $\mathrm{d}q$ against the existing p.d. $v = q/C$ is $\mathrm{d}W = v\,\mathrm{d}q$; integrating from 0 to $Q$:
\[ U = \int_0^Q \frac{q}{C}\,\mathrm{d}q = \frac{Q^2}{2C}. \]
For combinations: compute the energy of \emph{each} capacitor with its own $Q_i$ and $V_i$, then add — or use $U = \tfrac12 C_{\text{eq}} V_{\text{total}}^2$. Both must agree.
\end{propriete}

\begin{attention}[Energy is NOT Conserved in Charge Sharing]
When capacitors are reconnected (e.g. a charged capacitor joined to an uncharged one), \textbf{charge is conserved} but \textbf{energy is not} — the difference is dissipated as heat in the connecting wires (and radiation/spark). Never write "energy conserved" in such problems.
\end{attention}

\begin{exoresolu}[Energy before and after sharing charge]
A $C_1 = \SI{4}{\mu F}$ capacitor is charged to $\SI{100}{V}$ then disconnected from the supply. It is then connected across an uncharged $C_2 = \SI{8}{\mu F}$ capacitor. Find (a) the initial energy, (b) the common p.d. after connection, (c) the final total energy, (d) the energy dissipated.
\tcblower
\textbf{(a)} Initial charge: $Q = C_1 V_1 = (\SI{4e-6}{F})(\SI{100}{V}) = \SI{400}{\mu C}$.
\[ U_i = \tfrac12 C_1 V_1^2 = \tfrac12(\SI{4e-6}{F})(\SI{100}{V})^2 = \SI{2.0e-2}{J} = \SI{20}{mJ}. \]
\textbf{(b)} After connection the capacitors are in parallel; total charge is conserved and the p.d. $V'$ is common:
\[ V' = \frac{Q}{C_1 + C_2} = \frac{\SI{400}{\mu C}}{\SI{12}{\mu F}} = \SI{33.3}{V}. \]
\textbf{(c)} $U_f = \tfrac12 (C_1 + C_2) V'^2 = \tfrac12(\SI{12e-6}{F})(\SI{33.3}{V})^2 = \SI{6.67e-3}{J} = \SI{6.7}{mJ}$.
\textbf{(d)} Energy dissipated:
\[ \Delta U = U_i - U_f = 20 - 6.7 = \SI{13.3}{mJ}, \]
lost as heat in the wires (and a spark if the leads touch!). Notice that two-thirds of the energy vanished even though no charge was lost — energy is \emph{not} conserved in this process.
\end{exoresolu}

\courssub{Skill 6 — Charging a capacitor through a resistor}

\begin{popfigure}
\begin{tikzpicture}[european, scale=1]
\draw (0,0) to[battery1, l_={$E$}] (0,2)
      to[R=$R$] (3,2)
      to[C=$C$] (3,0)
      -- (0,0);
\draw (1.5,2) node[above] {$S$} node[below] {Switch};
\end{tikzpicture}
\legende{Series RC charging circuit. Switch closes at $t=0$.}
\end{popfigure}

\begin{propriete}[RC Charging Equations]
Circuit: battery of e.m.f. $E$, switch, resistor $R$, capacitor $C$ in series. Capacitor initially uncharged ($q(0)=0$).
Kirchhoff's loop law: $E = Ri + \dfrac{q}{C}$, with $i = \dfrac{\mathrm{d}q}{\mathrm{d}t}$.
Differential equation: $\dfrac{\mathrm{d}q}{\mathrm{d}t} + \dfrac{q}{RC} = \dfrac{E}{R}$.
Solution (with $q(0)=0$):
\[ q(t) = Q_0\left(1 - e^{-t/\tau}\right), \qquad Q_0 = CE, \qquad \tau = RC. \]
Hence
\[ u_C(t) = E\left(1 - e^{-t/\tau}\right), \qquad i(t) = \frac{E}{R}e^{-t/\tau}. \]
\textbf{Limits:}
\begin{itemize}
\item At $t=0$: $q=0$, $u_C=0$, $i = E/R$ (capacitor behaves like a wire / short circuit).
\item As $t\to\infty$: $q \to CE$, $u_C \to E$, $i \to 0$ (capacitor behaves like an open switch).
\end{itemize}
\textbf{Time constant $\tau = RC$:} after $t = \tau$, $q = 0.63\,Q_0$; after $t = 5\tau$ the charge exceeds $99\%$ of $Q_0$ — the capacitor is taken as fully charged.
\end{propriete}

\begin{methode}[$RC$ charging: setting up and using the equations]
\begin{enumerate}
\item Draw the circuit, label $E$, $R$, $C$, switch. Define current direction (usually from $+$ of battery).
\item Write Kirchhoff's loop equation: $E = Ri + q/C$.
\item Substitute $i = \mathrm{d}q/\mathrm{d}t$ to get the differential equation.
\item Write the standard solution $q(t) = CE(1-e^{-t/RC})$ and derive $u_C(t)$, $i(t)$.
\item Identify $\tau = RC$. Compute numerical values.
\item Interpret $t=0$ and $t\to\infty$ behaviour (short circuit / open circuit).
\end{enumerate}
\end{methode}

\begin{exoresolu}[Charging a capacitor]
A $\SI{100}{\mu F}$ capacitor, initially uncharged, is charged through a $\SI{47}{k\Omega}$ resistor from a $\SI{12}{V}$ supply. (a) Calculate the time constant. (b) Write the equations for $u_C(t)$ and $i(t)$. (c) Find $u_C$ and $i$ at $t = 0$ and at $t = \tau$. (d) After how long is the capacitor $99\%$ charged?
\tcblower
\textbf{(a)} $\tau = RC = (\SI{47e3}{\Omega})(\SI{100e-6}{F}) = \SI{4.7}{s}$.

\textbf{(b)} $Q_0 = CE = (\SI{100e-6}{F})(\SI{12}{V}) = \SI{1.2e-3}{C}$, and $I_0 = E/R = \SI{12}{V}/\SI{47e3}{\Omega} = \SI{2.55e-4}{A}$.
\[ u_C(t) = \SI{12}{V}\left(1 - e^{-t/4.7}\right), \qquad i(t) = \SI{2.55e-4}{A}\,e^{-t/4.7}. \]

\textbf{(c)} At $t = 0$: $u_C = 0$ and $i = I_0 = \SI{0.255}{mA}$ — all the e.m.f. appears across $R$.
At $t = \tau = \SI{4.7}{s}$: $u_C = 12(1 - e^{-1}) = 12(0.632) = \SI{7.6}{V}$; $i = I_0 e^{-1} = 0.255 \times 0.368 = \SI{0.094}{mA}$.

\textbf{(d)} We need $1 - e^{-t/\tau} = 0.99$, i.e. $e^{-t/\tau} = 0.01$, so
\[ t = \tau \ln 100 = 4.7 \times 4.605 = \SI{21.6}{s} \quad (\approx 4.6\,\tau). \]
\textbf{Interpretation:} charging is fast at first (large current, empty capacitor opposing nothing) and slows down as $u_C$ approaches $E$ and opposes the battery.
\end{exoresolu}

\courssub{Skill 7 — Discharging a capacitor through a resistor}

\begin{popfigure}
\begin{tikzpicture}[european, scale=1]
\draw (0,0) to[C=$C$, l_={$V_0$}] (0,2)
      to[R=$R$] (3,2)
      -- (3,0)
      -- (0,0);
\draw (1.5,2) node[above] {$S$} node[below] {Switch};
\end{tikzpicture}
\legende{Series RC discharging circuit. Capacitor pre-charged to $V_0$, switch closes at $t=0$ (battery removed).}
\end{popfigure}

\begin{propriete}[RC Discharging Equations]
Set-up: the charged capacitor ($Q_0 = CV_0$) is connected across $R$ (supply removed). Loop law: $Ri + \dfrac{q}{C} = 0$.
Differential equation: $\dfrac{\mathrm{d}q}{\mathrm{d}t} = -\dfrac{q}{RC}$, giving
\[ q(t) = Q_0 e^{-t/\tau}, \qquad u_C(t) = V_0 e^{-t/\tau}, \qquad i(t) = -\frac{V_0}{R}e^{-t/\tau}, \quad \tau = RC. \]
\textbf{Limits:} at $t = 0$, $u_C = V_0$ and $|i| = V_0/R$ (maximum); as $t \to \infty$, everything $\to 0$. After $t = \tau$, $u_C = 0.37\,V_0$.
\textbf{Log-linear analysis (measuring $C$):} $\ln u_C = \ln V_0 - \dfrac{t}{RC}$. Plotting $\ln u_C$ against $t$ gives a straight line of gradient $-\dfrac{1}{RC}$; knowing $R$, one obtains $C$.
\end{propriete}

\begin{methode}[$RC$ discharge and measuring $C$]
\begin{enumerate}
\item Write the loop equation for the discharge loop (no battery): $Ri + q/C = 0$.
\item Substitute $i = \mathrm{d}q/\mathrm{d}t$ (current leaves the positive plate, so $i = -\mathrm{d}q/\mathrm{d}t$ if $q$ is charge on positive plate; the equation becomes $\mathrm{d}q/\mathrm{d}t = -q/RC$).
\item Solution: $q = Q_0 e^{-t/RC}$, $u_C = V_0 e^{-t/RC}$, $|i| = (V_0/R)e^{-t/RC}$.
\item For experimental determination of $C$: measure $u_C$ at various $t$, plot $\ln u_C$ vs $t$, gradient $m = -1/RC \Rightarrow C = -1/(mR)$.
\end{enumerate}
\end{methode}

\begin{exoresolu}[Discharge curve and determination of $C$]
A capacitor charged to $V_0 = \SI{9.0}{V}$ is discharged through a resistor $R = \SI{220}{k\Omega}$. A data logger records $u_C$: at $t = \SI{10}{s}$, $u_C = \SI{3.3}{V}$. (a) Show that the decay is consistent with an exponential and estimate $\tau$. (b) Deduce the capacitance. (c) Find the initial discharge current and the time at which it has fallen to $10\%$ of that value.
\tcblower
\textbf{(a)} For an exponential, $u_C = V_0 e^{-t/\tau}$, so
\[ \tau = \frac{t}{\ln(V_0/u_C)} = \frac{10}{\ln(9.0/3.3)} = \frac{10}{\ln 2.727} = \frac{10}{1.003} \approx \SI{10.0}{s}. \]
\textbf{(b)} $C = \dfrac{\tau}{R} = \dfrac{10.0}{220\times10^{3}} = \SI{4.55e-5}{F} = \SI{45.5}{\mu F}$.

\textbf{(c)} Initial current: $I_0 = V_0/R = 9.0/(220\times10^{3}) = \SI{4.09e-5}{A} = \SI{41}{\mu A}$.
Current falls to $10\%$ when $e^{-t/\tau} = 0.10$:
\[ t = \tau \ln 10 = 10.0 \times 2.303 = \SI{23.0}{s}. \]
\textbf{Check by energy:} the energy initially stored, $\tfrac12 C V_0^2 = \tfrac12(\SI{45.5e-6}{F})(\SI{81}{V^2}) = \SI{1.84}{mJ}$, is entirely dissipated in $R$ during the discharge.
\end{exoresolu}

\courssub{Skill 8 — Full analysis of a charge/discharge problem (GCE synthesis)}

\begin{methode}[Tackling a two-way switch problem]
\begin{enumerate}
\item Identify the two switch positions: position 1 = charging (battery in circuit), position 2 = discharging (battery excluded).
\item Write the initial condition for each phase: the charge on $C$ (and hence $u_C$) is \emph{continuous} — it cannot jump instantaneously.
\item Apply the charging or discharging solution with the correct $\tau$ (the resistance \emph{in the loop} during that phase).
\item Sketch $u_C(t)$: exponential rise towards $E$ (charge), exponential fall towards $0$ (discharge).
\item Answer energy questions with $U = \tfrac12 C u_C^2$ evaluated at the instants asked.
\end{enumerate}
\end{methode}

\begin{popfigure}
\begin{tikzpicture}[european, scale=1]
% Two-way switch circuit
\draw (0,-2) to[battery1, l={$E$}] (0,0)
      to[R=$R_1$] (2,0) node[above] {1};
\draw (2,0) to[C=$C$] (2,-2) node[below] {2}
      -- (0,-2);
\draw (2,0) -- (2.5,0) node[above] {Switch};
\draw (2,-2) -- (2.5,-2);
\draw (2.5,0) -- (3,0.5) node[right] {Pos 1 (Charge)};
\draw (2.5,-2) -- (3,-2.5) node[right] {Pos 2 (Discharge)};
\draw (3,0.5) to[R=$R_2$] (3,-2.5);
\end{tikzpicture}
\legende{Two-way switch: Position 1 charges $C$ through $R_1$ from battery $E$; Position 2 discharges $C$ through $R_2$.}
\end{popfigure}

\begin{exoresolu}[Charge then discharge]
A two-way switch connects in position 1 a battery ($E = \SI{6.0}{V}$), a resistor $R_1 = \SI{10}{k\Omega}$ and a capacitor $C = \SI{220}{\mu F}$ in series. The capacitor is initially uncharged. (a) The switch is put in position 1 for $\SI{5.0}{s}$. Find $u_C$ at the end of this phase. (b) The switch is then flipped to position 2, connecting $C$ across $R_2 = \SI{4.7}{k\Omega}$ only. Find $u_C$ a further $\SI{2.0}{s}$ later. (c) Sketch qualitatively $u_C(t)$ for the whole experiment.
\tcblower
\textbf{(a) Charging phase.} $\tau_1 = R_1 C = (\SI{10e3}{\Omega})(\SI{220e-6}{F}) = \SI{2.2}{s}$.
After $t_1 = \SI{5.0}{s}$:
\[ u_C = E\left(1 - e^{-t_1/\tau_1}\right) = 6.0\left(1 - e^{-5.0/2.2}\right) = 6.0(1 - 0.103) = \SI{5.38}{V} \approx \SI{5.4}{V}. \]
\textbf{(b) Discharging phase.} The capacitor starts this phase at $u_C(0) = \SI{5.38}{V}$ (charge is continuous). New time constant:
\[ \tau_2 = R_2 C = (\SI{4.7e3}{\Omega})(\SI{220e-6}{F}) = \SI{1.034}{s}. \]
After a further $t_2 = \SI{2.0}{s}$:
\[ u_C = 5.38\,e^{-t_2/\tau_2} = 5.38\,e^{-2.0/1.034} = 5.38 \times 0.1445 = \SI{0.778}{V} \approx \SI{0.78}{V}. \]
\textbf{(c) Sketch:} from $0$ to $\SI{5}{s}$, a concave-down exponential rise from $0$ towards $\SI{6}{V}$, reaching $\SI{5.4}{V}$; from $\SI{5}{s}$ onwards, a steeper exponential decay ($\tau_2 < \tau_1$) from $\SI{5.4}{V}$ towards $0$, passing through $\SI{0.78}{V}$ at $t = \SI{7}{s}$. The slope is discontinuous at $t = \SI{5}{s}$ because the current changes abruptly, while $u_C$ itself is continuous.
\end{exoresolu}

\begin{savaistu}[Applications in Cameroon]
\begin{itemize}
\item \textbf{Power factor correction:} ENEO installs large capacitor banks (often $\SI{100}{\mu F}$ to $\SI{1}{mF}$, $400\ \mathrm{V}$) near industrial motors to reduce reactive power and lower electricity bills.
\item \textbf{Camera flash / Phone flash:} A battery charges a capacitor slowly ($\tau \sim \mathrm{s}$) through a high resistance; the capacitor then discharges in $\sim \SI{1}{ms}$ through a xenon tube (very low $R$) giving a huge current ($\sim \SI{100}{A}$) for a bright flash.
\item \textbf{Defibrillators:} Capacitors charged to $\sim \SI{5000}{V}$ store $\sim \SI{300}{J}$, discharged through paddles in $\sim \SI{5}{ms}$ to reset heart rhythm.
\end{itemize}
\end{savaistu}

\newpage

\coursec{Exercises}

\courssub{I check my knowledge}

\begin{exercice}[Multiple choice questions]
For each question, choose the single correct answer and justify your choice in one sentence.
\begin{enumerate}
\item The capacitance of a parallel-plate capacitor is doubled when:
\begin{enumerate}
\item the p.d. across the plates is doubled;
\item the charge on the plates is doubled;
\item the plate separation is halved;
\item the plate area is halved.
\end{enumerate}
\item Two capacitors of capacitances $2\ \mu\mathrm{F}$ and $4\ \mu\mathrm{F}$ are connected in series. Their equivalent capacitance is:
\begin{enumerate}
\item $6\ \mu\mathrm{F}$;
\item $3\ \mu\mathrm{F}$;
\item $1.33\ \mu\mathrm{F}$;
\item $0.75\ \mu\mathrm{F}$.
\end{enumerate}
\item The time constant of a circuit containing a resistor $R = 100\ \mathrm{k\Omega}$ and a capacitor $C = 10\ \mu\mathrm{F}$ is:
\begin{enumerate}
\item $0.1\ \mathrm{s}$;
\item $1.0\ \mathrm{s}$;
\item $10\ \mathrm{s}$;
\item $1000\ \mathrm{s}$.
\end{enumerate}
\item When a capacitor discharges through a resistor, after a time $t = \tau$ the p.d. across it has fallen to approximately:
\begin{enumerate}
\item $50\ \%$ of its initial value;
\item $63\ \%$ of its initial value;
\item $37\ \%$ of its initial value;
\item zero.
\end{enumerate}
\end{enumerate}
\end{exercice}

\begin{exercice}[True or false? Justify]
State whether each statement is true or false, and justify your answer with a law, a formula or a short argument.
\begin{enumerate}
\item The work done in moving a charge between two points in an electric field depends on the path followed.
\item The torque on an electric dipole in a uniform field is maximum when the dipole is perpendicular to the field.
\item Inserting a dielectric of relative permittivity $\varepsilon_r$ between the plates of an isolated charged capacitor increases the p.d. across it.
\item In a series combination of capacitors, each capacitor carries the same charge.
\item The energy stored in a capacitor is $\tfrac{1}{2}CV^2$, so doubling the p.d. doubles the stored energy.
\end{enumerate}
\end{exercice}

\begin{exercice}[Definitions and units]
\begin{enumerate}
\item Define the \cle{capacitance} of a capacitor and state its SI unit, expressing the farad in terms of the coulomb and the volt.
\item Define the \cle{relative permittivity} $\varepsilon_r$ of a dielectric. Why does $\varepsilon_r$ have no unit?
\item Define the \cle{time constant} $\tau$ of an $RC$ circuit and give two physical interpretations of its value during the discharge of a capacitor.
\item State what is meant by the \cle{electric dipole moment} $\vec{p}$ of a pair of charges $+q$ and $-q$ separated by a distance $d$, giving its magnitude, direction and SI unit.
\end{enumerate}
\end{exercice}

\begin{exercice}[Quick calculations]
\begin{enumerate}
\item A capacitor stores a charge of $60\ \mu\mathrm{C}$ when the p.d. across it is $12\ \mathrm{V}$. Calculate its capacitance.
\item Calculate the energy stored in a $100\ \mu\mathrm{F}$ capacitor charged to $50\ \mathrm{V}$.
\item A parallel-plate capacitor has plate area $A = 100\ \mathrm{cm^2}$ and separation $d = 2.0\ \mathrm{mm}$, with air between the plates ($\varepsilon_0 = 8.85\times10^{-12}\ \mathrm{F\,m^{-1}}$). Calculate its capacitance.
\item A capacitor of $47\ \mu\mathrm{F}$ discharges through a $220\ \mathrm{k\Omega}$ resistor. Calculate the time constant and the time taken for the p.d. to fall to $5\ \%$ of its initial value (take $\ln 20 \approx 3.0$).
\end{enumerate}
\end{exercice}

\courssub{I practise}

\begin{exercice}[The electric dipole in a uniform field]
An electric dipole consists of charges $+2.0\ \mu\mathrm{C}$ and $-2.0\ \mu\mathrm{C}$ separated by $2.0\ \mathrm{mm}$. It is placed in a uniform electric field of strength $E = 5.0\times10^{4}\ \mathrm{N\,C^{-1}}$, with its axis making an angle of $30^{\circ}$ with the field direction.
\begin{enumerate}
\item Calculate the dipole moment $p$.
\item Calculate the torque $\Gamma$ acting on the dipole.
\item Sketch the dipole in the field, showing the forces on each charge, and state the effect of the torque on the orientation of the dipole.
\item Identify the two equilibrium orientations of the dipole and state, with a reason, which one is stable.
\end{enumerate}
\end{exercice}

\begin{exercice}[Parallel-plate capacitor with a dielectric]
A parallel-plate capacitor has plate area $A = 200\ \mathrm{cm^2}$ and plate separation $d = 1.0\ \mathrm{mm}$. The space between the plates is completely filled with mica of relative permittivity $\varepsilon_r = 6$. Take $\varepsilon_0 = 8.85\times10^{-12}\ \mathrm{F\,m^{-1}}$.
\begin{enumerate}
\item Calculate the capacitance of the capacitor.
\item The capacitor is connected to a $100\ \mathrm{V}$ d.c. supply. Calculate the charge stored and the energy stored.
\item The capacitor is disconnected from the supply (so that its charge is trapped) and the mica sheet is then removed. Calculate the new capacitance and the new p.d. between the plates.
\item Explain why the p.d. changes, and calculate the change in stored energy. Suggest where this extra energy comes from.
\end{enumerate}
\end{exercice}

\begin{exercice}[Series and parallel combinations]
A $2\ \mu\mathrm{F}$ capacitor and a $4\ \mu\mathrm{F}$ capacitor are connected to a $12\ \mathrm{V}$ d.c. supply.
\begin{enumerate}
\item The capacitors are connected \textbf{in series}. Calculate (i) the equivalent capacitance, (ii) the charge on each capacitor, (iii) the p.d. across each capacitor, (iv) the total energy stored.
\item The same two capacitors are now connected \textbf{in parallel} across the same supply. Calculate the same four quantities.
\item Comment on the differences between the two arrangements, stating which combination stores more energy and why.
\end{enumerate}
\end{exercice}

\begin{exercice}[Charging a capacitor through a resistor]
A $100\ \mu\mathrm{F}$ capacitor, initially uncharged, is charged through a $47\ \mathrm{k\Omega}$ resistor from a $9.0\ \mathrm{V}$ battery of negligible internal resistance.
\begin{enumerate}
\item Calculate the time constant of the circuit.
\item Write down the equation giving the charge $q(t)$ on the capacitor during charging, and calculate the charge at $t = \tau$.
\item Calculate the initial charging current (at $t = 0$) and the current at $t = 2\tau$.
\item Sketch, on the same axes, the curves of capacitor p.d. $V_C(t)$ and current $I(t)$ during the charging, indicating the values at $t = 0$, $t = \tau$ and as $t \to \infty$.
\item How long does it take for the capacitor to reach $99\ \%$ of its final charge? (Take $\ln 100 \approx 4.6$.)
\end{enumerate}
\end{exercice}

\courssub{I prepare for the GCE}

\begin{exercice}[GCE-style structured question: combinations and energy]
Two capacitors, $C_1 = 2\ \mu\mathrm{F}$ and $C_2 = 4\ \mu\mathrm{F}$, are available in the laboratory, together with a $12\ \mathrm{V}$ d.c. supply, connecting leads and a two-way switch.
\begin{enumerate}
\item Define the capacitance of a capacitor and state its SI unit. \hfill (2 marks)
\item The capacitors are connected \textbf{in series} across the supply.
\begin{enumerate}
\item Draw the circuit diagram. \hfill (1 mark)
\item Calculate the equivalent capacitance of the combination. \hfill (2 marks)
\item Calculate the charge on each capacitor and the p.d. across each. \hfill (3 marks)
\item Calculate the total energy stored. \hfill (2 marks)
\end{enumerate}
\item The capacitors are now connected \textbf{in parallel} across the same supply. Calculate the equivalent capacitance, the charge on each capacitor and the total energy stored. \hfill (4 marks)
\item Explain, in terms of charge distribution, why the p.d. is the same across each capacitor in parallel but the charge is shared in proportion to capacitance. \hfill (2 marks)
\item A student claims: ``the series combination stores more energy because the charges add up''. Comment critically on this statement, using your results. \hfill (2 marks)
\end{enumerate}
\hfill (Total: 18 marks)
\end{exercice}

\begin{exercice}[GCE-style data analysis: discharge of a capacitor]
In an experiment to measure the capacitance of a capacitor, a charged capacitor is discharged through a resistor of resistance $R = 470\ \mathrm{k\Omega}$. The p.d. $V$ across the capacitor is measured at $10\ \mathrm{s}$ intervals, giving the following results:
\begin{center}
\begin{tabular}{|l|c|c|c|c|c|c|c|}
\hline
\rowcolor{popblueL} $t$ / s & 0 & 10 & 20 & 30 & 40 & 50 & 60 \\
\hline $V$ / V & 12.0 & 7.3 & 4.4 & 2.7 & 1.6 & 1.0 & 0.6 \\
\hline
\end{tabular}
\end{center}
\begin{enumerate}
\item State the equation describing the discharge of a capacitor through a resistor, defining each symbol. \hfill (2 marks)
\item Show that a graph of $\ln V$ against $t$ should be a straight line, and state what its gradient and intercept represent. \hfill (3 marks)
\item Complete a table of values of $\ln(V/\mathrm{V})$ and plot the graph of $\ln V$ against $t$ on graph paper, choosing suitable scales. \hfill (4 marks)
\item Use your graph to determine the time constant $\tau$ of the circuit, and hence calculate the capacitance $C$ of the capacitor. \hfill (4 marks)
\item The resistor is stated by the manufacturer to have a tolerance of $\pm 5\ \%$, and the gradient is determined to within $\pm 4\ \%$. Estimate the percentage uncertainty in your value of $C$, and state the result in the form $C = (\ldots \pm \ldots)\ \mu\mathrm{F}$. \hfill (3 marks)
\item Suggest one practical precaution that would improve the accuracy of the experiment. \hfill (1 mark)
\end{enumerate}
\hfill (Total: 17 marks)
\end{exercice}

\begin{exercice}[GCE-style problem: flash lamp and energy transfer]
\textbf{Part A.} A camera flash unit uses a capacitor of capacitance $2000\ \mu\mathrm{F}$ charged to $300\ \mathrm{V}$. The flash is produced by discharging the capacitor through the flash lamp, and the useful flash duration is approximately equal to one time constant, $\tau = 1.0\ \mathrm{ms}$.
\begin{enumerate}
\item Calculate the charge and the energy stored in the fully charged capacitor. \hfill (3 marks)
\item Estimate the resistance of the flash lamp during the discharge. \hfill (2 marks)
\item Calculate the initial discharge current, and explain why the brightness of the flash decreases with time. \hfill (3 marks)
\item Calculate the fraction of the initial energy still stored in the capacitor after one time constant. Comment on the efficiency of the flash. \hfill (3 marks)
\end{enumerate}
\textbf{Part B (extension).} In a separate experiment, a $10\ \mu\mathrm{F}$ capacitor is charged to $24\ \mathrm{V}$, then disconnected from the supply and connected across an identical, initially uncharged $10\ \mu\mathrm{F}$ capacitor.
\begin{enumerate}
\setcounter{enumi}{4}
\item Explain why charge is conserved in this process, and calculate the final p.d. across the combination. \hfill (3 marks)
\item Calculate the total energy stored before and after the connection, and show that energy appears to be ``lost''. \hfill (3 marks)
\item Explain carefully where the ``missing'' energy has gone, naming at least two mechanisms. \hfill (2 marks)
\end{enumerate}
\hfill (Total: 19 marks)
\end{exercice}

\newpage

\coursec{Solutions to the exercises}

\coursec{Corrigés des exercices — Chapitre PHY05 : Condensateurs}

\begin{corrige}[Exercice 1 — Questions à choix multiples]
\begin{enumerate}
\item \textbf{Bonne réponse : (c) la distance entre les armatures est divisée par deux.}\\
Justification : pour un condensateur plan $C = \dfrac{\varepsilon_0 \varepsilon_r A}{d}$ ; la capacité ne dépend que de la géométrie et du diélectrique, pas de $Q$ ni de $V$. Diviser $d$ par deux double $C$ ; diviser $A$ par deux divise $C$ par deux.
\item \textbf{Bonne réponse : (c) $1,33\ \mu\mathrm{F}$.}\\
Justification : en série, $\dfrac{1}{C_s} = \dfrac{1}{2} + \dfrac{1}{4} = \dfrac{3}{4}\ \mu\mathrm{F}^{-1}$, donc $C_s = \dfrac{4}{3} = 1,33\ \mu\mathrm{F}$. (Une association série est toujours plus petite que la plus petite capacité.)
\item \textbf{Bonne réponse : (b) $1,0\ \mathrm{s}$.}\\
Justification : $\tau = RC = (100\times10^{3})(10\times10^{-6}) = 1,0\ \mathrm{s}$.
\item \textbf{Bonne réponse : (c) $37\ \%$ de sa valeur initiale.}\\
Justification : $V = V_0 e^{-t/\tau}$ ; à $t = \tau$, $V = V_0 e^{-1} \approx 0,37\,V_0$.
\end{enumerate}
\end{corrige}

\begin{corrige}[Exercice 2 — Vrai ou faux ? Justifier]
\begin{enumerate}
\item \textbf{Faux.} Le champ électrostatique est conservatif : le travail $W = q(V_B - V_A)$ ne dépend que des potentiels des extrémités, pas du chemin. C'est pourquoi le potentiel électrique est bien défini.
\item \textbf{Vrai.} Le couple est $\Gamma = pE\sin\theta$ ; $\sin\theta$ est maximal ($=1$) pour $\theta = 90^{\circ}$, c'est-à-dire quand l'axe du dipôle est perpendiculaire à $\vec{E}$.
\item \textbf{Faux.} Pour un condensateur \emph{isolé}, $Q$ est constant. L'insertion du diélectrique multiplie $C$ par $\varepsilon_r$, donc $V = Q/C$ est \emph{divisé} par $\varepsilon_r$ : la tension \emph{diminue}. (Elle resterait constante seulement si le condensateur restait branché sur la source.)
\item \textbf{Vrai.} En série, la charge déplacée dans le circuit est la même partout : les armatures intérieures portent $+Q$ et $-Q$ par induction, donc chaque condensateur porte la même charge $Q$.
\item \textbf{Faux.} $U = \tfrac{1}{2}CV^2$ varie comme le \emph{carré} de la tension : doubler $V$ multiplie l'énergie par $4$, pas par $2$.
\end{enumerate}
\end{corrige}

\begin{corrige}[Exercice 3 — Définitions et unités]
\begin{enumerate}
\item La \cle{capacité} $C$ d'un condensateur est le rapport de la charge $Q$ stockée sur une armature à la tension $V$ entre les armatures : $C = \dfrac{Q}{V}$. Unité SI : le \textbf{farad} (F), avec $1\ \mathrm{F} = 1\ \mathrm{C\,V^{-1}}$ (un coulomb par volt).
\item La \cle{permittivité relative} $\varepsilon_r$ d'un diélectrique est le rapport de la capacité d'un condensateur rempli de ce diélectrique à la capacité du même condensateur sous vide : $\varepsilon_r = \dfrac{C}{C_0} = \dfrac{\varepsilon}{\varepsilon_0}$. Étant un rapport de deux grandeurs de même nature (deux capacités ou deux permittivités), les unités se simplifient : $\varepsilon_r$ est un nombre pur sans unité.
\item La \cle{constante de temps} d'un circuit $RC$ est $\tau = RC$. Interprétations en décharge : (i) c'est le temps après lequel la tension (ou la charge) est tombée à $e^{-1} \approx 37\ \%$ de sa valeur initiale ; (ii) c'est le temps que mettrait le condensateur pour se décharger complètement s'il continuait à se décharger à son \emph{taux initial} (la tangente en $t=0$ coupe l'axe des temps en $t = \tau$).
\item Le \cle{moment dipolaire électrique} $\vec{p}$ de deux charges $+q$ et $-q$ séparées d'une distance $d$ est le vecteur de norme $p = qd$, dirigé le long de l'axe du dipôle \emph{de la charge négative vers la charge positive}. Unité SI : coulomb-mètre ($\mathrm{C\,m}$).
\end{enumerate}
\end{corrige}

\begin{corrige}[Exercice 4 — Calculs rapides]
\begin{enumerate}
\item $C = \dfrac{Q}{V} = \dfrac{60\times10^{-6}}{12} = 5,0\times10^{-6}\ \mathrm{F} = 5,0\ \mu\mathrm{F}$.
\item $U = \tfrac{1}{2}CV^2 = \tfrac{1}{2}(100\times10^{-6})(50)^2 = 0,125\ \mathrm{J} \approx 0,13\ \mathrm{J}$.
\item Conversion des unités : $A = 100\ \mathrm{cm^2} = 1,0\times10^{-2}\ \mathrm{m^2}$, $d = 2,0\times10^{-3}\ \mathrm{m}$.
\[ C = \frac{\varepsilon_0 A}{d} = \frac{(8,85\times10^{-12})(1,0\times10^{-2})}{2,0\times10^{-3}} = 4,4\times10^{-11}\ \mathrm{F} = 44\ \mathrm{pF}. \]
\item $\tau = RC = (220\times10^{3})(47\times10^{-6}) = 10,3\ \mathrm{s} \approx 10\ \mathrm{s}$.\\
Pour $V = 0,05\,V_0$ : $0,05 = e^{-t/\tau} \Rightarrow t = \tau\ln 20 \approx 10,3 \times 3,0 \approx 31\ \mathrm{s}$.
\end{enumerate}
\end{corrige}

\begin{corrige}[Exercice 5 — Dipôle électrique dans un champ uniforme]
\begin{enumerate}
\item $p = qd = (2,0\times10^{-6})(2,0\times10^{-3}) = 4,0\times10^{-9}\ \mathrm{C\,m}$.
\item $\Gamma = pE\sin\theta = (4,0\times10^{-9})(5,0\times10^{4})\sin 30^{\circ} = (2,0\times10^{-4})(0,5) = 1,0\times10^{-4}\ \mathrm{N\,m}$.
\item Schéma : la force sur $+q$ est $\vec{F}_+ = q\vec{E}$ (parallèle à $\vec{E}$) ; la force sur $-q$ est $\vec{F}_- = -q\vec{E}$ (antiparallèle). Les deux forces sont égales et opposées mais non confondues : elles forment un couple.
\begin{popfigure}
\begin{center}
\begin{tikzpicture}[scale=1.1]
\foreach \x in {0,1.2,2.4,3.6} {\draw[-latex,popblue] (\x,-1.1) -- (\x,1.3);}
\node[popblue] at (4.1,1.1) {$\vec{E}$};
\draw[thick,popdark] (1.3,-0.55) -- (2.3,0.55);
\fill[poporange] (2.3,0.55) circle (0.12); \node[above] at (2.3,0.6) {$+q$};
\fill[poppurple] (1.3,-0.55) circle (0.12); \node[below] at (1.3,-0.65) {$-q$};
\draw[-latex,thick,poporange] (2.3,0.55) -- (3.1,0.55) node[right] {$\vec{F}_+$};
\draw[-latex,thick,poppurple] (1.3,-0.55) -- (0.5,-0.55) node[left] {$\vec{F}_-$};
\draw[popmuted] (1.8,0) arc (0:30:0.5); \node at (2.35,0.12) {$\theta$};
\end{tikzpicture}
\end{center}
\legende{Dipôle dans un champ uniforme : le couple $(\vec{F}_+,\vec{F}_-)$ produit un torque $\Gamma = pE\sin\theta$.}
\end{popfigure}
La force résultante est nulle (pas de translation), mais le couple fait tourner le dipôle pour \emph{aligner} $\vec{p}$ avec $\vec{E}$ ($\theta \to 0$).
\item L'équilibre impose $\Gamma = 0$, soit $\theta = 0$ ou $\theta = 180^{\circ}$. $\theta = 0$ ($\vec{p}$ parallèle à $\vec{E}$) est l'équilibre \textbf{stable} : un petit déplacement produit un couple de rappel qui ramène le dipôle. $\theta = 180^{\circ}$ est instable : tout petit déplacement produit un couple qui l'éloigne (c'est aussi un maximum de l'énergie potentielle $U = -pE\cos\theta$).
\end{enumerate}
\end{corrige}

\begin{corrige}[Exercice 6 — Condensateur plan avec diélectrique]
Données : $A = 200\ \mathrm{cm^2} = 2,0\times10^{-2}\ \mathrm{m^2}$ ; $d = 1,0\times10^{-3}\ \mathrm{m}$ ; $\varepsilon_r = 6$.
\begin{enumerate}
\item $C = \dfrac{\varepsilon_0 \varepsilon_r A}{d} = \dfrac{(8,85\times10^{-12})(6)(2,0\times10^{-2})}{1,0\times10^{-3}} = 1,06\times10^{-9}\ \mathrm{F} \approx 1,1\ \mathrm{nF}$.
\item $Q = CV = (1,06\times10^{-9})(100) = 1,06\times10^{-7}\ \mathrm{C} \approx 0,11\ \mu\mathrm{C}$.\\
$U = \tfrac{1}{2}CV^2 = \tfrac{1}{2}(1,06\times10^{-9})(100)^2 = 5,3\times10^{-6}\ \mathrm{J} = 5,3\ \mu\mathrm{J}$.
\item Alimentation débranchée $\Rightarrow$ $Q$ piégé : $Q = 1,06\times10^{-7}\ \mathrm{C}$ constant. Retrait de la mica divise la capacité par $\varepsilon_r$ :
\[ C' = \frac{C}{\varepsilon_r} = \frac{1,06\times10^{-9}}{6} = 1,77\times10^{-10}\ \mathrm{F} \approx 0,18\ \mathrm{nF}. \]
Nouvelle tension : $V' = \dfrac{Q}{C'} = \varepsilon_r V = 6 \times 100 = 600\ \mathrm{V}$.
\item La tension augmente car la même charge est stockée sur une capacité plus faible ($V = Q/C$). Nouvelle énergie :
\[ U' = \tfrac{1}{2}C'V'^2 = \tfrac{1}{2}(1,77\times10^{-10})(600)^2 = 3,2\times10^{-5}\ \mathrm{J} = 32\ \mu\mathrm{J}. \]
$\Delta U = U' - U \approx 27\ \mu\mathrm{J}$ (l'énergie est multipliée par $\varepsilon_r = 6$). L'énergie supplémentaire provient du \textbf{travail mécanique fourni par l'opérateur} qui retire la plaque de mica : le diélectrique polarisé est attiré dans le champ, il faut donc fournir un travail contre cette attraction.
\end{enumerate}
\end{corrige}

\begin{corrige}[Exercice 7 — Associations série et parallèle]
\begin{enumerate}
\item \textbf{Série} ($C_1 = 2\ \mu\mathrm{F}$, $C_2 = 4\ \mu\mathrm{F}$, $V = 12\ \mathrm{V}$) :
\begin{enumerate}
\item $\dfrac{1}{C_s} = \dfrac{1}{2} + \dfrac{1}{4} = \dfrac{3}{4} \Rightarrow C_s = \dfrac{4}{3} = 1,33\ \mu\mathrm{F}$.
\item Même charge sur chacun : $Q = C_s V = (1,33\times10^{-6})(12) = 1,6\times10^{-5}\ \mathrm{C} = 16\ \mu\mathrm{C}$.
\item $V_1 = \dfrac{Q}{C_1} = \dfrac{16}{2} = 8,0\ \mathrm{V}$ ; \quad $V_2 = \dfrac{Q}{C_2} = \dfrac{16}{4} = 4,0\ \mathrm{V}$. (Vérification : $V_1 + V_2 = 12\ \mathrm{V}$.)
\item $U_s = \tfrac{1}{2}C_s V^2 = \tfrac{1}{2}(1,33\times10^{-6})(144) = 9,6\times10^{-5}\ \mathrm{J} = 96\ \mu\mathrm{J}$.
\end{enumerate}
\item \textbf{Parallèle} : même tension $V = 12\ \mathrm{V}$ aux bornes de chacun.
\begin{enumerate}
\item $C_p = C_1 + C_2 = 6,0\ \mu\mathrm{F}$.
\item $Q_1 = C_1 V = 24\ \mu\mathrm{C}$ ; \quad $Q_2 = C_2 V = 48\ \mu\mathrm{C}$ (total $72\ \mu\mathrm{C}$).
\item $V_1 = V_2 = 12\ \mathrm{V}$.
\item $U_p = \tfrac{1}{2}C_p V^2 = \tfrac{1}{2}(6,0\times10^{-6})(144) = 4,3\times10^{-4}\ \mathrm{J} = 432\ \mu\mathrm{J}$.
\end{enumerate}
\item L'association parallèle stocke plus d'énergie ($432\ \mu\mathrm{J} \gg 96\ \mu\mathrm{J}$) : les deux condensateurs sont chargés à la tension d'alimentation, donc la charge totale prélevée ($72\ \mu\mathrm{C}$) est bien plus grande qu'en série ($16\ \mu\mathrm{C}$), et $U = \tfrac{1}{2}QV$ est d'autant plus grande. En série la capacité équivalente est réduite en dessous de la plus petite valeur, ce qui limite la charge et l'énergie stockées.
\end{enumerate}
\end{corrige}

\begin{corrige}[Exercice 8 — Charge d'un condensateur à travers une résistance]
\begin{enumerate}
\item $\tau = RC = (47\times10^{3})(100\times10^{-6}) = 4,7\ \mathrm{s}$.
\item $q(t) = Q_0\left(1 - e^{-t/\tau}\right)$ avec $Q_0 = CE = (100\times10^{-6})(9,0) = 9,0\times10^{-4}\ \mathrm{C}$.\\
À $t = \tau$ : $q = Q_0(1 - e^{-1}) = 0,63\,Q_0 = 5,7\times10^{-4}\ \mathrm{C}$.
\item Courant initial : $I_0 = \dfrac{E}{R} = \dfrac{9,0}{47\times10^{3}} = 1,9\times10^{-4}\ \mathrm{A} \approx 0,19\ \mathrm{mA}$.\\
$I(t) = I_0 e^{-t/\tau}$ ; à $t = 2\tau$ : $I = I_0 e^{-2} = 0,135\,I_0 \approx 2,6\times10^{-5}\ \mathrm{A} \approx 26\ \mu\mathrm{A}$.
\item Schéma :
\begin{popfigure}
\begin{center}
\begin{tikzpicture}[scale=1.0]
\draw[-latex] (0,0) -- (6.4,0) node[below] {$t$};
\draw[-latex] (0,0) -- (0,3.0) node[left] {};
\draw[domain=0:6,smooth,thick,popblue] plot (\x,{2.4*(1-exp(-\x/1.5))});
\draw[domain=0:6,smooth,thick,poporange] plot (\x,{2.4*exp(-\x/1.5)});
\draw[dashed,popmuted] (0,2.4) -- (6.2,2.4) node[right,popblue] {$E$};
\draw[dashed,popmuted] (1.5,0) node[below] {$\tau$} -- (1.5,1.51);
\draw[dashed,popmuted] (0,1.51) node[left] {$0,63E$} -- (1.5,1.51);
\draw[dashed,popmuted] (0,0.88) node[left] {$0,37I_0$} -- (1.5,0.88);
\node[left] at (0,2.4) {$I_0$};
\node[popblue] at (4.6,2.1) {$V_C$};
\node[poporange] at (4.6,0.45) {$I$};
\end{tikzpicture}
\end{center}
\legende{Charge : $V_C$ monte vers $E$ ; $I$ décroît de $I_0 = E/R$ vers zéro. À $t=\tau$ : $V_C = 0,63E$, $I = 0,37I_0$.}
\end{popfigure}
À $t=0$ : $V_C = 0$, $I = I_0$ ; quand $t\to\infty$ : $V_C \to E = 9,0\ \mathrm{V}$, $I \to 0$.
\item $q = 0,99\,Q_0 \Rightarrow 1 - e^{-t/\tau} = 0,99 \Rightarrow e^{-t/\tau} = 0,01 \Rightarrow t = \tau\ln 100 = 4,7 \times 4,6 \approx 22\ \mathrm{s}$.
\end{enumerate}
\end{corrige}

\begin{corrige}[Exercice 9 — Question de type GCE : associations et énergie]
\begin{enumerate}
\item \textbf{(2 points)} Capacité $C = Q/V$ : charge stockée par unité de tension entre les armatures \textbf{(1)}. Unité SI : le farad (F), $1\ \mathrm{F} = 1\ \mathrm{C\,V^{-1}}$ \textbf{(1)}.
\item \textbf{Série :}
\begin{enumerate}
\item \textbf{(1 point)} Schéma : source, interrupteur et les deux condensateurs en série dans une seule maille.
\begin{popfigure}
\begin{center}
\begin{circuitikz}[european]
\draw (0,0) to[battery1,l={$12\ \mathrm{V}$}] (0,2.4) to[nos] (1.6,2.4) to[C=$C_1$] (3.2,2.4) to[C=$C_2$] (4.8,2.4) -- (4.8,0) -- (0,0);
\end{circuitikz}
\end{center}
\legende{$C_1$ et $C_2$ en série aux bornes de la source $12\ \mathrm{V}$.}
\end{popfigure}
\item \textbf{(2 points)} $\dfrac{1}{C_s} = \dfrac{1}{C_1} + \dfrac{1}{C_2} = \dfrac{1}{2} + \dfrac{1}{4} = \dfrac{3}{4}\ \mu\mathrm{F}^{-1}$ \textbf{(1)} ; $C_s = 1,33\ \mu\mathrm{F}$ \textbf{(1)}.
\item \textbf{(3 points)} $Q = C_s V = 1,33\times12 = 16\ \mu\mathrm{C}$ sur \emph{chaque} condensateur \textbf{(1)} ; $V_1 = Q/C_1 = 8,0\ \mathrm{V}$ \textbf{(1)} ; $V_2 = Q/C_2 = 4,0\ \mathrm{V}$ \textbf{(1)}.
\item \textbf{(2 points)} $U = \tfrac{1}{2}C_s V^2 = \tfrac{1}{2}(1,33\times10^{-6})(12)^2$ \textbf{(1)} $= 9,6\times10^{-5}\ \mathrm{J}$ \textbf{(1)}.
\end{enumerate}
\item \textbf{Parallèle (4 points) :} $C_p = 2 + 4 = 6,0\ \mu\mathrm{F}$ \textbf{(1)} ; $Q_1 = C_1 V = 24\ \mu\mathrm{C}$ \textbf{(1)} ; $Q_2 = C_2 V = 48\ \mu\mathrm{C}$ \textbf{(1)} ; $U = \tfrac{1}{2}C_p V^2 = 4,3\times10^{-4}\ \mathrm{J}$ \textbf{(1)}.
\item \textbf{(2 points)} En parallèle, les deux paires d'armatures sont reliées directement aux deux mêmes bornes de la source, donc la tension aux bornes de chacune doit être égale à la tension de la source \textbf{(1)}. Comme $Q = CV$ avec le même $V$, la charge prélevée est proportionnelle à la capacité : le condensateur le plus grand stocke proportionnellement plus de charge \textbf{(1)}.
\item \textbf{(2 points)} L'affirmation est fausse. En série les charges ne s'additionnent \emph{pas} : chaque condensateur porte la \emph{même} charge $16\ \mu\mathrm{C}$, et l'énergie totale n'est que $96\ \mu\mathrm{J}$, contre $432\ \mu\mathrm{J}$ en parallèle \textbf{(1)}. L'association parallèle stocke plus d'énergie car les deux condensateurs sont à la pleine tension $12\ \mathrm{V}$ et la charge totale stockée ($72\ \mu\mathrm{C}$) est bien plus grande \textbf{(1)}.\\
\emph{Attentes du correcteur : utilisation correcte de la formule des inverses avec étapes montrées ; énoncé clair que la charge série est commune ; réponses numériques avec unités ; commentaire critique appuyé sur les propres résultats du candidat.}
\end{enumerate}
\end{corrige}

\begin{corrige}[Exercice 10 — Question de type GCE : analyse de données (décharge)]
\begin{enumerate}
\item \textbf{(2 points)} $V = V_0 e^{-t/\tau}$ avec $\tau = RC$ \textbf{(1)} ; $V$ = tension à l'instant $t$, $V_0$ = tension initiale, $R$ = résistance, $C$ = capacité \textbf{(1)}.
\item \textbf{(3 points)} En prenant les logarithmes népériens : $\ln V = \ln V_0 - \dfrac{t}{RC}$ \textbf{(1)}, qui est de la forme $y = c + mx$ : une droite de pente $-\dfrac{1}{RC} = -\dfrac{1}{\tau}$ \textbf{(1)} et d'ordonnée à l'origine $\ln V_0$ \textbf{(1)}.
\item \textbf{(4 points)} Tableau de valeurs :
\begin{center}
\begin{tabular}{|l|c|c|c|c|c|c|c|}
\hline
\rowcolor{popblueL} $t$ / s & 0 & 10 & 20 & 30 & 40 & 50 & 60 \\
\hline $V$ / V & 12,0 & 7,3 & 4,4 & 2,7 & 1,6 & 1,0 & 0,6 \\
\hline $\ln(V/\mathrm{V})$ & 2,48 & 1,99 & 1,48 & 0,99 & 0,47 & 0,00 & $-0,51$ \\
\hline
\end{tabular}
\end{center}
Échelles : par ex. $2\ \mathrm{cm} : 10\ \mathrm{s}$ sur l'axe $t$ et $2\ \mathrm{cm} : 0,5$ unité sur l'axe $\ln V$. Tracer les points et tracer la droite de meilleur ajustement \textbf{(1 point pour le tableau, 1 pour échelles/axes, 1 pour points, 1 pour droite)}.
\begin{popfigure}
\begin{center}
\begin{tikzpicture}[scale=1.0]
\draw[-latex] (0,0) -- (6.6,0) node[below] {$t$/s};
\draw[-latex] (0,0) -- (0,3.4) node[left] {$\ln(V/\mathrm{V})$};
\foreach \t/\y in {0/2.48,10/1.99,20/1.48,30/0.99,40/0.47,50/0.00,60/-0.51}
  \fill[poporange] (\t/10,{(\y+0.6)*0.9}) circle (0.06);
\draw[thick,popblue] (0,3.1) -- (6.3,0.05);
\foreach \y/\l in {0.99/1,1.89/2,2.79/3} \node[left] at (0,\y) {$\l$};
\foreach \x/\l in {1/10,2/20,3/30,4/40,5/50,6/60} \node[below] at (\x,0) {$\l$};
\end{tikzpicture}
\end{center}
\legende{$\ln V$ en fonction de $t$ : droite de pente $-1/\tau$.}
\end{popfigure}
\item \textbf{(4 points)} Pente lue sur la droite : par ex. en utilisant $(0, 2,48)$ et $(60\ \mathrm{s}, -0,51)$ :
\[ \text{pente} = \frac{-0,51 - 2,48}{60 - 0} = -5,0\times10^{-2}\ \mathrm{s^{-1}} \textbf{(2 points pour grand triangle et valeur)}. \]
$\tau = -\dfrac{1}{\text{pente}} = \dfrac{1}{0,050} = 20\ \mathrm{s}$ \textbf{(1)}.\\
$C = \dfrac{\tau}{R} = \dfrac{20}{470\times10^{3}} = 4,3\times10^{-5}\ \mathrm{F} = 43\ \mu\mathrm{F}$ \textbf{(1)}.
\item \textbf{(3 points)} Comme $C = \tau/R$ et $\tau$ vient de la pente, les incertitudes relatives s'additionnent :
\[ \frac{\Delta C}{C} = 4\% + 5\% = 9\% \textbf{(1)}; \quad \Delta C = 0,09 \times 43 \approx 4\ \mu\mathrm{F} \textbf{(1)}; \]
\[ C = (43 \pm 4)\ \mu\mathrm{F} \textbf{(1)}. \]
\item \textbf{(1 point)} N'importe laquelle : utiliser un voltmètre à haute impédance (numérique) pour ne pas décharger le condensateur ; prendre des mesures à intervalles plus courts / plus de points ; répéter et faire la moyenne ; utiliser un enregistreur de données (data logger) pour acquérir $V$ et $t$ simultanément ; s'assurer que le condensateur est complètement chargé avant de commencer.\\
\emph{Attentes du correcteur : linéarisation explicitée ; valeurs des logs données à 2--3 chiffres significatifs ; pente lue sur un grand triangle de la droite de meilleur ajustement (pas sur les points bruts) ; incertitudes additionnées pour un quotient ; réponse finale avec précision cohérente.}
\end{enumerate}
\end{corrige}

\begin{corrige}[Exercice 11 — Problème de type GCE : lampe flash et transfert d'énergie]
\textbf{Partie A.}
\begin{enumerate}
\item \textbf{(3 points)} $Q = CV = (2000\times10^{-6})(300) = 0,60\ \mathrm{C}$ \textbf{(1)} ; $U = \tfrac{1}{2}CV^2 = \tfrac{1}{2}(2000\times10^{-6})(300)^2$ \textbf{(1)} $= 90\ \mathrm{J}$ \textbf{(1)}.
\item \textbf{(2 points)} $\tau = RC \Rightarrow R = \dfrac{\tau}{C} = \dfrac{1,0\times10^{-3}}{2000\times10^{-6}}$ \textbf{(1)} $= 0,50\ \Omega$ \textbf{(1)}.
\item \textbf{(3 points)} $I_0 = \dfrac{V}{R} = \dfrac{300}{0,50} = 600\ \mathrm{A}$ \textbf{(1)}. Pendant la décharge $V = V_0 e^{-t/\tau}$ et $I = I_0 e^{-t/\tau}$ décroissent exponentiellement \textbf{(1)} ; la puissance dissipée dans la lampe ($P = VI$) chute donc rapidement, ce qui fait diminuer la luminosité avec le temps \textbf{(1)}.
\item \textbf{(3 points)} Après $t = \tau$ : $V = V_0 e^{-1}$, et comme $U \propto V^2$,
\[ \frac{U}{U_0} = e^{-2} = 0,135 \approx 14\% \textbf{(2)}. \]
Environ $86\%$ de l'énergie stockée est restituée pendant le flash utile : le flash est assez efficace pour extraire l'énergie en une constante de temps \textbf{(1)}.
\end{enumerate}
\textbf{Partie B.}
\begin{enumerate}
\setcounter{enumi}{4}
\item \textbf{(3 points)} Le système des deux condensateurs est isolé (débranché de la source), donc la charge totale est conservée : la charge se redistribue simplement entre les condensateurs jusqu'à ce qu'ils atteignent une tension commune \textbf{(1)}. Charge initiale $Q = CV = (10\ \mu\mathrm{F})(24\ \mathrm{V}) = 240\ \mu\mathrm{C}$. Tension commune finale :
\[ V_f = \frac{Q}{C + C} = \frac{240\ \mu\mathrm{C}}{20\ \mu\mathrm{F}} = 12\ \mathrm{V} \textbf{(2)}. \]
\item \textbf{(3 points)} Avant : $U_i = \tfrac{1}{2}(10\times10^{-6})(24)^2 = 2,88\times10^{-3}\ \mathrm{J} = 2,9\ \mathrm{mJ}$ \textbf{(1)}.\\
Après : $U_f = \tfrac{1}{2}(20\times10^{-6})(12)^2 = 1,44\times10^{-3}\ \mathrm{J} = 1,4\ \mathrm{mJ}$ \textbf{(1)}.\\
$U_f = \tfrac{1}{2}U_i$ : la moitié de l'énergie ($1,44\ \mathrm{mJ}$) a apparemment ``disparu'' \textbf{(1)}.
\item \textbf{(2 points)} L'énergie manquante est dissipée lors du passage du courant dans le circuit de connexion : (i) \textbf{effet Joule} dans la résistance des fils et des armatures \textbf{(1)} ; (ii) énergie rayonnée sous forme d'\textbf{ondes électromagnétiques} par l'étincelle/courant transitoire (et éventuellement son d'une étincelle) \textbf{(1)}. L'énergie est globalement conservée ; elle n'est plus stockée dans les condensateurs.\\
\emph{Attentes du correcteur : distinction nette entre conservation de la charge (toujours) et ``perte'' d'énergie (dissipation) ; utilisation correcte de $U \propto V^2$ pour la fraction $e^{-2}$ ; mécanismes physiques plausibles cités pour l'énergie manquante.}
\end{enumerate}
\end{corrige}

\newpage

\coursec{Summary sheet}

\begin{popfigure}
\begin{tikzpicture}[
  mindmap,
  concept color=popblue,
  level 1/.style={level distance=3.5cm, sibling angle=360/7},
  level 2/.style={level distance=2.5cm, sibling angle=360/7},
  every node/.style={font=\small, text=white, align=center},
  grow cyclic
]
\node[concept, minimum size=2.2cm] {Capacitors}
  child[concept color=popblue] { node[concept, minimum size=1.8cm, align=center]{Electric Potential\\ \& Work} }
  child[concept color=popgreen] { node[concept, minimum size=1.8cm] {Electric Dipole} }
  child[concept color=poporange] { node[concept, minimum size=1.8cm, align=center]{Capacitance\\ \& Permittivity} }
  child[concept color=poppurple] { node[concept, minimum size=1.8cm, align=center]{Series \& Parallel\\ Combinations} }
  child[concept color=poppink] { node[concept, minimum size=1.8cm] {Energy Stored} }
  child[concept color=popgold] { node[concept, minimum size=1.8cm, align=center]{Charging\\ \& Discharging} }
  child[concept color=popdark] { node[concept, minimum size=1.8cm, align=center]{Time Constant\\ $\tau = RC$} };
\end{tikzpicture}
\legende{Mind map of Chapter PHY05: Capacitors}
\end{popfigure}

\begin{bilan}

\paragraph{Key Definitions}
\begin{itemize}
  \item \cle{Electric potential} $V$ at a point: work done per unit positive charge to bring it from infinity to that point. $V = \frac{W}{q}$, unit: volt (V) = J/C.
  \item \cle{Potential difference} (p.d.) between two points: work done per unit charge moving between them. $V_{AB} = V_A - V_B = \frac{W_{AB}}{q}$.
  \item \cle{Electric dipole}: two equal and opposite charges $\pm q$ separated by distance $2a$. Dipole moment $\vec{p} = q(2\vec{a})$, direction from $-q$ to $+q$.
  \item \cle{Torque on a dipole} in uniform field $\vec{E}$: $\vec{\tau} = \vec{p} \times \vec{E}$, magnitude $\tau = pE\sin\theta$.
  \item \cle{Capacitance} $C$ of a conductor: ratio of charge $Q$ to its potential $V$. $C = \frac{Q}{V}$, unit: farad (F) = C/V.
  \item \cle{Parallel-plate capacitor}: $C = \frac{\varepsilon_0 \varepsilon_r A}{d}$ where $\varepsilon_0 = 8.85\times10^{-12}\ \mathrm{F/m}$, $\varepsilon_r$ = relative permittivity (dielectric constant), $A$ = plate area, $d$ = separation.
  \item \cle{Time constant} $\tau = RC$ for an RC circuit: time for charge/current/voltage to change by factor $1/e$ ($\approx 37\%$) during discharge, or to reach $\approx 63\%$ of final value during charging.
\end{itemize}

\paragraph{Laws, Formulas \& Theorems to Know}
\begin{itemize}
  \item Work done moving charge $q$ through p.d. $V$: $W = qV$.
  \item Energy stored in a capacitor: $U = \frac{1}{2}QV = \frac{1}{2}CV^2 = \frac{Q^2}{2C}$.
  \item Capacitors in \cle{parallel}: $C_{\text{eq}} = C_1 + C_2 + \cdots$; same $V$, charges add.
  \item Capacitors in \cle{series}: $\frac{1}{C_{\text{eq}}} = \frac{1}{C_1} + \frac{1}{C_2} + \cdots$; same $Q$, voltages add.
  \item Charging a capacitor through resistor $R$ from DC source $V_0$:
    \begin{align*}
      Q(t) &= C V_0 \left(1 - e^{-t/RC}\right), \\
      V_C(t) &= V_0 \left(1 - e^{-t/RC}\right), \\
      I(t) &= \frac{V_0}{R} e^{-t/RC}.
    \end{align*}
  \item Discharging a capacitor through resistor $R$:
    \begin{align*}
      Q(t) &= Q_0 e^{-t/RC}, \\
      V_C(t) &= V_0 e^{-t/RC}, \\
      I(t) &= -\frac{V_0}{R} e^{-t/RC} \quad (\text{magnitude } I_0 e^{-t/RC}).
    \end{align*}
  \item At $t=0$: capacitor acts as short circuit (uncharged) or open circuit (charged, no current). As $t\to\infty$: capacitor acts as open circuit (charged to $V_0$) or short circuit (discharged).
  \item Half-life for discharge: $t_{1/2} = \tau \ln 2 \approx 0.693\, \tau$.
\end{itemize}

\paragraph{Methods}
\begin{enumerate}
  \item \textbf{Finding equivalent capacitance}: Identify series/parallel groups, reduce step by step. Redraw circuit after each reduction.
  \item \textbf{Energy problems}: Use $U = \frac{1}{2}CV^2$ when $V$ known; $U = \frac{Q^2}{2C}$ when $Q$ known. For combinations, compute total energy from equivalent $C$ and total $V$ (or $Q$), or sum individual energies.
  \item \textbf{RC transient analysis}: Write Kirchhoff's loop equation $\to$ differential equation $\to$ solve with initial condition. Memorise the standard exponential forms above.
  \item \textbf{Graph interpretation}: Plot $\ln I$ vs $t$ (discharge) $\to$ straight line slope $=-1/RC$. Plot $\ln(V_0-V_C)$ vs $t$ (charging) $\to$ slope $=-1/RC$.
  \item \textbf{Experimental determination of $C$}: Use reed switch (charge/discharge at frequency $f$, measure average current $I = C V f$) or RC time constant method (measure $\tau$ with known $R$).
\end{enumerate}

\paragraph{Common Mistakes}
\begin{itemize}
  \item Confusing series/parallel formulas: remember \emph{capacitors add in parallel} (like resistors in series) and \emph{reciprocals add in series}.
  \item Forgetting that $Q$ is the same on all capacitors in series, but $V$ splits; in parallel $V$ is same, $Q$ splits.
  \item Using $U = QV$ instead of $\frac{1}{2}QV$ for capacitor energy.
  \item Mixing up charging/discharging equations: check $t=0$ and $t\to\infty$ limits.
  \item Ignoring the minus sign in discharge current (direction opposite to charging).
  \item Using $\tau = 1/RC$ instead of $\tau = RC$; dimensional check: $\Omega \cdot \mathrm{F} = \mathrm{s}$.
  \item Forgetting to convert $\mu\mathrm{F}$, $\mathrm{nF}$, $\mathrm{pF}$ to farads before calculation.
  \item In dipole torque, using $\cos\theta$ instead of $\sin\theta$; torque is maximum at $\theta=90^\circ$.
\end{itemize}

\paragraph{GCE Tips}
\begin{itemize}
  \item \textbf{Definitions} are frequently asked verbatim (e.g. capacitance, farad, relative permittivity, time constant). Learn them precisely.
  \item \textbf{Derivations}: Be able to derive $C = \varepsilon_0 A/d$ for parallel plates, energy $U = \frac{1}{2}CV^2$, and the charging/discharging equations from Kirchhoff's law.
  \item \textbf{Graph work}: Know how to linearise $I = I_0 e^{-t/RC}$ by plotting $\ln I$ vs $t$. The slope gives $-1/RC$. Label axes with units.
  \item \textbf{Numerical questions}: Always show the formula, substitution with units, and final answer with unit. Use $g=9.81$ or $10$ as instructed. Keep 3 significant figures unless told otherwise.
  \item \textbf{Practical}: Describe the reed switch method or the RC circuit with a stopwatch/voltmeter. Mention safety (discharge capacitor before handling).
  \item \textbf{Multiple choice}: Watch for distractors with wrong units (e.g. $\tau = R/C$), wrong exponential form, or confusion between $Q$, $V$, $I$ at $t=0$ and $t=\infty$.
  \item \textbf{Structured questions}: Often combine topics: e.g. a dipole in a uniform field created by a capacitor; or energy loss when connecting charged capacitors (charge redistributes, energy not conserved due to radiation/heat).
\end{itemize}

\end{bilan}

\findecours

\end{document}
